\documentclass[11pt]{article}
% preamble.tex -- shared by all section files of the short paper. Do not add macros in section files;
% ask the architect. Notation: short/OUTLINE.md, Section 4.
\usepackage[margin=1in]{geometry}
\usepackage[T1]{fontenc}
\usepackage{lmodern}
\usepackage{amsmath,amssymb,amsthm,mathtools}
\usepackage{booktabs,array}
\usepackage[hidelinks]{hyperref}

\numberwithin{equation}{section}
\renewcommand{\topfraction}{0.9}\renewcommand{\textfraction}{0.1}\renewcommand{\floatpagefraction}{0.8}

% ---------------------------------------------------------------- theorem environments (shared counter, per section)
\theoremstyle{plain}
\newtheorem{theorem}{Theorem}[section]
\newtheorem{proposition}[theorem]{Proposition}
\newtheorem{lemma}[theorem]{Lemma}
\newtheorem{corollary}[theorem]{Corollary}
\theoremstyle{definition}
\newtheorem{definition}[theorem]{Definition}
\theoremstyle{remark}
\newtheorem{remark}[theorem]{Remark}

% A named hypothesis/condition label: \namedlabel{hyp:end}{(S1)} makes \ref{hyp:end} print "(S1)".
\makeatletter
\newcommand{\namedlabel}[2]{\begingroup\def\@currentlabel{#2}\phantomsection\label{#1}\endgroup}
\makeatother
% Item tag for hypotheses: \hypitem{S1} gives "(S1)".
\newcommand{\hyptag}[1]{\textup{(#1)}}

% ---------------------------------------------------------------- general
\newcommand{\R}{\mathbb{R}}
\newcommand{\C}{\mathbb{C}}
\newcommand{\eps}{\varepsilon}
\newcommand{\epsS}{\varepsilon_*}           % 4e-11
\newcommand{\vz}{v_0}                        % 6e-5
\newcommand{\rhat}{\hat r}
\newcommand{\uhat}{\hat u}                   % frozen unit vector; leg direction \uhat_{j+1}
\newcommand{\nhat}{\hat n}                   % run-2 heading / touch-down tangent direction
\newcommand{\xz}{x^0}                        % curve at touch-down
\newcommand{\vsig}{\varsigma}                % T . e^perp
\newcommand{\Vu}{V_{\hat u}}                 % int |T - u|^2
\newcommand{\aD}{\alpha_{\mathrm D}}         % 1 + 1/K + 2/K^2
\newcommand{\toff}{\mathfrak t}              % time offset t - t_td
\DeclareMathOperator{\conv}{conv}
\DeclareMathOperator{\dist}{dist}

% ---------------------------------------------------------------- chase and end theorem (Sections 3-4)
\newcommand{\ttd}{t_{\mathrm{td}}}           % touch-down time
\newcommand{\Tstop}{T_{\mathrm{stop}}}
\newcommand{\Fam}{\mathcal F}                % family
\newcommand{\FamS}{\mathcal F_*}             % admissible family
\newcommand{\We}{W_{\mathrm e}}              % excuse length
\newcommand{\jW}{j_W}
\newcommand{\tW}{t_W}
\newcommand{\fF}{f_{\mathrm F}}
\newcommand{\dL}{d_{\mathrm L}}
\newcommand{\qc}{q_{\mathrm c}}
\newcommand{\chig}{\chi_{\mathrm g}}
\newcommand{\Keps}{K_\varepsilon}
\newcommand{\CTL}{C_{\mathrm{TL}}}
\newcommand{\CmuA}{C_{\mu\mathrm A}}
\newcommand{\ConeA}{C_{1\mathrm A}}
\newcommand{\camp}{c_{\mathrm{amp}}}
\newcommand{\PsiI}{\Psi_{\mathrm I}}
\newcommand{\TI}{T_{\mathrm I}}
\newcommand{\Ta}{T_\alpha}
\newcommand{\Tb}{T_\beta}
\newcommand{\DI}{D_{\mathrm I}}
\newcommand{\DIIs}{D_{\mathrm{IIs}}}
\newcommand{\DIIl}{D_{\mathrm{IIl}}}
\newcommand{\Xfar}{X_{\mathrm{far}}}
\newcommand{\LJ}{L_J}
\newcommand{\pA}{p_{\mathrm A}}              % stage-A tuple
\newcommand{\gD}{\gamma_{\mathrm D}}

% ---------------------------------------------------------------- heading, departure, assembly (Sections 5, 8)
\newcommand{\UA}{U_{\mathrm A}}
\newcommand{\UE}{U_{\mathrm E}}
\newcommand{\hH}{h_{\mathrm H}}
\newcommand{\Dbar}{\bar D}
\newcommand{\etaD}{\eta_{\mathrm D}}
\newcommand{\etaN}{\eta_{\mathrm N}}
\newcommand{\kD}{k_{\mathrm D}}
\newcommand{\Tsharp}{T_0^\sharp}

% ---------------------------------------------------------------- straight runs, design T, touch-down (Sections 6-7)
\newcommand{\sg}{\mathsf g}                  % sin(beta)/tan^v(beta/2)
\newcommand{\gmax}{\mathsf g_{\max}}
\newcommand{\gb}{\mathsf g_{\mathrm b}}
\newcommand{\sk}{\mathsf k}
\newcommand{\sa}{\mathsf a}
\newcommand{\sfc}{\mathsf c}
\newcommand{\sbb}{\mathsf b}
\newcommand{\sr}{\mathsf r}
\newcommand{\sA}{\mathsf A}
\newcommand{\TS}{\mathfrak S}
\newcommand{\ta}{t_{\mathrm a}}
\newcommand{\tc}{t_{\mathrm c}}
\newcommand{\as}{\alpha_{\mathrm s}}
\newcommand{\circab}[1]{|#1|_\circ}          % min(alpha, 2pi - alpha)
\newcommand{\xe}{\xi_{\mathrm e}}
\newcommand{\be}{b_{\mathrm e}}
\newcommand{\vbar}{\bar v}
\newcommand{\epsbar}{\bar\varepsilon}
\newcommand{\mub}{\bar\mu}
\newcommand{\muT}{\mu_{\mathrm T}}
\newcommand{\muTb}{\bar\mu_{\mathrm T}}
\newcommand{\vk}{\varkappa}
\newcommand{\kc}{\kappa_{\mathrm c}}
\newcommand{\czlo}{c_{\mathrm{zlo}}}
\newcommand{\czhi}{c_{\mathrm{zhi}}}
\newcommand{\Ups}{\Upsilon}
\newcommand{\Ghost}{\mathcal G}
\newcommand{\xend}{x_{\mathrm{end}}}


\title{A path fixed in advance that finds a slow fleeing point}
\author{}
\date{}

\begin{document}
\maketitle

\begin{abstract}
A point $S$ (the Snitch) starts at distance $d$ from the origin in an unknown direction $\alpha$ and always moves with speed $v$
directly away from a pursuer $H$ of speed at most $1$, which starts at the origin and must fix its path in advance; $S$ counts as
found once it is within distance $\eps$ of $H$. We prove that for $v\le6\cdot10^{-5}$ there is, for every $\eps\in(0,d)$, an
explicit path that finds the Snitch, whatever $\alpha$ is, by time $10.13\,d$; the time bound does not depend on $\eps$. By
compactness the same holds with $\eps=0$: a single path of speed at most $1$ catches the Snitch exactly, in every direction, by
time $10.13\,d$. The path runs once through the set of possible positions, turns so as to touch this set tangentially, and then
follows a greedy chase along a regularised (``virtual'') version of it. The proof is a bootstrap for the two-point excess (arc
minus chord) of the virtual front, uniform in $\eps$, combined with exact formulas for the front after the two straight runs.
All constants are explicit and are verified by a short script.
\end{abstract}

% sec1.tex -- writer W1.
\section{Introduction}\label{sec:intro}

\paragraph{The problem.} Identify $\R^2$ with $\C$. A \emph{pursuer path} is a map $H:[0,\infty)\to\R^2$ with $H(0)=0$ and
$|H'|\le1$ a.e., i.e.\ a $1$-Lipschitz path from the origin. A point $S$, the Snitch, starts at distance $d>0$ from the origin in
an unknown direction $\alpha\in[0,2\pi)$ and always moves with speed $v\in(0,1)$ directly away from the pursuer:
\begin{equation}\label{eq:snitch}
S_\alpha'(t)=v\,\frac{S_\alpha(t)-H(t)}{|S_\alpha(t)-H(t)|},\qquad S_\alpha(0)=d\,e^{i\alpha}.
\end{equation}
The pursuer knows $d$ and $v$ but not $\alpha$, and sees the Snitch only when it is within distance $\eps\in[0,d)$ (the
\emph{vision radius}); his path is therefore fixed in advance, i.e.\ it may not depend on $\alpha$. For a continuous $H$ the
right-hand side of \eqref{eq:snitch} is continuous in $t$ and locally Lipschitz in $S$ away from $H(t)$; we call $S_\alpha$, the
\emph{candidate} $\alpha$, the maximal solution, defined up to the first time at which it meets $H$ (since $|S_\alpha'|=v$ and
$\bigl||S_\alpha-H|'\bigr|\le1+v$, a maximal solution defined only on a bounded interval has a limit at its end, where it meets
$H$). Candidate $\alpha$ is \emph{caught} if $S_\alpha(t)=H(t)$ for some $t$ and \emph{seen} if $|S_\alpha(t)-H(t)|\le\eps$ for
some $t$ in its interval of existence or at its end. With $\eps>0$ the pursuer wins once the Snitch is seen; with $\eps=0$ he
must catch it. Since \eqref{eq:snitch} is invariant under $(S,H,t)\mapsto(\lambda S,\lambda H,\lambda t)$, one may take $d=1$.

The question is whether, for given $v$, a single path sees (or catches) \emph{every} direction, and in what time. A path that
runs straight to the circle of radius $d$ and once around it sees every direction when $(1+2\pi)vd\le\eps$: it passes $de^{i\alpha}$ before time $(1+2\pi)d$, when the Snitch is still within $(1+2\pi)vd$ of that point; the difficulty is to obtain a range of
$v$, and a time, that do not degenerate as $\eps\to0$.

\begin{theorem}\label{thm:main}
Let $\vz:=5.8\cdot10^{-3}$, $T_0:=13.1$ and $\epsS:=4\cdot10^{-11}$. For every $d>0$, every $\eps\in(0,d)$ and every
$v\in(0,\vz]$ the path $H=H_{v,\eps,d}$ --- design~T \textup{(Definition~\ref{def:T})} followed by the greedy chase
\textup{(Definition~\ref{def:chase})} on the working front of Theorem~\ref{thm:touchdown}, both for the vision radius
$\eps':=\min(\eps,\epsS d)$ and scaled to distance $d$ --- satisfies $H(0)=0$, $|H'|\le1$ a.e., depends only on $(v,\eps,d)$, and
for every $\alpha\in[0,2\pi)$ there is $t\le T_0d$ with $|S_\alpha(t)-H(t)|\le\eps$. Neither $\vz$ nor $T_0$ depends on $\eps$.
\end{theorem}

\begin{corollary}\label{cor:exact}
Let $d>0$ and $v\in(0,5.8\cdot10^{-3}]$. There is a path $H$ with $H(0)=0$, $|H'|\le1$ a.e., such that for every
$\alpha\in[0,2\pi)$ there is $t\le13.1\,d$ with $S_\alpha(t)=H(t)$.
\end{corollary}

Both are proved in Section~\ref{sec:proof}; the corollary follows from the theorem and the compactness Lemma~\ref{lem:compact},
because the time bound of Theorem~\ref{thm:main} is uniform in $\eps$.

\begin{remark}\label{rem:numerics}
The parameters of Appendix~\ref{app:const} were found by numerical optimisation; the script \texttt{verify\_constants.py}
checks every side condition in $50$-digit arithmetic; the near-pair constant $D'$ (Definition~\ref{def:Dprime}), a supremum
over a two-parameter family, is bounded there by an evaluation in outward-rounded arithmetic. At the chosen parameters several
side conditions are nearly tight at once (Table~\ref{tab:side}), so the method as it stands does not reach much beyond the
stated $\vz$. Simulations suggest that the same path works for much larger $v$.
\end{remark}

\paragraph{Outline of the proof.} We take $d=1$. Inside the vision disc the field $x/|x|$ of \eqref{eq:snitch} is singular, so
we replace it by a $C^1$ field $F$ that coincides with $x/|x|$ outside the disc of radius $\eps$ (Definition~\ref{def:field}).
Curves of \emph{virtual} candidates, moved by $vF(x-P)$, form $C^1$ fronts whose arclength never decreases and whose chords
never shrink (Lemmas~\ref{lem:flow}, \ref{lem:sep}), and a virtual candidate that comes within $\eps$ of the path is a true
candidate that has been seen (Lemma~\ref{lem:seen}). The pursuer can therefore follow such a front without knowing which of
its points carry the Snitch: the \emph{greedy chase} (Definition~\ref{def:chase}) walks along it in straight legs, each aimed
at the point an arclength $\eps$ further on, and every point passed is seen.

The chase finishes in a time independent of $\eps$ only if the front stays nearly straight on every scale between $\eps$ and a
fixed length $R$. This is measured by the two-point \emph{excess} $E=l-c$, arc length minus chord, of pairs of labels
(Definition~\ref{def:excess}); its evolution under the flow is governed by the exact identities of Lemma~\ref{lem:excess}, in
which the stretching $T\cdot AT$ of the front and the turning of its chord appear explicitly. The end theorem
(Theorem~\ref{thm:end}) is a bootstrap: if at the start the excess of a pair of label distance $\lambda$ is at most a small
multiple of $\lambda^{1+2\gamma}$, with $\gamma>\frac12$, and chords are comparable to $\lambda$, then the length bound
$l\le\kappa\lambda$ and the excess bound $N=(2E)^{1/2}\le M\lambda^{1/2+\gamma}$ persist along the whole chase, every
leg advances the label by at least $\eps/\kappa$, and the chase ends after time $\kappa U/(1-v)$, where $U$ is the length of
the front. Pairs far ahead of the pursuer are controlled by a Duhamel inequality (Lemma~\ref{lem:duhamel}), pairs near him by
integrating the excess identity over the legs, with the pursuer's position averaged along each leg (Section~\ref{sec:farnear}).
No constant depends on $\eps$.

The start data come from \emph{design~T} (Definition~\ref{def:T}). The pursuer first runs straight through the set of possible
positions, catching direction $0$ and seeing a thin sector around it; he then turns and runs along a straight line that
touches the remaining front tangentially, at the label $\alpha_0$, at the touch-down time $\ttd$. Along a straight run the
motion of the candidates has two exact integrals (Proposition~\ref{prop:runs}), and their derivatives in the direction label
transport the tangent and the curvature of the front exactly (Lemma~\ref{lem:transport}). This yields the front at touch-down
explicitly enough (Theorem~\ref{thm:touchdown}): it is a $C^1$ curve of bounded length that starts at the pursuer, is
flat at \emph{both} ends with tangent close to the run-2 heading, has excess $N\le c_N\lambda^{3/2-\muTb}$ on small scales and
chords bounded below on large ones (Theorem~\ref{thm:touchdown}, Lemma~\ref{lem:T2star}).

The front ends close to its start: its last point lies a distance at most of order $\eps^{1/2}$ behind the pursuer at $\ttd$. Pairs
with one point at the beginning of the front and the other near its end do not satisfy the start conditions. They are
\emph{excused}: the end theorem claims nothing about them until the chase has advanced by a fixed length $\We$. Meanwhile the
pursuer moves away from the end of the front with a controlled heading (Lemma~\ref{lem:heading}, obtained from a second
application of the end theorem to a short initial arc), and the departure theorem (Theorem~\ref{thm:departure}) shows that the
excused pairs are then far pairs with the start data the bootstrap needs. Assembling these pieces with explicit constants
(Section~\ref{sec:proof}, Appendix~\ref{app:const}) gives Theorem~\ref{thm:main}; letting $\eps\to0$ along a sequence of these
paths gives Corollary~\ref{cor:exact}.

\paragraph{Organisation.} Section~\ref{sec:prelim} collects the calculus of virtual fronts and of the excess. Section~\ref{sec:end}
defines the chase and states and organises the proof of the end theorem, which is completed in Section~\ref{sec:farnear}.
Section~\ref{sec:dep} proves the heading lemma and the departure theorem. Section~\ref{sec:designT} treats straight runs and
defines design~T, and Section~\ref{sec:touch} analyses the front at touch-down. Section~\ref{sec:proof} assembles the proof, and
Appendix~\ref{app:const} lists the constants and the side conditions they satisfy.

\paragraph{Notation.} Vectors in $\R^2=\C$: $a\cdot b:=\operatorname{Re}(\bar ab)$, $a\times b:=\operatorname{Im}(\bar ab)$,
$a^\perp:=ia$; thus $a^\perp\cdot b=a\times b$. Throughout Sections~\ref{sec:prelim}--\ref{sec:proof} (except in the last two proofs
of Section~\ref{sec:proof}) $d=1$, $0<v\le\vz$, $0<\eps\le\epsS$ (in Sections~\ref{sec:designT} and~\ref{sec:touch}:
$0<\eps\le10^{-3}$), and $\mu:=v/(1-v)$. Named constants are numbers, listed in Appendix~\ref{app:const}. The notation for fronts is
fixed in Section~\ref{sec:prelim}.

% sec2.tex -- writer W1.
\section{Virtual fronts and two-point calculus}\label{sec:prelim}

Throughout, $d=1$, $0<\eps<1$, and $P$ denotes a $1$-Lipschitz path with $P(0)=0$. External input: for an ordinary
differential equation $x'=G(t,x)$ with $G$ continuous and $D_xG$ continuous and bounded, solutions exist globally, are unique,
depend $C^1$ on the initial point, and $\partial x/\partial x_0$ solves the variational equation (see, e.g.,
\cite[Ch.~V]{Hartman}).

\subsection{The field and virtual flows}

\begin{definition}\label{def:field}
$f(\rho):=1$ for $\rho\ge\eps$ and $f(\rho):=(\rho/\eps)(2-\rho/\eps)$ for $0\le\rho\le\eps$; $F(x):=f(|x|)\,x/|x|$, $F(0):=0$;
$A:=DF$; $a^*(\rho):=1/\rho$ for $\rho\ge\eps$ and $a^*(\rho):=2/\eps$ for $\rho<\eps$.
\end{definition}

\begin{lemma}\label{lem:field}
\textup{(i)} $F\in C^1(\R^2)$, $|F|\le1$, and $F(x)=x/|x|$ for $|x|\ge\eps$.
\textup{(ii)} With $\rho=|x|$, $\rhat=x/\rho$, $\Pi=I-\rhat\rhat^{\mathsf T}$: $A=f'(\rho)\rhat\rhat^{\mathsf T}+(f(\rho)/\rho)\Pi$;
$A$ is symmetric, $0\le A\le a^*(\rho)I$, $A=\Pi/\rho$ for $\rho\ge\eps$, and $A(0)=(2/\eps)I$. Moreover
$f(\rho)/\rho-f'(\rho)=\rho/\max(\rho,\eps)^2$.
\textup{(iii)} For every unit vector $e$: $|e^\perp\cdot A(x)e|\le|x\times e|/\max(\rho,\eps)^2$.
\textup{(iv)} $(F(x)-F(y))\cdot(x-y)\ge0$ for all $x,y$; and $(x/|x|-y/|y|)\cdot(x-y)\ge0$ for $x,y\ne0$.
\end{lemma}

\begin{proof}
$F(x)=\sigma(|x|)x$ with $\sigma(\rho)=(2\eps-\rho)/\eps^2$ on $[0,\eps]$ and $\sigma=1/\rho$ on $[\eps,\infty)$. Both pieces and
their derivatives $-1/\eps^2$, $-1/\rho^2$ agree at $\eps$, so $\sigma\in C^1(0,\infty)$ and $F\in C^1(\R^2\setminus\{0\})$; near
$0$, $F(x)=(2\eps x-|x|x)/\eps^2$, and $x\mapsto|x|x$ is $C^1$ with derivative $|x|I+xx^{\mathsf T}/|x|\to0$. This gives (i)
($|F|=f\le1$) and $A(0)=(2/\eps)I$. For $x\ne0$, $DF=\sigma I+\rho\sigma'\rhat\rhat^{\mathsf T}=(f/\rho)\Pi+f'\rhat\rhat^{\mathsf T}$,
since $\sigma+\rho\sigma'=(\rho\sigma)'=f'$. For $\rho<\eps$: $f'=(2/\eps)(1-\rho/\eps)\in[0,2/\eps]$,
$f/\rho=(2-\rho/\eps)/\eps\in[1/\eps,2/\eps]$, $f/\rho-f'=\rho/\eps^2$; for $\rho\ge\eps$: $f'=0$, $f/\rho=1/\rho$. This is (ii).
(iii) By (ii), $A=(f/\rho)I-(f/\rho-f')\rhat\rhat^{\mathsf T}$, so
$e^\perp\cdot Ae=-\rho\max(\rho,\eps)^{-2}(e\cdot\rhat)(e^\perp\cdot\rhat)$, and $|e^\perp\cdot\rhat|=|e\times\rhat|=|x\times e|/\rho$;
at $x=0$ both sides vanish. (iv) $(F(x)-F(y))\cdot(x-y)=\int_0^1(x-y)\cdot A(y+s(x-y))(x-y)\,ds\ge0$ by (ii); and
$(x/|x|-y/|y|)\cdot(x-y)=|x|+|y|-(x\cdot y)(1/|x|+1/|y|)\ge0$ because $x\cdot y\le|x||y|$.
\end{proof}

\begin{definition}\label{def:flow}
Let $x^\circ:[u_-,u_+]\to\R^2$ be a $C^1$ curve with $\partial_ux^\circ\ne0$, and $t_\circ\ge0$. The \emph{virtual flow} of $P$
from $x^\circ$ at time $t_\circ$ is the solution of
\[
\partial_tx(u,t)=v\,F\bigl(x(u,t)-P(t)\bigr),\qquad x(u,t_\circ)=x^\circ(u),\qquad t\ge t_\circ .
\]
We write $x_u(t):=x(u,t)$, $r=r_u:=x-P$, $\rho:=|r|$, $\rhat:=r/\rho$, $w:=|\partial_ux|$, $T:=\partial_ux/w=e^{i\psi}$ (the unit
tangent, with angle $\psi$) and $\phi:=\arg r-\psi$, so that $\cos\phi=T\cdot\rhat$, $\sin\phi=T\times\rhat$.
\begin{enumerate}
\item[(a)] \emph{Virtual candidates}: $t_\circ=0$, $x^\circ(\theta)=e^{i\theta}$; we write $x_\theta(t)$ and $\rho_\theta$.
\item[(b)] \emph{Working front}: $t_\circ=\ttd$ (the touch-down time of Definition~\ref{def:T}), $x^\circ=\xz$ with
$|\partial_u\xz|\equiv1$ and labels $u\in[0,U]$; $u':=U-u$.
\item[(c)] A label $u$ of a working front is \emph{true with direction $\theta$} if the true candidate $S_\theta$ exists on
$[0,\ttd]$ and $\xz(u)=S_\theta(\ttd)$; otherwise it is a \emph{ghost label}.
\item[(d)] The \emph{true candidate} $S_\theta$ is the solution of $S'=v(S-P)/|S-P|$, $S(0)=e^{i\theta}$, up to the first time
it meets $P$ (that is, \eqref{eq:snitch} with $H=P$, $d=1$).
\end{enumerate}
\end{definition}

\begin{lemma}\label{lem:flow}
For a virtual flow:
\begin{enumerate}
\item[(a)] It exists and is unique on $[t_\circ,\infty)$, $|\partial_tx|\le v$, $x(\cdot,t)\in C^1$, $\partial_ux$ is continuous in
$(u,t)$, and $\partial_t\partial_ux=vA(r)\partial_ux$. The stretch $w$ is non-decreasing in $t$ (so $w\ge1$ when
$w(\cdot,t_\circ)\equiv1$), and lengths $l$ of label intervals are non-decreasing in $t$.
\item[(b)] $w_t=vw\,T\cdot AT$ and $T_t=v(AT-(T\cdot AT)T)$. Hence $\psi_t=-v\rho\max(\rho,\eps)^{-2}\sin\phi\cos\phi$ \textup(read
as $0$ if $\rho=0$\textup), so $|\psi_t|\le v\rho|\sin\phi|/\max(\rho,\eps)^2\le v/\max(\rho,\eps)$; and for $\rho\ge\eps$:
$w_t/w=v\sin^2\!\phi/\rho$, $\psi_t=-(v/\rho)\sin\phi\cos\phi$.
\item[(c)] For fixed labels, $l$, $c$, $E$ \textup(Definition~\ref{def:excess}\textup) are $C^1$ in $t$ \textup($c$ and $E$ on
time intervals where $c>0$\textup); their derivatives involve $P$ only through $P(t)$, so they do not jump when $P$ turns.
\end{enumerate}
\end{lemma}

\begin{proof}
(a) The field $(t,x)\mapsto vF(x-P(t))$ is continuous, and $C^1$ in $x$ with $\|D_x\|\le2v/\eps$ by Lemma~\ref{lem:field}(ii);
the ODE facts quoted above give everything except the monotonicity, and $|\partial_tx|=v|F|\le v$. By the variational equation
and $A\ge0$, $\partial_tw^2=2v\,\partial_ux\cdot A\,\partial_ux\ge0$; $l=\int w\,du$.
(b) $\partial_t(w^2)=2vw^2\,T\cdot AT$ gives $w_t$; then $T_t=\partial_t(\partial_ux/w)=vAT-(w_t/w)T$. Since $T_t\perp T$,
$\psi_t=T^\perp\cdot T_t=v\,T^\perp\cdot AT$. As $\rhat=e^{i\phi}T$, we have $T\cdot\rhat=\cos\phi$ and $T^\perp\cdot\rhat=\sin\phi$,
so by Lemma~\ref{lem:field}(ii), $T^\perp\cdot AT=-(f/\rho-f')(T^\perp\cdot\rhat)(T\cdot\rhat)=-\rho\max(\rho,\eps)^{-2}\sin\phi\cos\phi$
(and $T^\perp\cdot A(0)T=0$). For $\rho\ge\eps$, $A=\Pi/\rho$ and $T\cdot\Pi T=1-\cos^2\phi$.
(c) $|w_t|\le(2v/\eps)w$ and $w_t$ is continuous in $(u,t)$, so $l=\int_a^bw\,du$ is $C^1$; $x_a,x_b$ are $C^1$ in $t$ with
derivatives $vF(x-P(t))$.
\end{proof}

\begin{lemma}\label{lem:seen}
\textup{(a)} If $\rho_\theta(t)\le\eps$ for a virtual candidate, the true candidate $\theta$ is seen by time $t$.
\textup{(b)} If a label $u$ of a working front is true with direction $\theta$ and $\rho_u(t)\le\eps$ for some $t\ge\ttd$, then
$\theta$ is seen by time $t$.
\end{lemma}

\begin{proof}
(a) Since $\rho_\theta(0)=|e^{i\theta}-P(0)|=1>\eps$, $t':=\inf\{s\le t:\rho_\theta(s)\le\eps\}$ is positive and
$\rho_\theta\ge\eps$ on $[0,t']$, $\rho_\theta(t')=\eps$. There $F(x_\theta-P)=(x_\theta-P)/\rho_\theta$, so $x_\theta$ solves the
true equation; as that equation has a locally Lipschitz field away from $P$, $S_\theta=x_\theta$ on $[0,t']$, and
$|S_\theta(t')-P(t')|=\eps$. (b) If $S_\theta$ was within $\eps$ of $P$ at some time $\le\ttd$ there is nothing to prove.
Otherwise $\rho_u(\ttd)=|S_\theta(\ttd)-P(\ttd)|>\eps$, and the argument of (a), started at $\ttd$, applies.
\end{proof}

\begin{lemma}\label{lem:sep}
Let $X_a,X_b$ solve $X'=vV(X-P)$ on $[t_1,t_2]$, each with its own $V\in\{F,\ x\mapsto x/|x|\}$ \textup(for $x/|x|$: a true
candidate not meeting $P$ on $[t_1,t_2]$\textup).
\textup{(a)} If both use $F$, $|X_b-X_a|$ is non-decreasing.
\textup{(b)} In general, if $|X_b(t_1)-X_a(t_1)|\ge2\eps$, then $|X_b(t)-X_a(t)|\ge|X_b(t_1)-X_a(t_1)|$ on $[t_1,t_2]$.
\textup{(c)} In particular, chords $c(a,b;t)$ of a virtual flow are non-decreasing in $t$.
\end{lemma}

\begin{proof}
Put $D:=X_b-X_a$ and $p_\bullet:=X_\bullet-P$, so $D=p_b-p_a$ and $(|D|^2)'=2v\bigl(V_b(p_b)-V_a(p_a)\bigr)\cdot(p_b-p_a)$. If
$V_a=V_b$ this is $\ge0$ by Lemma~\ref{lem:field}(iv); this proves (a), and (c) follows. If $V_a\ne V_b$ and $|D|\ge2\eps$, then
$|p_a|\ge\eps$ or $|p_b|\ge\eps$; at that point both fields agree (Lemma~\ref{lem:field}(i)), so $(|D|^2)'$ is the corresponding
expression for a single monotone field, and $(|D|^2)'\ge0$. Thus in all cases $(|D|^2)'\ge0$ whenever $|D|\ge2\eps$.

Let $c_0:=|D(t_1)|$. If $c_0>2\eps$ and $|D|<c_0$ somewhere, let $t^*$ be the infimum of such times. Then $|D|\ge c_0$ on
$[t_1,t^*]$ and $|D|>2\eps$ on $[t_1,t^*+\delta)$ for some $\delta>0$, where $|D|^2$ is non-decreasing; so $|D|\ge c_0$ on
$[t_1,t^*+\delta)$, contradicting the choice of $t^*$. If $c_0=2\eps$, let $X_b^\eta$ solve the equation of $X_b$ with
$X_b^\eta(t_1)=X_b(t_1)+\eta D(t_1)/c_0$, $\eta>0$. Since $X_b$ stays at positive distance from $P$ when $V_b=x/|x|$, continuous
dependence on initial data gives, for small $\eta$, that $X_b^\eta$ exists on $[t_1,t_2]$ and $X_b^\eta\to X_b$ uniformly as
$\eta\to0$; the case just proved gives $|X_b^\eta-X_a|\ge c_0+\eta$, and we let $\eta\to0$.
\end{proof}

\subsection{Two-point quantities}

\begin{definition}\label{def:excess}
For a regular $C^1$ curve $u\mapsto x_u$ (such as a virtual flow at a fixed time), with arclength $ds=w\,du$ and unit tangent $T$,
and labels $a<b$: $l(a,b):=\int_a^bds$, $c(a,b):=|x_b-x_a|$, $E:=l-c\ge0$, $N:=(2E)^{1/2}$, and, when $c>0$,
$e=e_{ab}:=(x_b-x_a)/c$. For a virtual flow the time is written as a subscript or third argument: $l_t(a,b)$, $c(a,b;t)$. We put
$\vsig:=T\cdot e^\perp$ and $q:=e^\perp\cdot A(r)e$ \textup(functions of $s\in[a,b]$\textup). For a fixed unit vector $\uhat$ and a
label interval $J$: $\Vu(J):=\int_J|T-\uhat|^2ds$.
\end{definition}

\begin{lemma}\label{lem:static}
Let $u\mapsto x_u$ be a regular $C^1$ curve and $J=(a,b)$ a label interval.
\begin{enumerate}
\item[(a)] $E(a',b')\le E(a,b)$ for $a\le a'<b'\le b$.
\item[(b)] $\Vu(J)=2[l(J)-(x_b-x_a)\cdot\uhat]=N(J)^2+c(J)|e_J-\uhat|^2$. Hence $N^2\le\Vu$, $|e_J-\uhat|^2\le\Vu/c$,
$\int_J|T-e_J|^2ds=N(J)^2$, $\Vu$ is additive over adjacent intervals, and
$|\uhat\times(x_b-x_a)|\le\int_J|T-\uhat|\,ds\le(l\,\Vu)^{1/2}$.
\item[(c)] If the tangent angle is $K$-Lipschitz in arclength on $J$ and $Kl\le\pi$, then $N\le(K/\sqrt{12})\,l^{3/2}$.
\item[(d)] \textup{(sub-chord tilt)} $|e_{a'b'}-e_{ab}|\le[l(a',b')^{1/2}N(a,b)+E(a',b')]/c(a',b')$ for $a\le a'<b'\le b$.
\end{enumerate}
\end{lemma}

\begin{proof}
(a) $l(a,b)=l(a,a')+l(a',b')+l(b',b)$ and $c(a,b)\le c(a,a')+c(a',b')+c(b',b)$; subtract and use $E\ge0$ on the outer pieces.
(b) $|T-\uhat|^2=2-2T\cdot\uhat$ and $\int_JT\,ds=x_b-x_a=c\,e_J$ give the first equality; then
$2(l-c\,e_J\cdot\uhat)=2(l-c)+c|e_J-\uhat|^2$. With $\uhat=e_J$ this is $\int_J|T-e_J|^2=N^2$. Finally
$\uhat\times(x_b-x_a)=\int_J\uhat\times(T-\uhat)\,ds$, and Cauchy--Schwarz.
(c) We may take $K>0$. Let $\psi$ be the tangent angle as a function of arclength $s\in[0,l]$, $\psi_{\mathrm m}:=\psi(l/2)$. Then
$|\psi-\psi_{\mathrm m}|\le K|s-l/2|\le\pi/2$, so
$c\ge(x_b-x_a)\cdot e^{i\psi_{\mathrm m}}=\int_0^l\cos(\psi-\psi_{\mathrm m})\,ds\ge\int_0^l\cos(K|s-l/2|)\,ds=(2/K)\sin(Kl/2)\ge l-K^2l^3/24$;
thus $N^2=2E\le K^2l^3/12$.
(d) $c(a',b')e_{a'b'}-l(a',b')e_{ab}=\int_{a'}^{b'}(T-e_{ab})\,ds$ has norm $\le l(a',b')^{1/2}N(a,b)$ by Cauchy--Schwarz and (b),
and equals $c(a',b')(e_{a'b'}-e_{ab})-E(a',b')e_{ab}$.
\end{proof}

\begin{lemma}\label{lem:excess}
For a virtual flow, fixed labels $a<b$ and any time $t$ with $c(a,b;t)>0$ \textup(primes are $t$-derivatives, integrals are over
$s\in[a,b]$, $A=A(r(s))$\textup):
\begin{enumerate}
\item[(i)] $l'=v\int T\cdot AT\,ds$ and $c'=v\int e\cdot AT\,ds$.
\item[(ii)] $E'=v\int(T-e)\cdot A(T-e)\,ds-\frac v2\int(e\cdot Ae)|T-e|^2ds+v\int\vsig q\,ds$; moreover $\int|T-e|^2ds=N^2$,
$\int\vsig\,ds=0$ and $|\vsig|\le|T-e|$.
\item[(iii)] For every constant $\bar q$: $E'\le v\int a^*(\rho)|T-e|^2ds+v\int\vsig(q-\bar q)\,ds$.
\item[(iv)] $E'\le vN\|AT\|_{L^2(ds)}$.
\item[(v)] \textup{(far inequality)} If $\rho_m:=\dist(P,\conv\{x_s:a\le s\le b\})\ge\eps$ and $p_a:=r_a\times e$, then
\[
E'\le\frac v{\rho_m}\Bigl(1+\frac l{\rho_m}+\frac{2l^2}{\rho_m^2}\Bigr)N^2+\frac{2vl^{3/2}|p_a|N}{\rho_m^3},\qquad
l'\le\frac{2v}{\rho_m}\Bigl[\frac{l(|p_a|+l^{1/2}N)^2}{\rho_m^2}+N^2\Bigr].
\]
\end{enumerate}
\end{lemma}

\begin{proof}
(i) $l'=\int w_t\,du$ (Lemma~\ref{lem:flow}(c)) and $w_t=vw\,T\cdot AT$. Next $c'=e\cdot(x_b'-x_a')=v\,e\cdot(F(r_b)-F(r_a))$ and
$F(r_b)-F(r_a)=\int_a^bA(r)\partial_ux\,du=\int AT\,ds$.
(ii) By (i), $E'=v\int(T-e)\cdot AT$, and $(T-e)\cdot AT=(T-e)\cdot A(T-e)+(T-e)\cdot Ae$. Decomposing $T-e$ in the basis
$e,e^\perp$ and using the symmetry of $A$: $(T-e)\cdot Ae=[(T-e)\cdot e](e\cdot Ae)+[(T-e)\cdot e^\perp]\,q$, where
$(T-e)\cdot e=-|T-e|^2/2$ and $(T-e)\cdot e^\perp=\vsig$. Further, $\int|T-e|^2=N^2$ by Lemma~\ref{lem:static}(b),
$\int\vsig=(x_b-x_a)\cdot e^\perp=0$, and $|\vsig|=|(T-e)\cdot e^\perp|\le|T-e|$.
(iii) In (ii), $(T-e)\cdot A(T-e)\le a^*|T-e|^2$ and $e\cdot Ae\ge0$ (Lemma~\ref{lem:field}(ii)), and $\int\vsig q=\int\vsig(q-\bar q)$.
(iv) Cauchy--Schwarz in $E'=v\int(T-e)\cdot AT$, with $\int|T-e|^2=N^2$.

(v) Let $\mathcal K:=\conv\{x_s:a\le s\le b\}-P$; on $\mathcal K$, $|y|\ge\rho_m\ge\eps$, so $A(y)=\Pi(y)/|y|$ and $q(s)=Q(r(s))$
with $Q(y):=-(y\cdot e)(y\cdot e^\perp)/|y|^3$. In the coordinates $y_1:=y\cdot e$, $y_2:=y\cdot e^\perp$,
$\partial_{y_2}Q=-y_1(y_1^2-2y_2^2)/|y|^5$ and $\partial_{y_1}Q=-y_2(y_2^2-2y_1^2)/|y|^5$; hence $|\partial_{y_2}Q|\le|y|^{-2}$
(as $\max_{|z|\le1}|z(3z^2-2)|=1$) and $|\partial_{y_1}Q|\le2|y_2|/|y|^3$. For $s\in[a,b]$ the segment
$y(\sigma)=r_a+\sigma(x_s-x_a)$, $\sigma\in[0,1]$, lies in $\mathcal K$; along it $dy_2=\delta(s)\,d\sigma$,
$dy_1=((x_s-x_a)\cdot e)\,d\sigma$,
where $\delta(s):=(x_s-x_a)\cdot e^\perp=\int_a^s\vsig$ satisfies $|\delta(s)|\le\int|T-e|\le l^{1/2}N$; moreover
$|(x_s-x_a)\cdot e|\le l$ and $|y_2|\le|p_a|+|\delta(s)|$ (note $|r_a\cdot e^\perp|=|r_a\times e|=|p_a|$). Therefore
\[
|q(s)-q(a)|\le\frac{l^{1/2}N}{\rho_m^2}+\frac{2l(|p_a|+l^{1/2}N)}{\rho_m^3}.
\]
Apply (iii) with $\bar q=q(a)$, $a^*(\rho)\le1/\rho_m$ on the arc, $\int|T-e|^2=N^2$ and $\int|\vsig|\le l^{1/2}N$:
$E'\le(v/\rho_m)N^2+vl^{1/2}N\sup_s|q(s)-q(a)|$, which is the first claim. For the length, $T\cdot AT=\sin^2\!\phi/\rho$ and
$\rho\sin\phi=T\times r=e\times r+(T-e)\times r$, where $|e\times r(s)|=|r(s)\cdot e^\perp|=|\delta(s)-p_a|\le|p_a|+l^{1/2}N$. Hence
$\sin^2\!\phi/\rho\le2(|p_a|+l^{1/2}N)^2/\rho^3+2|T-e|^2/\rho$; integrate with $\rho\ge\rho_m$ and use (i).
\end{proof}

\begin{lemma}\label{lem:duhamel}
Let $a<b$ be labels of a virtual flow with $c(a,b;t_1)>0$, and suppose $\rho_m(t)\ge\eps$ on $[t_1,t_2]$. With
$\alpha:=(v/\rho_m)(1+l/\rho_m+2l^2/\rho_m^2)$ and $\beta:=2vl^{3/2}|p_a|/\rho_m^3$ \textup(functions of $t$\textup),
\[
N(t_2)\le e^{\int_{t_1}^{t_2}\alpha}N(t_1)+\int_{t_1}^{t_2}e^{\int_t^{t_2}\alpha}\beta(t)\,dt .
\]
If moreover $\rho_m\ge Kl$ on $[t_1,t_2]$, then $\alpha\le\aD(K)\,v/\rho_m$, where $\aD(K):=1+1/K+2/K^2$.
\end{lemma}

\begin{proof}
By Lemma~\ref{lem:sep}(c), $c>0$ on $[t_1,t_2]$; $\alpha,\beta$ are continuous. For $\eta>0$ let $N_\eta:=(N^2+\eta^2)^{1/2}$, a
$C^1$ function (Lemma~\ref{lem:flow}(c)). By Lemma~\ref{lem:excess}(v), $N_\eta N_\eta'=E'\le\alpha N^2+\beta N\le\alpha
N_\eta^2+\beta N_\eta$, so $N_\eta'\le\alpha N_\eta+\beta$, i.e.\ $(N_\eta e^{-\int_{t_1}^t\alpha})'\le\beta e^{-\int_{t_1}^t\alpha}$.
Integrate and let $\eta\to0$. The last claim is $l/\rho_m\le1/K$.
\end{proof}

\subsection{Vanishing vision radius}

\begin{lemma}\label{lem:compact}
Let $d>0$, let $D_1\subset D_2\subset\cdots$ be finite sets of directions whose union is dense in $[0,2\pi)$, and $\eps_N\to0$.
Suppose that for each $N$ a path $H_N$ with $H_N(0)=0$, $|H_N'|\le1$ on $[0,\bar T]$ comes within $\eps_N$ of every $S_\alpha$,
$\alpha\in D_N$, by time $\bar T$ \textup($\bar T$ independent of $N$; $S_\alpha$ is the candidate of \eqref{eq:snitch} with $H=H_N$\textup).
Then there is a path $H$ with $H(0)=0$, $|H'|\le1$ on $[0,\bar T]$, which catches every direction by time $\bar T$.
\end{lemma}

\begin{proof}
By Arzel\`a--Ascoli a subsequence, still written $H_N$, converges uniformly on $[0,\bar T]$ to a $1$-Lipschitz path $H$, and
$H(0)=0$. For a $1$-Lipschitz path $G$ write $\rho^G_\alpha:=|S_\alpha-G|$ for the candidate $\alpha$ under $G$; it is defined until
the catch (Section~\ref{sec:intro}) and $|(\rho^G_\alpha)'|\le1+v$. We use $|y/|y|-z/|z||\le2|y-z|/\max(|y|,|z|)$, which follows
from $|y/|y|-z/|y||+|z/|y|-z/|z||\le2|y-z|/|y|$ and symmetry.

\emph{Claim.} Let $G,G_k$ be $1$-Lipschitz paths on $[0,\bar T]$ with $\|G_k-G\|_\infty\to0$, and let $\rho^G_\alpha\ge c>0$ on
$[0,\tau]$, $\tau\le \bar T$. Put $\tau':=\min(\tau+c/(2(1+v)),\bar T)$. Then there are a neighbourhood $\mathcal U$ of $\alpha$ and $k_0$ such
that for $\alpha'\in\mathcal U$ and $k\ge k_0$ the candidate $\alpha'$ under $G_k$ is defined and has $\rho^{G_k}_{\alpha'}\ge c/4$ on
$[0,\tau']$. Indeed, as $|(\rho^G_\alpha)'|\le1+v$, the candidate $\alpha$ under $G$ is defined with $\rho^G_\alpha\ge c/2$ on
$[0,\tau']$. Write $S:=S_\alpha$ under $G$ and $\tilde S:=S_{\alpha'}$ under $G_k$.
While $|S-G|,|\tilde S-G_k|\ge c/4$, the inequality above gives
$|\tilde S-S|'\le(8v/c)(|\tilde S-S|+\|G_k-G\|_\infty)$, hence by Gr\"onwall
$|\tilde S-S|\le b_k:=\bigl(d|e^{i\alpha'}-e^{i\alpha}|+(8v/c)\bar T\|G_k-G\|_\infty\bigr)e^{8v\bar T/c}$. Choose $\mathcal U$, $k_0$ with
$b_k+\|G_k-G\|_\infty<c/4$. Let $I$ be the set of $s\in[0,\tau']$ such that $\tilde S$ is defined on $[0,s]$ with $|\tilde S-G_k|\ge c/4$
there. Then $0\in I$ ($|\tilde S(0)-G_k(0)|=d\ge c$), $I$ is closed, and for $s\in I$ the bound gives
$|\tilde S-G_k|\ge|S-G|-b_k-\|G_k-G\|_\infty>c/4$ on $[0,s]$, so $I$ contains a neighbourhood of $s$ in $[0,\tau']$. Hence
$I=[0,\tau']$.

Fix $\alpha\in D_M$. For $N\ge M$, $\alpha\in D_N$; let $t_N\le \bar T$ be the first time with $\rho^{H_N}_\alpha(t_N)\le\eps_N$ (it
exists, as a catch time is such a time), and pass to a subsequence with $t_N\to t^*\le \bar T$. If $H$ did not catch $\alpha$ by time
$t^*$, then $\rho^H_\alpha$ would be defined and positive on $[0,t^*]$, hence $\ge c>0$ there, and the Claim ($G=H$, $G_k=H_N$,
$\alpha'=\alpha$, $\tau=t^*$) would give $\rho^{H_N}_\alpha\ge c/4$ on $[0,\min(t^*+c/(2(1+v)),\bar T)]\ni t_N$ for large $N$,
contradicting $\eps_N\to0$. So $H$ catches every direction of $\bigcup_ND_N$ by time $\bar T$. Finally, the Claim with $G_k=G=H$ and
$\tau=\bar T$ shows that the set of directions not caught by $H$ on $[0,\bar T]$ is open. It is disjoint from the dense set
$\bigcup_ND_N$, hence empty.
\end{proof}

% sec3.tex -- writer W2. Section 3: the greedy chase and the end theorem.
\section{The greedy chase and the end theorem}\label{sec:end}

Throughout Sections~\ref{sec:end} and~\ref{sec:farnear}, $0<v\le\vz$, $\mu=v/(1-v)$, and $x$ is a working front
(Definition~\ref{def:flow}(b)) on $[0,U]$ from $\xz$ at $\ttd$, with $|\partial_u\xz|\equiv1$, so that a label is its arclength
coordinate at $\ttd$. For labels $a<b$ we write $\lambda:=b-a$ (the \emph{label length}) and call $(a,b)$ a \emph{pair}; for a pair
$J$ we write $\lambda_J,l(J),c(J),e_J,N(J)$.

\subsection{The chase}

\begin{definition}[Greedy chase]\label{def:chase}
Let $P$ be a $1$-Lipschitz path on $[0,\ttd]$ with $\xz(0)=P(\ttd)$, and let $b_{\mathrm c}\in(0,U]$ (the \emph{cut}). Put
$t_0:=\ttd$, $a_0:=0$. Given a leg start $t_j$ with $P(t_j)=x_{a_j}(t_j)$, let
\[
a_{j+1}:=\max\{u\in[a_j,b_{\mathrm c}]:\ l_{t_j}(a_j,u)\le\eps\}.
\]
If $a_{j+1}=b_{\mathrm c}$, the chase \emph{stops at} $t_j=:\Tstop$: leg $j+1$ is \emph{final}, it is not walked, and $P$ is
constant after $\Tstop$. Otherwise $P$ runs at unit speed along the ray from $P(t_j)$ through $x_{a_{j+1}}(t_j)$ until it meets
$x_{a_{j+1}}$, at time $t_{j+1}$; this requires $c(a_j,a_{j+1};t_j)>0$ (the \emph{proviso}). Leg $j+1$ is the interval
$[t_j,t_{j+1}]$, and $\uhat_{j+1}$ is its direction. The working front is computed with the path so extended.
\end{definition}

For $t\in(t_j,t_{j+1}]$, and for $t=t_0$ with $j=0$, we write $j(t):=j$; $a_{j(t)}$ is the \emph{base} at time $t$. A leg start
$t_j$, $j\ge1$, belongs to the legs $j$ and $j+1$: statements ``in leg $j+1$'' may be applied at $t_j$ with base $a_j$, and
statements ``in leg $j$'' at $t_j$ with base $a_{j-1}$; both are valid since $P(t_j)=x_{a_j}(t_j)$. Likewise, for a tail pair the
time $\tW$ is treated as a point of leg $\jW+1$, with base $a_{\jW}$ (Lemma~\ref{lem:local} and Section~\ref{sec:farnear}). Since $x(u,t)$ depends only on $\xz(u)$ and on $P|_{[\ttd,t]}$, the chase cut at
$b_{\mathrm c}$ depends only on $\xz|_{[0,b_{\mathrm c}]}$ and on the path up to the current time.

\begin{lemma}[The chase]\label{lem:chase}
Assume the proviso at every non-final leg start reached.
\begin{enumerate}
\item[(0)] Either the chase stops or $t_j\to\infty$; $P$ is $1$-Lipschitz.
\item[(1)] During leg $j+1$: $x_{a_{j+1}}=P+\rho\uhat_{j+1}$ with $\rho'=-1+vf(\rho)\le-(1-v)$, and $x_{a_j}=P-\sigma\uhat_{j+1}$ with
$\sigma\ge0$. So $x_{a_j},P,x_{a_{j+1}}$ are collinear in this order, and the leg lasts at most
$c(a_j,a_{j+1};t_j)/(1-v)\le\eps/(1-v)$.
\item[(2)] Every true label $u\in[a_j,a_{j+1}]$ has its direction seen by time $t_j$; this includes the final leg.
\item[(3)] During leg $j+1$: $l_t(a_j,a_{j+1})\le e^{2\mu}\eps$, $|P-x_{a_j}|\le e^{2\mu}\eps$, and $l_t(a,b)\le e^{2\mu}l_{t_j}(a,b)$
for all $a<b$.
\item[(4)] If a path $\tilde P$ equals $P$ on $[0,t_j]$, the working fronts of $P$ and $\tilde P$ agree on $[\ttd,t_j]$.
\end{enumerate}
\end{lemma}

\begin{proof}
By Lemma~\ref{lem:flow}(b) and Lemma~\ref{lem:field}(ii), $w_t=vw\,T\cdot AT\le(2v/\eps)w$ and, by Lemma~\ref{lem:excess}(i),
$l'\le(2v/\eps)l$.
(0) On $[\ttd,T]$, $w\le e^{2v(T-\ttd)/\eps}$. A non-final leg $j+1$ has $l_{t_j}(a_j,a_{j+1})=\eps$ (maximality and continuity of $l$
in the label), so if $t_j\le T$ it advances the base by at least $\eps e^{-2v(T-\ttd)/\eps}$; hence finitely many legs start
before $T$. $P$ is piecewise linear with speed $\le1$.
(1) Let $\rho$ solve $\rho'=-1+vf(\rho)$, $\rho(t_j)=c(a_j,a_{j+1};t_j)>0$. As $F(\rho\uhat)=f(\rho)\uhat$, $y:=P+\rho\uhat_{j+1}$ solves
$y'=vF(y-P)$ with $y(t_j)=x_{a_{j+1}}(t_j)$, so $y=x_{a_{j+1}}$ by uniqueness while $\rho>0$; $\rho'\le-(1-v)$ as $f\le1$. Likewise
$P-\sigma\uhat_{j+1}$ with $\sigma'=1+vf(\sigma)$, $\sigma(t_j)=0$ solves the equation of $x_{a_j}$ (note $F(-\sigma\uhat)=-f(\sigma)\uhat$),
and $\sigma\ge0$. Finally $c\le l_{t_j}(a_j,a_{j+1})\le\eps$.
(2) $\rho_u(t_j)=c(a_j,u;t_j)\le l_{t_j}(a_j,u)\le\eps$; apply Lemma~\ref{lem:seen}(b).
(3) The leg lasts at most $\eps/(1-v)$, so lengths grow at most by the factor $e^{(2v/\eps)\eps/(1-v)}=e^{2\mu}$; and
$|P-x_{a_j}|=\sigma\le c(a_j,a_{j+1};t)\le l_t(a_j,a_{j+1})$ by (1).
(4) Uniqueness for the equation of Definition~\ref{def:flow}.
\end{proof}

\subsection{Parameters and the end theorem}

\begin{definition}[Parameters]\label{def:param}
A \emph{parameter tuple} is
\[
p=(\gamma,R,m,s_0,\fF,\kappa,c_A,K_R,K_f,\Keps,\eta_1,\qc,\nu,\dL,\delta_1,A_0).
\]
Its \emph{derived constants} are
\begingroup\allowdisplaybreaks
\begin{align*}
&M:=mR^{-\gamma},\quad \chig:=\kappa\dL,\quad c_6:=\frac{1+\delta_1}{c_A},\quad \sigma_q:=\frac1{1-\qc^{-\gamma}},\quad
r_*:=\max\Bigl(\qc,\frac{2\qc}{\qc-1}\Bigr),\\
&\mathsf s(r):=\frac{\kappa^{1/2}}{c_A}\cdot\begin{cases}r/2,&1\le r\le2,\\(r-1)^{1/2},&r\ge2,\end{cases}\qquad
\omega_r:=\Bigl(1-\frac{m^2r}{4c_A}\Bigr)^{-1/2},\\
&\CT:=\omega_{r_*}\max\Bigl(\sigma_q\,\mathsf s(\qc),\,\mathsf s\Bigl(\frac{2\qc}{\qc-1}\Bigr)\Bigr),\quad \Cmu:=1+\kappa\CT^2,\\
&c_\rho^*:=c_A-\frac{2\delta_1}{K_f},\quad c_\eta:=c_A-\frac{\delta_1}{\Keps\eta_1},\quad c_R^*:=c_AK_R-\frac{\delta_1}{\Keps},\\
&K_1:=\frac{c_\rho^*K_f}{\kappa},\quad K_2:=\frac{c_R^*}{\kappa},\quad \camp:=\frac{\aD(K_1)\chig}{c_\rho^*},\\
&B_{2U}':=\frac{2\kappa^{3/2}(1+1/K_2)}{\pi},\quad
B_{2\eta}^*:=\frac{2\kappa^{3/2}(1+1/K_1)}{\pi c_\eta^2}\Bigl(\frac{\chig}{\eta_1}+\frac{\dL}{\Keps\eta_1^2}\Bigr),\\
&C_p':=\kappa^{1/2}\Bigl(1+\frac1{K_f}\Bigr)^{1/2+\gamma}K_f^{-1/2}+\kappa\CT\Bigl(1+\frac1{K_f}\Bigr)^\gamma,\quad
C_p^V:=(\kappa\Cmu)^{1/2}\Bigl(1+\frac1{K_f}\Bigr)^{1/2+\gamma},\\
&C_p^*:=\min(C_p',C_p^V),\quad C_{\mathrm{ar}}:=\frac{\kappa^{3/2}(1+1/K_f)^{3/2+\gamma}}{c_A},\quad
\hat C(x):=\min(C_p^*,C_{\mathrm{ar}}x^{1/2}),\\
&B_3':=\frac{2\kappa^{3/2}K_f^{-\nu}}{\pi(c_\rho^*)^3}\Bigl[2\dL K_f^{\nu+\gamma-2}\hat C(K_f)+\chig\int_{K_f}^\infty x^{\nu+\gamma-2}\hat C(x)\,dx\Bigr],\\
&\PsiI:=2+\log K_f,\quad\TI:=\dL(\kappa c_6+2)+c_6\chig\PsiI,\\
&B_5^*:=\frac1{c_\rho^*}\Bigl[\frac{C_p^*+\kappa^{1/2}K_f^{-1-\gamma}/\pi}{c_\rho^*}+K_f^{-\gamma}\Bigr]^2,\quad
I_5:=\frac{\chig(\eta_1+1/\Keps)^{2\gamma}}{2\gamma},
\end{align*}
where $\aD$ is the function of Lemma~\ref{lem:duhamel}; and $\bar\theta:=\fF^2+2\vz D'$, with the near-pair constant $D'$ of
Definition~\ref{def:Dprime} (it depends only on the tuple). With $x_*:=(C_p^*/C_{\mathrm{ar}})^2$ and $\hat p:=\nu+\gamma-2\in(-\frac32,-1)$,
the integral in $B_3'$ is $C_{\mathrm{ar}}(x_*^{\hat p+3/2}-K_f^{\hat p+3/2})/(\hat p+\frac32)+C_p^*x_*^{\hat p+1}/(-\hat p-1)$ if $x_*>K_f$
(then $\hat C(K_f)=C_{\mathrm{ar}}K_f^{1/2}$), and $C_p^*K_f^{\hat p+1}/(-\hat p-1)$ otherwise (then $\hat C(K_f)=C_p^*$). For a length $U_0>0$ and a
non-decreasing function $\Cfl:[0,\infty)\to[0,\infty)$ put $\rho_{\Cfl}(y):=\max(c_AK_RR,\Cfl(y))-\delta_1R/\Keps$ and, for $k=1,2$,
\begin{align*}
&J_k^{\Cfl}(U_0):=\frac{2\dL R}{\Keps(c_R^*R)^k}+\chig\Bigl[\int_{K_RR}^{\max(K_RR,U_0/2)}\rho_{\Cfl}^{-k}+\int_0^{U_0/2}\rho_{\Cfl}^{-k}\Bigr],\\
&J_k(U_0):=\min\Bigl\{\frac{\chig U_0}{(c_R^*R)^k},\,J_k^{\Cfl}(U_0)\Bigr\},\\
&\Xfar(U_0):=\vz\Bigl[J_1(U_0)+\frac1{c_\eta}\Bigl(\frac{\dL}{\Keps\eta_1}+\chig\log\frac{K_R}{\eta_1}\Bigr)\Bigr],\\
&\LJ(U_0):=\exp\bigl(\Xfar(U_0)+B_5^*I_5\vz m^2\bigr)+\vz\kappa\TI .
\end{align*}
\endgroup
\namedlabel{hyp:P}{\textup{(P0)--(P8)}}\textbf{Side conditions} for $(p,U_0,\Cfl)$:
\begin{itemize}
\item[(P0)] $1/2<\gamma$, $\gamma+\nu<1$; $0<R\le1$; $m>0$; $0<s_0\le\fF<1$; $1<\kappa\le4/3$; $0<c_A\le1$; $K_f\ge3$,
$K_f+1\le K_R$; $\Keps\ge2$; $0<\eta_1\le K_f/(K_f+1)$; $\qc>1$; $\nu\ge0$; $\dL,\delta_1,A_0\ge1$; $0<\vz<1$;
\item[(P1)] $1/(1-\vz)\le\dL$ and $e^{2\vz/(1-\vz)}\le\delta_1$;
\item[(P2)] $c_\rho^*K_f\ge2$, $c_R^*\Keps\ge1$, $c_6\le K_R\Keps$, $c_\eta>0$;
\item[(P3)] $\camp\vz\le\nu$;
\item[(P4)] $\vz\bigl[2\aD(K_1)\dL/(c_\rho^*K_f)+\aD(K_2)J_1(U_0)\bigr]\le\log A_0$;
\item[(P5)] $A_0\Bigl[s_0+\Bigl(\dfrac{K_R}{K_f}\Bigr)^{\!\nu}\dfrac{\vz}m\bigl(B_{2U}'RJ_2(U_0)+B_{2\eta}^*\bigr)+B_3'\vz\Bigr]\le\fF$;
\item[(P6)] $\bar\theta<1$;\qquad (P7) $\LJ(U_0)<\kappa$;\qquad (P8) $m^2r_*<4c_A$.
\end{itemize}
\end{definition}
By (P2), $c_\rho^*>0$ and $c_R^*>0$; so $\rho_{\Cfl}\ge c_R^*R>0$, and $J_k$ is non-decreasing in $U_0$. By (P1), $1/(1-v)\le\dL$
and $e^{2\mu}\le\delta_1$. Note $c_6\ge2$, $K_R\ge4$, that $\aD$ is decreasing, and that by (P8) $\omega_r$ is defined and
increasing for $1\le r\le r_*$, and $\mathsf s$ is increasing.

\medskip\noindent\namedlabel{hyp:end}{(H)}\textbf{Setting (H).} The working front on $[0,U]$ and the chase cut at $b_{\mathrm c}=U$;
a non-decreasing function $\Cfl:[0,\infty)\to[0,\infty)$;
$(p,U_0,\Cfl)$ satisfies (P0)--(P8); $U\le U_0$; $\eps\le R/\Keps$; an \emph{excuse length} $\We\ge0$ with $\We=0$ or
$U\ge2\We+R+\eps$; and a number $\delta_t\ge0$ (in the application, $\delta_t=e^{\vz K_w}-1$ of
Section~\ref{sec:dep}). (S1), (S1G) and (S2) are hypotheses at the time $\ttd$. A pair $(a,b)$, $0\le a<b\le U$, is a \emph{body pair} if $U-b\ge\We$, a \emph{tail pair} if $U-b<\We$, and
\emph{good} if $a\ge\We$ or $U-b\ge\We$. A pair that is not good has $\lambda>U-2\We\ge R+\eps$; so \emph{every pair with
$\lambda\le R$ is good}. The \emph{family of leg $j+1$} consists of the pairs with $\eps/2\le\lambda\le R$ and ($a=a_j$ or
$a\ge a_{j+1}$); the \emph{family} $\Fam(t)$ is the family of leg $j(t)+1$ (so $\Fam(t_j)$, $j\ge1$, is the family of leg $j$, which
contains that of leg $j+1$). The \emph{gap} of a pair of $\Fam(t)$ is $g:=a-a_{j(t)}\ge0$; the pair is \emph{far} if $g\ge K_f\lambda$
and \emph{near} otherwise. Let $\jW$ be the first $j$ with $a_j\ge\We$ and $\tW:=t_{\jW}$ (if the chase reaches it). The
\emph{admissible family} $\FamS(t)$ consists of the body pairs of $\Fam(t)$ if $t<\tW$ (or $\tW$ is not reached), and of all pairs
of $\Fam(t)$ if $t\ge\tW$. Hypotheses:
\begin{itemize}
\item[(S1)] $c(a,b;\ttd)\ge c_A\min(\lambda,K_RR)$ for all good pairs with $\lambda\ge\eps/2$;
\item[(S1G)] $c(a,b;\ttd)\ge\Cfl\bigl(\min(\lambda,a+U-b)\bigr)$ for all pairs with $\lambda>K_RR$;
\item[(S2)] $N(a,b;\ttd)\le s_0M\lambda^{1/2+\gamma}\min\bigl(1,K_f\lambda/\min(a,K_RR)\bigr)^\nu$ for all body pairs with
$\eps/2\le\lambda\le R$ (the last factor is $1$ if $a=0$);
\item[(TW)] (only if $\We>0$) if $\tW$ is reached, every tail pair with $\eps/2\le\lambda\le R$ satisfies, with $g_W:=a-a_{\jW}$,
\[
g_W\ge K_f\lambda,\quad l_{\tW}(a,b)\le(1+\delta_t)\lambda,\quad
N(a,b;\tW)\le s_0M\lambda^{1/2+\gamma}\min\Bigl(1,\frac{K_f\lambda}{\min(g_W,K_RR)}\Bigr)^{\nu};
\]
\item[(PT)] (only if $\We>0$) $(1+\delta_t)\LJ(U_0)<\kappa$.
\end{itemize}
If $\We=0$, then $\jW=0$, $\tW=\ttd$, there are no tail pairs, $\FamS=\Fam$, and (TW), (PT) are void. Families shrink in $t$ (a pair
of the family of leg $k+1$ has $a\ge a_k\ge a_{j+1}$ for $k>j$), and the gap of a fixed pair is non-increasing in $t$.

\begin{theorem}[End theorem]\label{thm:end}
In the setting \textup{(H)}, the proviso holds at every non-final leg start, and:
\begin{enumerate}
\item[(a)] \emph{(invariants)} for every $t\in[\ttd,\Tstop]$ and every pair of $\FamS(t)$,
\begin{equation}\label{eq:J}
\textup{(J1)}\quad l_t(a,b)\le\kappa\lambda,\qquad\qquad \textup{(J2)}\quad N(a,b;t)\le M\lambda^{1/2+\gamma};
\end{equation}
if $\We>0$, then $\tW$ is reached. Moreover, for all good pairs with $\lambda\ge\eps/2$ and all $t\ge\ttd$,
\begin{equation}\label{eq:CA}
\textup{(CA)}\qquad c(a,b;t)\ge c_A\min(\lambda,K_RR).
\end{equation}
\item[(b)] \emph{(the end)} every non-final leg advances the base by $a_{j+1}-a_j\in[\eps/\kappa,\eps]$; the chase stops at
$\Tstop\le\ttd+\kappa U/(1-v)$; the direction of every true label of $[0,U]$ is seen by $\Tstop$.
\item[(c)] \emph{(tilts)} at every $t\in[\ttd,\Tstop]$, for pairs $J\subset I$ of $\FamS(t)$: $|e_J-e_I|\le\CT m(\lambda_I/R)^\gamma$.
\item[(d)] \emph{(distance from the base)} for $t$ in leg $j+1$ and a label $s\ge a_j$ with ($U-s\ge\We$ or $a_j\ge\We$):
$\rho_s(t)\ge c_A\min(s-a_j,K_RR)-\delta_1\eps$ and $\max(\rho_s,\eps)\ge\min(s-a_j,K_RR)/c_6$; and, under \textup{(S1G)}, if
$s-a_j>K_RR$: $\rho_s(t)\ge\max\bigl(c_AK_RR,\Cfl(\min(s-a_j,a_j+U-s))\bigr)-\delta_1\eps$.
\item[(e)] \emph{(causality)} if a path agrees with the chase on $[0,t_j]$, then \textup{(a)}, \textup{(c)}, \textup{(d)} hold for it on
$[\ttd,t_j]$.
\end{enumerate}
No constant depends on $\eps$, $v$ or $U$.
\end{theorem}

The proof occupies the rest of this section and Section~\ref{sec:farnear}; it is completed in Section~\ref{sec:endproof}.

\subsection{Legs, gaps, distances, tilts}

\emph{Proof of \eqref{eq:CA}.} Chords are non-decreasing in $t$ (Lemma~\ref{lem:sep}(c)); apply (S1). In particular \eqref{eq:CA}
holds for every pair with $\eps/2\le\lambda\le R$ (such pairs are good, and $R\le K_RR$). \qed

\medskip
Let $T_{\mathrm{def}}\le\infty$ be the supremum of the times up to which the chase is defined (it is defined on $[\ttd,t_j]$ up to
each leg start $t_j$ reached).

\begin{lemma}[Legs]\label{lem:legs}
Let $t'\le T_{\mathrm{def}}$, and suppose \textup{(J1)} holds at every leg start $t_i\le t'$ for the pairs of $\FamS(t_i)$. Then for
every leg start $t_j\le t'$:
\begin{enumerate}
\item[(i)] if $U-a_j\ge\eps/2$, leg $j+1$ is final, or $a_{j+1}-a_j\in[\eps/\kappa,\eps]$ and $l_{t_j}(a_j,a_{j+1})=\eps$;
\item[(ii)] if $U-a_j<\eps/2$, leg $j+1$ is final; if $a_j<\We$, leg $j+1$ is not final;
\item[(iii)] if leg $j+1$ is not final, $c(a_j,a_{j+1};t_j)\ge c_A(a_{j+1}-a_j)>0$: the proviso holds;
\item[(iv)] at most $\kappa U/\eps$ non-final legs start in $[\ttd,t']$, each lasts at most $\dL\eps$, and every $t\le t'$ satisfies
$t-\ttd\le T_B:=\chig U$.
\end{enumerate}
The pairs used are $(a_j,s)$ with $a_j<s\le a_j+\eps$, and $(a_{j-1},U)$ in case \textup{(ii)}; they belong to $\FamS(t_j)$.
\end{lemma}

\begin{proof}
Always $a_{j+1}-a_j\le l_{t_j}(a_j,a_{j+1})\le\eps$, since $w\ge1$. \emph{Admissibility.} If $t_j<\tW$ then $a_j<\We$, and every
$s\le a_j+\eps$ has $U-s>U-\We-\eps\ge\We+R$: the pairs $(a_j,s)$ are body pairs; and $U-a_j>\eps$. If $t_j\ge\tW$,
$\FamS(t_j)=\Fam(t_j)$. Every pair $(a_j,s)$ with $\eps/2\le s-a_j\le R$ is in $\Fam(t_j)$ (for $j=0$ it has first point $a_0$; for
$j\ge1$, $a_j\ge a_j$ in the family of leg $j$).
(i) The label $s:=a_j+\eps/2$ is $\le U$, and $\eps/2\le R$ ($\Keps\ge2$); (J1) gives $l_{t_j}(a_j,s)\le\kappa\eps/2\le\eps$, so
$a_{j+1}\ge s$. If leg $j+1$ is non-final, $a_{j+1}<U$ and $l_{t_j}(a_j,a_{j+1})=\eps$ by maximality; the pair $(a_j,a_{j+1})$ has
$\eps/2\le\lambda\le\eps\le R$, and (J1) gives $\eps\le\kappa\lambda$.
(ii) A final leg $j+1$ has $U-a_j\le l_{t_j}(a_j,U)\le\eps$, so $a_j\ge U-\eps>\We$ if $\We>0$: legs with $a_j<\We$ are not final.
Now let $U-a_j<\eps/2$; then $t_j\ge\tW$. Induction on $j$. For $j=0$, $l_{t_0}(0,U)=U<\eps/2$. For $j\ge1$, leg $j$ was not final, so
$U-a_{j-1}\ge\eps/2$ by the induction hypothesis, and by (i) $\lambda_1:=a_j-a_{j-1}\in[\eps/\kappa,\eps]$ with
$l_{t_{j-1}}(a_{j-1},a_j)=\eps$. With $\lambda_2:=U-a_j<\eps/2$, the pair $(a_{j-1},U)$ has label length in $[\eps/2,3\eps/2]\subset
[\eps/2,R]$ and lies in $\Fam(t_j)$ (family of leg $j$). By (J1) and the monotonicity of lengths,
$l_{t_j}(a_j,U)\le\kappa(\lambda_1+\lambda_2)-\eps<\frac32\kappa\eps-\eps\le\eps$, as $\kappa\le4/3$; so $a_{j+1}=U$.
(iii) By (i), $\lambda:=a_{j+1}-a_j\in[\eps/2,R]$, and \eqref{eq:CA} gives $c\ge c_A\lambda$.
(iv) Each non-final leg starting in $[\ttd,t']$ advances the base by at least $\eps/\kappa$ (by (i), (ii)), and the advances add up to
at most $U$. Each leg lasts at most $\eps/(1-v)\le\dL\eps$ (Lemma~\ref{lem:chase}(1), (P1)). So $t-\ttd\le(\kappa U/\eps)\dL\eps=T_B$.
\end{proof}

\begin{lemma}[Gap--time integrals]\label{lem:gapint}
Let $t'$ be as in Lemma~\ref{lem:legs}, $a$ a label, and $\mathcal T\subset[\ttd,t']$ a set of times on which
$g(\tau):=a-a_{j(\tau)}\ge0$. The gap is constant on each leg and, by Lemma~\ref{lem:legs}(i), (ii), its values on distinct legs
differ by at least $\eps/\kappa$.
\begin{enumerate}
\item[(i)] If $h\ge0$ is non-increasing on $[g_{\min},\infty)$, where $g_{\min}$ is the smallest value $\ge g_0$ of $g$ on
$\mathcal T$, then
\[
\textstyle\int_{\{\tau\in\mathcal T:\,g(\tau)\ge g_0\}}h(g(\tau))\,d\tau\le\dL\eps\,h(g_{\min})+\chig\int_{g_{\min}}^\infty h(g)\,dg .
\]
\item[(ii)] If $h\ge0$ is non-decreasing on $[0,\bar g+\eps/\kappa]$ and $g\le\bar g$ on $\mathcal T$, then
\[
\textstyle\int_{\mathcal T}h(g(\tau))\,d\tau\le\chig\int_0^{\bar g+\eps/\kappa}h(g)\,dg .
\]
\item[(iii)] At most $\kappa x/\eps+1$ legs meet $\{\tau\in\mathcal T:g(\tau)<x\}$.
\item[(iv)] If $g(\tau)\in(g_0,Y]$ on $\mathcal T$, and $H\ge0$ is non-increasing on $(0,Y/2]$ and non-decreasing on $[Y/2,\infty)$, with
$H\le H_{\max}$ on $(g_0,Y]$, then \textup(an integral over an empty range is $0$\textup)
\[
\textstyle\int_{\mathcal T}H(g(\tau))\,d\tau\le2\dL\eps H_{\max}+\chig\bigl(\int_{g_0}^{Y/2}H(g)\,dg+\int_{Y/2}^YH(g)\,dg\bigr).
\]
\end{enumerate}
\end{lemma}

\begin{proof}
The bases of non-final legs increase by at least $\eps/\kappa$ (Lemma~\ref{lem:legs}), and each leg lasts at most $\dL\eps$.
(i) Let $g_1>\dots>g_n=g_{\min}$ be the values $\ge g_0$. For $i<n$, $g_i-\eps/\kappa\ge g_{i+1}$ and, $h$ being non-increasing,
$h(g_i)\le(\kappa/\eps)\int_{g_i-\eps/\kappa}^{g_i}h$; these intervals are disjoint and lie in $[g_{\min},\infty)$. So the integral is
$\le\dL\eps\sum_ih(g_i)\le\dL\eps\,h(g_{\min})+\dL\kappa\int_{g_{\min}}^\infty h$. (ii) Likewise $h(g_i)\le(\kappa/\eps)
\int_{g_i}^{g_i+\eps/\kappa}h$ for every value $g_i$, with disjoint intervals in $[0,\bar g+\eps/\kappa]$. (iii) The values in $[0,x)$
are spaced by at least $\eps/\kappa$. (iv) Let $g_1>\dots>g_n$ be the values of $g$ on $\mathcal T$. For each value $\le Y/2$ except
the smallest, $H(g_i)\le(\kappa/\eps)\int_{g_i-\eps/\kappa}^{g_i}H$ ($H$ is non-increasing on $(0,Y/2]$ and
$g_i-\eps/\kappa\ge g_{i+1}>g_0$); these intervals are disjoint and lie in $(g_0,Y/2]$. For each value $>Y/2$ except the largest,
$H(g_i)\le(\kappa/\eps)\int_{g_i}^{g_i+\eps/\kappa}H$ ($H$ is non-decreasing on $[Y/2,\infty)$ and $g_i+\eps/\kappa\le g_{i-1}\le Y$); these
intervals are disjoint and lie in $[Y/2,Y]$. The two exceptional values contribute at most $2\dL\eps H_{\max}$, the others at most
$\dL\eps(\kappa/\eps)(\int_{g_0}^{Y/2}H+\int_{Y/2}^YH)$.
\end{proof}

\begin{lemma}[Distance from the base]\label{lem:dist}
Let $t\le T_{\mathrm{def}}$ lie in leg $j+1$ \textup($t\in[t_j,t_{j+1}]$, or $t=t_j=\Tstop$\textup), and let $s\ge a_j$ be a label
with $U-s\ge\We$ or $a_j\ge\We$. Then
\begin{enumerate}
\item[(i)] $\rho_s(t)\ge c_A\min(s-a_j,K_RR)-\delta_1\eps$;
\item[(ii)] $\max(\rho_s(t),\eps)\ge\min(s-a_j,K_RR)/c_6$;
\item[(iii)] under \textup{(S1G)}, if $w:=s-a_j>K_RR$: $\rho_s(t)\ge\max\bigl(c_AK_RR,\Cfl(\min(w,a_j+U-s))\bigr)-\delta_1\eps$.
\end{enumerate}
We call (i) the first and (ii) the second form.
\end{lemma}

\begin{proof}
Let $w:=s-a_j$. If $w<\eps/2$, the first right side is negative ($\delta_1\ge1>c_A/2$). Otherwise $(a_j,s)$ is good, and by
\eqref{eq:CA} and $|P-x_{a_j}|\le e^{2\mu}\eps\le\delta_1\eps$ (Lemma~\ref{lem:chase}(3), (P1); at $t=t_j$, $P=x_{a_j}$):
$\rho_s\ge c(a_j,s;t)-|P-x_{a_j}|\ge c_A\min(w,K_RR)-\delta_1\eps$; here $c(a_j,s;t)\ge c(a_j,s;\ttd)$ (Lemma~\ref{lem:sep}(c)). Put
$y:=\min(w,K_RR)$. If $c_Ay-\delta_1\eps\ge y/c_6$ we are done; otherwise, as $c_A-1/c_6=c_A\delta_1/(1+\delta_1)$,
$y\,c_A\delta_1/(1+\delta_1)<\delta_1\eps$, i.e.\ $y<c_6\eps$. (iii) As in (i), $\rho_s\ge c(a_j,s;\ttd)-\delta_1\eps$, and
$c(a_j,s;\ttd)\ge\Cfl(\min(w,a_j+U-s))$ by (S1G) for the pair $(a_j,s)$.
\end{proof}

\begin{lemma}[Tilts]\label{lem:tilt}
Let $t\le T_{\mathrm{def}}$ lie in leg $j+1$, let $I=(a,b)$ be a pair of the family of leg $j+1$, and suppose that \textup{(J1)},
\textup{(J2)} hold at time $t$ for every pair of the family of leg $j+1$ contained in $[a_j,\max(b,a_{j+1})]$. Then:
\begin{enumerate}
\item[(o)] for pairs $J\subset I$ of the family of leg $j+1$ with $\lambda_I\le r\lambda_J$ and $1\le r<4c_A/m^2$:
$|e_J-e_I|\le\omega_r\mathsf s(r)M\lambda_I^\gamma$;
\item[(i)] for pairs $J\subset I$ of the family of leg $j+1$: $|e_J-e_I|\le\CT M\lambda_I^\gamma=\CT m(\lambda_I/R)^\gamma$;
\item[(ii)] for $a<s\le b$ and $\lambda_s:=\max(s-a,\eps/2)$,
\begin{equation}\label{eq:Emu}
\textup{(E-$\mu$)}\quad\int_a^s|T-e_{ab}|^2ds\le M^2\lambda^{2\gamma}\min(\Cmu\lambda_s,\lambda),\quad
\int_a^s|T-e_{ab}|\,ds\le M\lambda^\gamma\bigl(\kappa\lambda_s\min(\Cmu\lambda_s,\lambda)\bigr)^{1/2}.
\end{equation}
\end{enumerate}
\end{lemma}

\begin{proof}
All pairs below lie in $[a_j,\max(b,a_{j+1})]$ and in the family of leg $j+1$, so (J1), (J2) and \eqref{eq:CA} apply to them.
(o) Let $J\ne I$ (otherwise there is nothing to prove). By Lemma~\ref{lem:split}(i), $Z^2\le N(I)^2f(l(J))$ with
$f(\ell):=\ell(1-\ell/l(I))$; $f$ increases in $l(I)$, and $l(I)\le\kappa\lambda_I$, $l(J)\le\kappa\lambda_J$ by (J1). So
$f(l(J))\le\max_{0<\ell\le\kappa\lambda_J}\ell(1-\ell/(\kappa\lambda_I))=\kappa\lambda_J\tilde\rho(\rho)$, where $\rho:=\lambda_I/\lambda_J$
and $\tilde\rho(\rho)=\rho/4$ for $\rho\le2$, $\tilde\rho(\rho)=1-1/\rho$ for $\rho\ge2$. By (J2) and \eqref{eq:CA} ($c(J)\ge c_A\lambda_J$),
$Z/c(J)\le M\lambda_I^{1/2+\gamma}(\kappa\lambda_J\tilde\rho)^{1/2}/(c_A\lambda_J)=(\kappa^{1/2}/c_A)(\rho\tilde\rho(\rho))^{1/2}M\lambda_I^\gamma
=\mathsf s(\rho)M\lambda_I^\gamma\le\mathsf s(r)M\lambda_I^\gamma$, since $\rho\tilde\rho(\rho)$ is $\rho^2/4$ resp.\ $\rho-1$, $\mathsf s$ is increasing and
$\rho\le r$. Further $N(I)^2/(4c(J))\le M^2\lambda_I^{1+2\gamma}/(4c_A\lambda_J)\le m^2r/(4c_A)<1$ ($M\lambda_I^\gamma\le MR^\gamma=m$),
and Lemma~\ref{lem:split}(ii) gives the claim.

(i) Let $J=(a',b')\subset I$, $q:=\qc$, $\lambda_k:=\lambda_Iq^{-k}$ and $K:=\min\{k\ge1:\lambda_k\le\lambda_J\}$.
\emph{Case A: $a\ge a_{j+1}$, or $a'=a$, or $\lambda_I(1-1/q)\ge\eps$.} For $0\le k\le K-1$ put $\alpha_k:=\min(a',b-\lambda_k)$,
$I_k:=(\alpha_k,\alpha_k+\lambda_k)$, and $I_K:=J$. (1) $J\subset I_k\subset I$: $\alpha_k\ge\min(a',b-\lambda_I)=a$ and $\alpha_k\le a'$; the
right end is $\min(a'+\lambda_k,b)\le b$, and it is $\ge b'$ as $\lambda_k>\lambda_J$ for $k\le K-1$. (2) $\alpha_k$ is non-decreasing in $k$
and the right end non-increasing, so $I_{k+1}\subset I_k$. (3) Each $I_k$ is in the family: its length $\lambda_k\in(\lambda_J,\lambda_I]
\subset[\eps/2,R]$, and its first point is $a'$ (the first point of the family pair $J$) or $b-\lambda_k<a'$; in the second case $k\ge1$
($I_0=I$) and $b-\lambda_k\ge a+\lambda_I(1-1/q)$, which is $\ge a\ge a_{j+1}$ if $a\ge a_{j+1}$, and $\ge a_j+\eps\ge a_{j+1}$ if
$\lambda_I(1-1/q)\ge\eps$ ($a_{j+1}-a_j\le l_{t_j}(a_j,a_{j+1})\le\eps$, as $w\ge1$); if $a'=a$, then $\alpha_k=a$ for all $k$.
(4) $\lambda(I_k)/\lambda(I_{k+1})=q$ for $k+1\le K-1$, and $\lambda(I_{K-1})/\lambda_J\in[1,q]$. By (o) with $r=q$,
$|e_{I_{k+1}}-e_{I_k}|\le\omega_q\mathsf s(q)M\lambda_k^\gamma$ ($q\le r_*<4c_A/m^2$ by (P8)), and summing,
$|e_J-e_I|\le\omega_q\mathsf s(q)M\lambda_I^\gamma\sum_{k=0}^{K-1}q^{-k\gamma}<\omega_{r_*}\sigma_q\mathsf s(q)M\lambda_I^\gamma\le\CT M\lambda_I^\gamma$.
\emph{Case B: $a=a_j$, $a'>a$ and $\lambda_I(1-1/q)<\eps$.} Then $\lambda_I/\lambda_J<(\eps q/(q-1))/(\eps/2)=2q/(q-1)\le r_*$, and
(o) with $r=2q/(q-1)$ gives $|e_J-e_I|\le\omega_{r_*}\mathsf s(2q/(q-1))M\lambda_I^\gamma\le\CT M\lambda_I^\gamma$. The two cases are exhaustive.

(ii) $J:=(a,a+\lambda_s)$ is a pair of the family (first point $a$, length in $[\eps/2,\lambda]$) contained in $I$. By
Lemma~\ref{lem:static}(b) with $\uhat=e:=e_{ab}$, (J2), (J1) and (i),
$\int_a^s|T-e|^2\le\int_J|T-e|^2=N(J)^2+c(J)|e_J-e|^2\le M^2\lambda_s^{1+2\gamma}+\kappa\lambda_s\CT^2M^2\lambda^{2\gamma}\le\Cmu M^2\lambda^{2\gamma}\lambda_s$;
also $\int_a^s|T-e|^2\le\int_I|T-e|^2=N(I)^2\le M^2\lambda^{1+2\gamma}$. The second bound follows by Cauchy--Schwarz, as
$l(a,s)\le l(J)\le\kappa\lambda_s$.
\end{proof}

\begin{lemma}[Locality]\label{lem:local}
Let $t\le T_{\mathrm{def}}$, $(a,b)\in\FamS(t)$, and $\ttd\le t'\le t$, with $t'>\tW$ if $(a,b)$ is a tail pair. Then
\begin{enumerate}
\item[(i)] every pair of $\Fam(t')$ contained in $[a_{j(t')},\max(b,a_{j(t')+1})]$ belongs to $\FamS(t')$;
\item[(ii)] every $s\in[a,b]$ satisfies the hypothesis of Lemma~\ref{lem:dist} at $t'$ with the base $a_{j(t')}$; in particular the
pair $(a_{j(t')},s)$ is good if $s-a_{j(t')}\ge\eps/2$;
\item[(iii)] the pairs used by Lemma~\ref{lem:legs} at a leg start $t_i$ belong to $\FamS(t_i)$;
\item[(iv)] if $\We>0$: a tail pair with $\lambda\le R$ has $a>a_{\jW}$; the pairs of $\Fam(\tW)$ with first point $a_{\jW-1}$ are body
pairs.
\end{enumerate}
For a tail pair at $t'=\tW$, \textup{(i)} and \textup{(ii)} hold with $j(t')$ replaced by $\jW$ \textup($\tW$ belongs to leg
$\jW+1$, and $\FamS(\tW)=\Fam(\tW)$ contains the family of leg $\jW+1$\textup).
\end{lemma}

\begin{proof}
Note $a\ge a_{j(t)}\ge a_{j(t')}$. (i) If $t'\ge\tW$, $\FamS(t')=\Fam(t')$. If $t'<\tW$, then $(a,b)$ is a body pair and
$a_{j(t')}<\We$, so $a_{j(t')+1}<\We+\eps$; a pair $(a'',b'')$ in the interval has $U-b''\ge\min(U-b,U-\We-\eps)\ge\We$. (ii) For a
body pair $U-s\ge U-b\ge\We$; for a tail pair $t'>\tW$ gives $a_{j(t')}\ge a_{\jW}\ge\We$. (iii) is shown in the proof of
Lemma~\ref{lem:legs}. (iv) A tail pair with $\lambda\le R$ has $a\ge b-R>U-\We-R\ge\We+\eps>a_{\jW}$, since $a_{\jW}\le a_{\jW-1}+\eps$.
A pair with first point $a_{\jW-1}<\We$ and $\lambda\le R$ has $U-b>U-\We-R\ge\We$. The last sentence: $a_{\jW}\ge\We$.
\end{proof}

\noindent\textbf{Locality principle.} The proofs of Propositions~\ref{prop:far}--\ref{prop:length} for a pair $(a,b)\in\FamS(t)$
use the family bounds (J1), (J2) only for the pairs of Lemma~\ref{lem:local}(i), (iii) (at the times $t'$ allowed there),
\eqref{eq:CA} only for these pairs and for the pairs $(a_{j(t')},s)$, $s\in[a,b]$, of Lemma~\ref{lem:local}(ii) (through
Lemma~\ref{lem:dist}), the hypothesis (S1G) at $\ttd$ only for these pairs $(a_{j(t')},s)$ (through Lemma~\ref{lem:dist}(iii)),
and, beyond the bounds of the pairs of Lemma~\ref{lem:local}(i) (which include $(a,b)$ itself), only the start data of $(a,b)$ at
$t_\circ$. Each invocation in Section~\ref{sec:farnear} is of this kind.

\begin{lemma}[Continuity scheme]\label{lem:cont}
Consider the statement, for a time $t$:
\begin{quote}
\textup{(Imp)} every pair of $\FamS(t)$ satisfies $l_t(a,b)\le(1+\delta)\LJ(U_0)\lambda$ and
$E(a,b;t)\le\frac12\bar\theta M^2\lambda^{1+2\gamma}$, where $\delta:=0$ for body pairs and $\delta:=\delta_t$ for tail pairs.
\end{quote}
Suppose that, whenever $t^*\le T_{\mathrm{def}}$ is such that the chase has not stopped before $t^*$ and \textup{(J1)}, \textup{(J2)}
hold for all $t\in[\ttd,t^*]$ and all pairs of $\FamS(t)$, \textup{(Imp)} holds at every $t\in[\ttd,t^*]$. Then the proviso holds
at every non-final leg start, the chase stops, \textup{(J1)}, \textup{(J2)} hold for all $t\in[\ttd,\Tstop]$ and all pairs of
$\FamS(t)$, and $\tW$ is reached if $\We>0$.
\end{lemma}

\begin{proof}
Let $\tau^*$ be the supremum of the $t'\le T_{\mathrm{def}}$ such that the chase has not stopped before $t'$ and (J1), (J2) hold on
$[\ttd,t']$ for $\FamS$. At $\ttd$: $l=\lambda<\kappa\lambda$, and $\FamS(\ttd)$ consists of body pairs ($a_0=0<\We$ if $\We>0$),
for which (S2) gives $N<M\lambda^{1/2+\gamma}$ ($s_0<1$). So $\tau^*\ge\ttd$.

\emph{The bounds hold at $\tau^*$ \textup($<\infty$\textup).} Let $\tau^*>\ttd$. If $\tau^*\in(t_j,t_{j+1}]$, then $\Fam(t)=\Fam(\tau^*)$
for $t\in(t_j,\tau^*)$, and $\FamS(t)=\FamS(\tau^*)$ unless $\tau^*=\tW$; as $l$, $E$ are continuous in $t$
(Lemma~\ref{lem:flow}(c)), (J1), (J2) hold at $\tau^*$. If $\tau^*=\tW$, the remaining pairs are the tail pairs, and (TW), (PT) give
$l\le(1+\delta_t)\lambda<\kappa\lambda$ ($\LJ\ge1$) and $N\le s_0M\lambda^{1/2+\gamma}<M\lambda^{1/2+\gamma}$. Hence (Imp) holds at
$\tau^*$, and by (P7), (PT), (P6) its bounds are strict: $l<\kappa\lambda$, $E<\frac12M^2\lambda^{1+2\gamma}$.

\emph{Extension.} Suppose the chase does not stop at $\tau^*$. If $\tau^*$ is a leg start $t_j$, Lemma~\ref{lem:legs} (with
$t'=\tau^*$) shows that leg $j+1$ is non-final and satisfies the proviso; so in all cases the chase is defined on
$[\tau^*,\tau^*+\delta_0]$ for some $\delta_0>0$, and we may take $\delta_0$ so small that $(\tau^*,\tau^*+\delta_0]$ contains no leg
start. For $t\in(\tau^*,\tau^*+\delta_0]$, $\FamS(t)\subset\FamS(\tau^*)$ (one leg; at a leg start the family contains that of the
next leg; $\tW\notin(\tau^*,\tau^*+\delta_0]$). For a fixed pair, $l'\le(2v/\eps)l$ and $|E'|=v|\int(T-e)\cdot AT|\le(4v/\eps)l$
(Lemma~\ref{lem:excess}(i), Lemma~\ref{lem:field}(ii)). With $\Lambda:=(1+\delta_t)\LJ(U_0)$ if $\We>0$, $\Lambda:=\LJ(U_0)$ otherwise,
choose $\delta\in(0,\delta_0]$ with $e^{2v\delta/\eps}\Lambda<\kappa$ and $\delta(4v/\eps)\kappa e^{2v\delta/\eps}\le\frac12(1-\bar\theta)M^2(\eps/2)^{2\gamma}$.
Then for $t\in[\tau^*,\tau^*+\delta]$ and every pair of $\FamS(t)$ ($\lambda\ge\eps/2$):
$l_t\le e^{2v\delta/\eps}\Lambda\lambda<\kappa\lambda$ and $E(t)\le\frac12\bar\theta M^2\lambda^{1+2\gamma}+\delta(4v/\eps)\kappa
e^{2v\delta/\eps}\lambda\le\frac12M^2\lambda^{1+2\gamma}$, so \textup{(J1)}, \textup{(J2)} hold on $[\ttd,\tau^*+\delta]$, contradicting the definition of $\tau^*$.

So the chase stops at $\tau^*$, unless $\tau^*=\infty$; but then (J1) holds at all leg starts, the chase never stops and has
infinitely many legs (Lemma~\ref{lem:chase}(0)), against Lemma~\ref{lem:legs}(iv). Hence $\tau^*=\Tstop<\infty$, the proviso held at
every non-final leg start (Lemma~\ref{lem:legs}(iii)), and (J1), (J2) hold on $[\ttd,\Tstop]$. Finally, if $\We>0$, the final leg
$J+1$ has $U-a_J\le l_{t_J}(a_J,U)\le\eps$, so $a_J\ge U-\eps>\We$, and $\jW\le J$.
\end{proof}

% sec4.tex -- writer W2. Section 4: far pairs, near pairs, length; end of the proof of Theorem thm:end.
\section{Far pairs, near pairs and length}\label{sec:farnear}

\noindent\textbf{Standing assumptions.} We are in the setting (H). Let $t^*\le T_{\mathrm{def}}$ be such that the chase has not
stopped before $t^*$ and (J1), (J2) hold for all $t\in[\ttd,t^*]$ and all pairs of $\FamS(t)$. Then Lemma~\ref{lem:legs} holds with
$t'=t^*$, and Lemma~\ref{lem:gapint} for sets of times in $[\ttd,t^*]$. Fix $t\in[\ttd,t^*]$ and $(a,b)\in\FamS(t)$, and put
$e:=e_{ab}$. For a tail pair at $t=\tW$ the bounds of (Imp) follow from (TW): $l\le(1+\delta_t)\lambda\le(1+\delta_t)\LJ(U_0)\lambda$
and $E\le\frac12s_0^2M^2\lambda^{1+2\gamma}\le\frac12\bar\theta M^2\lambda^{1+2\gamma}$; so for tail pairs we assume $t>\tW$.

The \emph{start time} is $t_\circ:=\ttd$ for body pairs and $t_\circ:=\tW$ for tail pairs. For $\tau\in[t_\circ,t]$ the \emph{base} is
$\beta(\tau):=a_{j(\tau)}$ and the \emph{next base} $\beta^+(\tau):=a_{j(\tau)+1}$, except that for a tail pair $\beta(\tW):=a_{\jW}$,
$\beta^+(\tW):=a_{\jW+1}$ ($\tW$ belongs to leg $\jW+1$); the gap is $g(\tau):=a-\beta(\tau)$, and $g:=g(t)$, $\hat g:=\min(g,K_RR)$.
Since $a\ge a_{j(t)}$, the pair belongs, for every $\tau\in[t_\circ,t]$, to the family of the leg with base $\beta(\tau)$ and to
$\FamS(\tau)$; for a tail pair at $\tW$ this uses $g_W\ge K_f\lambda>\eps\ge a_{\jW+1}-a_{\jW}$. The gap is non-increasing on
$[t_\circ,t]$, constant on each leg, and $g(t_\circ)=a$ for body pairs, $g(t_\circ)=g_W$ for tail pairs. By
Lemma~\ref{lem:local}, at every $\tau\in[t_\circ,t]$: Lemma~\ref{lem:dist} applies at every $s\in[a,b]$ with the base $\beta(\tau)$;
the pairs of the family of the leg of $\tau$ contained in $[\beta(\tau),\max(b,\beta^+(\tau))]$ are admissible, so they satisfy (J1),
(J2), and Lemma~\ref{lem:tilt} applies to them; by Lemma~\ref{lem:chase}(1), $P(\tau)$ lies on the segment
$[x_{\beta(\tau)},x_{\beta^+(\tau)}]$. Finally $\sigma_s:=s-a$, $w_s:=s-\beta(\tau)=g(\tau)+\sigma_s$ and $D(s):=\max(\rho_s,\eps)$.

\begin{lemma}[Two integrals]\label{lem:int}
For $\lambda>0$ and $z>0$ let
\[
h_1(z):=\int_z^\infty\min(\lambda,y)\,y^{-2}dy,\qquad h_3(z):=\int_z^\infty\min(\lambda,y)\,y^{-3}dy .
\]
Then
$h_1(z)=1+\log(\lambda/z)$ for $z\le\lambda$, $h_1(z)=\lambda/z$ for $z>\lambda$; $h_3(z)=1/z-1/(2\lambda)$ for $z\le\lambda$,
$h_3(z)=\lambda/(2z^2)$ for $z>\lambda$. Both are non-increasing; $z\,h_1(z)\le\lambda$; $\int_0^{K\lambda}h_1=\lambda(2+\log K)$ for
$K\ge1$; and $z\,h_3(z)=G_3(z/\lambda)$ with $G_3(x)=1-x/2$ for $x\le1$, $G_3(x)=1/(2x)$ for $x>1$.
\end{lemma}

\begin{proof}
Direct integration. For $x=z/\lambda\le1$, $z\,h_1(z)=\lambda x(1-\log x)\le\lambda$, as $x(1-\log x)$ increases on $(0,1]$ to $1$; for
$z>\lambda$, $z\,h_1(z)=\lambda$. Finally $\int_0^\lambda(1+\log(\lambda/z))\,dz=2\lambda$ and $\int_\lambda^{K\lambda}\lambda z^{-1}dz=\lambda\log K$.
\end{proof}

\begin{proposition}[Far pairs]\label{prop:far}
If $(a,b)$ is far at $t$, it is far on $[t_\circ,t]$, and $N(a,b;t)\le\fF M\lambda^{1/2+\gamma}$. Moreover, at every $\tau\in[t_\circ,t]$,
with $g=g(\tau)$: $\rho_m\ge c_\rho g$ if $g\le K_RR$, $\rho_m\ge c_RR$ if $g>K_RR$, and $|p_a|\le C_pMg^{1+\gamma}$ if $g\le\eta_1R$.
\end{proposition}

\begin{proof}
The gap is non-increasing, so the pair is far on $[t_\circ,t]$. Fix $\tau\in[t_\circ,t]$, $g=g(\tau)\ge K_f\lambda\ge K_f\eps/2>\eps$; so
$a>\beta^+(\tau)$.

\emph{Distance.} The arc has length $l\le\kappa\lambda$ (J1), so its convex hull lies in the ball of radius $l$ about $x_a$, and by
Lemma~\ref{lem:dist} at $s=a$: $\rho_m\ge\rho_a-l\ge c_A\min(g,K_RR)-\delta_1\eps-\kappa\lambda$. If $g\le K_RR$, then with $\eps\le2\lambda$
and $\lambda\le g/K_f$: $\rho_m\ge c_Ag-(2\delta_1+\kappa)\lambda\ge c_\rho g\ge c_\rho K_f\lambda$. If $g>K_RR$, then with
$\eps\le R/\Keps$ and $\lambda\le R$: $\rho_m\ge c_AK_RR-\delta_1R/\Keps-\kappa R=c_RR\ge c_R\lambda$. In both cases $\rho_m\ge\eps$, by
(P2) ($c_\rho K_f\eps/2\ge\eps$, resp.\ $c_RR\ge c_R\Keps\eps\ge\eps$), and $\rho_m\ge K_1l$, resp.\ $\rho_m\ge K_2l$. So
Lemma~\ref{lem:duhamel} applies on $[t_\circ,t]$, with $\alpha\le\aD(K_1)v/(c_\rho g)$ while $g\le K_RR$ and $\alpha\le\aD(K_2)v/(c_RR)$ while
$g>K_RR$.

\emph{Amplification.} For $t_\circ\le\tau\le t$:
$\exp\int_\tau^t\alpha\le A_0\bigl(\min(g(\tau),K_RR)/\hat g\bigr)^\nu$. Indeed, the times with $g>K_RR$ form an initial segment of
length $\le T_B=\chig U$ (Lemma~\ref{lem:legs}(iv)), on which $\int\alpha\le\aD(K_2)v\chig U/(c_RR)$. On the times of $[\tau,t]$ with
$g\le K_RR$ (present only if $g(t)\le K_RR$), Lemma~\ref{lem:gapint}(i) with $h(y)=y^{-1}\mathbf 1\{y\le\min(g(\tau),K_RR)\}$ and
$g_{\min}=g(t)\ge K_f\eps/2$ gives
\[
\int\alpha\le\frac{\aD(K_1)v}{c_\rho}\Bigl[\frac{\dL\eps}{g(t)}+\chig\log\frac{\min(g(\tau),K_RR)}{g(t)}\Bigr]
\le\frac{2\aD(K_1)\dL v}{c_\rho K_f}+\camp v\log\frac{\min(g(\tau),K_RR)}{g(t)} .
\]
As $v\le\vz$ and $U\le U_0$, (P4) and (P3) (the ratio is $\ge1$) give the claim.

\emph{Forcing} $\beta_{\mathrm f}:=2v\,l^{3/2}|p_a|/\rho_m^3$, where $|p_a|=\dist(P,\text{line}(x_a,x_b))$.
\emph{Case $g>\eta_1R$.} $|p_a|\le\rho_a\le\rho_m+l\le\rho_m(1+1/K_1)$ if $g\le K_RR$, resp.\ $\le\rho_m(1+1/K_2)$ if
$g>K_RR$ (as $\rho_m\ge K_1l$, resp.\ $\rho_m\ge K_2l$). So $\beta_{\mathrm f}\le2v\kappa^{3/2}\lambda^{3/2}(1+1/K_1)/(c_\rho g)^2$
if $\eta_1R<g\le K_RR$, and $\beta_{\mathrm f}\le2v\kappa^{3/2}\lambda^{3/2}(1+1/K_2)/(c_RR)^2$ if $g>K_RR$. The second regime lasts at
most $\chig U$. For the first, Lemma~\ref{lem:gapint}(i) with $g_0=\eta_1R$, $h(y)=y^{-2}\mathbf 1\{y\le K_RR\}$ gives
$\int g^{-2}d\tau\le\dL\eps/(\eta_1R)^2+\chig/(\eta_1R)\le(\dL/(\Keps\eta_1^2)+\chig/\eta_1)/R$. Hence
\[
\int_{\{g>\eta_1R\}}\beta_{\mathrm f}\,d\tau\le v\lambda^{3/2}\bigl[B_{2U}\chig U/R^2+B_{2\eta}/R\bigr].
\]
\emph{Case $g\le\eta_1R$.} As $P$ lies on $[x_{\beta},x_{\beta^+}]$ and the distance to a line is convex,
\[
|p_a|\le\max_{a'\in\{\beta,\beta^+\}}c(a',a)|e_{a'a}\times e|\le\max_{a'}c(a',a)\bigl(|e_{a'a}-e_{a'b}|+|e_{a'b}-e|\bigr).
\]
The pairs
$(a',a)$, $(a',b)$, $(a,b)$ lie in $[\beta,b]$ and belong to the family of the leg: their first points are $a'$ or $a$;
$\lambda(a',a)\ge g-\eps\ge(K_f-2)\lambda\ge\eps/2$; $\lambda(a',b)\le g+\lambda\le\eta_1R(1+1/K_f)\le R$ by (P0). Since
$\lambda(a',b)/\lambda(a',a)\le1+1/(K_f-2)=r_f$, the one-step bound of Lemma~\ref{lem:tilt} gives $|e_{a'a}-e_{a'b}|\le c_t(r_f)Mg^\gamma$;
same right end: $|e_{a'b}-e|\le C_BM(g+\lambda)^\gamma\le C_B(1+1/K_f)^\gamma Mg^\gamma$; and $c(a',a)\le\kappa g$. So
$|p_a|\le C_pMg^{1+\gamma}$, and with $\rho_m\ge c_\rho g$: $\beta_{\mathrm f}\le2v(\kappa\lambda)^{3/2}C_pMg^{1+\gamma}/(c_\rho g)^3
=B_3vM\lambda^{3/2}g^{\gamma-2}$.

\emph{Result.} By Lemma~\ref{lem:duhamel} on $[t_\circ,t]$ and the amplification bound (on $\{g\le\eta_1R\}$, $\min(g(\tau),K_RR)=g(\tau)$
and $\hat g=g$),
\[
N(t)\le A_0\Bigl(\frac{\min(g(t_\circ),K_RR)}{\hat g}\Bigr)^{\!\nu}N(t_\circ)+A_0\Bigl(\frac{K_RR}{\hat g}\Bigr)^{\!\nu}
\int_{\{g>\eta_1R\}}\beta_{\mathrm f}+A_0\int_{\{g\le\eta_1R\}}\Bigl(\frac{g(\tau)}{g}\Bigr)^{\!\nu}\beta_{\mathrm f}(\tau)\,d\tau .
\]
\emph{First term.} By (S2) ($g(t_\circ)=a$) resp.\ (TW) ($g(t_\circ)=g_W$), and as $g(t_\circ)\ge g>0$,
\[
N(t_\circ)\le s_0M\lambda^{1/2+\gamma}\min\bigl(1,K_f\lambda/\min(g(t_\circ),K_RR)\bigr)^\nu .
\]
The weight of $s_0M\lambda^{1/2+\gamma}$ is
$(\min(g(t_\circ),K_RR,K_f\lambda)/\hat g)^\nu$, which is $\le(K_f\lambda/g)^\nu\le1$ if $g\le K_RR$ and $\le1$ otherwise. So the term is
$\le A_0s_0M\lambda^{1/2+\gamma}$.
\emph{Second term.} $(K_RR/\hat g)^\nu\le(K_R/K_f)^\nu(R/\lambda)^\nu$ (if $g\le K_RR$ use $g\ge K_f\lambda$; otherwise the left side is $1$,
and $K_R>K_f$, $R\ge\lambda$). As $\gamma+\nu<1$ and $\lambda\le R$, $\lambda^{3/2}(R/\lambda)^\nu=\lambda^{1/2+\gamma}\lambda^{1-\gamma-\nu}R^\nu\le
\lambda^{1/2+\gamma}R^{1-\gamma}$. With $R^{-\gamma}=M/m$ the term is $\le A_0(K_R/K_f)^\nu(v/m)(B_{2U}\chig U/R+B_{2\eta})M\lambda^{1/2+\gamma}$.
\emph{Third term} (present only if $g\le\eta_1R$). It is $\le A_0B_3vM\lambda^{3/2}g^{-\nu}\int g(\tau)^{\nu+\gamma-2}d\tau$, and
Lemma~\ref{lem:gapint}(i) ($h(y)=y^{\nu+\gamma-2}$, $g_{\min}=g$) gives $\int\le\dL\eps g^{\nu+\gamma-2}+\chig g^{\nu+\gamma-1}/(1-\gamma-\nu)$. With
$g\ge K_f\lambda$, $\eps\le2\lambda$ and $\gamma<1$ the term is
$\le A_0B_3K_f^{\gamma-1}\bigl(2\dL/K_f+\chig/(1-\gamma-\nu)\bigr)vM\lambda^{1/2+\gamma}$.
With $v\le\vz$, $U\le U_0$, (P5) gives $N\le\fF M\lambda^{1/2+\gamma}$.
\end{proof}

\begin{proposition}[Near pairs]\label{prop:near}
If $(a,b)$ is near at $t$, then $E(a,b;t)\le\frac12\bar\theta M^2\lambda^{1+2\gamma}$.
\end{proposition}

\begin{proof}
\emph{Start value.} The gap is constant on legs and non-increasing, so the set of times of $[t_\circ,t]$ at which the pair is near
is either $[t_\circ,t]$ or $(t_n,t]$ for a leg start $t_n\ge t_\circ$ at which it is far. In the first case $t_\circ=\ttd$ (tail pairs
are far at $\tW$), and $E(\ttd)\le\frac12s_0^2M^2\lambda^{1+2\gamma}$ by (S2); we put $t_n:=\ttd$. In the second case
Proposition~\ref{prop:far} at time $t_n$ (or (TW) if $t_n=\tW$) gives $E(t_n)\le\frac12\fF^2M^2\lambda^{1+2\gamma}$. We bound the
increase on $(t_n,t]$, where $g<K_f\lambda$, using Lemma~\ref{lem:excess}(iii) with $\bar q=0$: $E'\le\mathrm I+\mathrm{II}$,
$\mathrm I:=v\int_a^ba^*(\rho)|T-e|^2ds$, $\mathrm{II}:=v\int_a^b|\vsig||q|\,ds$.

\emph{Term \textup{I}.} Since $a^*(\rho)=\eps^{-1}\mathbf 1\{\rho<\eps\}+\int_{\max(\rho,\eps)}^\infty y^{-2}dy$,
$\mathrm I=v[\omega(\eps)/\eps+\int_\eps^\infty\omega(y)y^{-2}dy]$ with $\omega(y):=\int_{\{s\in(a,b):\rho_s<y\}}|T-e|^2ds$. Let $y\ge\eps$.
If $(y+\delta_1\eps)/c_A\le K_RR$, then by Lemma~\ref{lem:dist} $\rho_s<y$ forces $c_A\min(w_s,K_RR)-\delta_1\eps<y$, i.e.\
$\min(w_s,K_RR)<(y+\delta_1\eps)/c_A\le K_RR$, hence $w_s<(y+\delta_1\eps)/c_A\le(1+\delta_1)y/c_A=c_6y$, so these $s$ lie in
$(a,a+\lambda_y)$, $\lambda_y:=\min(\lambda,\max(c_6y-g,\eps/2))\le\min(\lambda,c_6y)$, and \eqref{eq:Emu} gives
$\omega(y)\le\CmuA M^2\lambda^{2\gamma}\min(\lambda,c_6y)$. If $(y+\delta_1\eps)/c_A>K_RR$, then $c_6y>K_RR\ge\lambda$ and
$\omega(y)\le N^2\le M^2\lambda^{1+2\gamma}\le\CmuA M^2\lambda^{2\gamma}\min(\lambda,c_6y)$. Also $\omega(y)=0$ for $y\le\min\rho$.
The piece $\omega(\eps)/\eps$ is non-zero only if $\min_{[a,b]}\rho<\eps$: then $c_A\min(w_s,K_RR)-\delta_1\eps<\eps$ for some $s$, i.e.\
$\min(w_s,K_RR)<c_6\eps\le K_RR$ ((P2), $\eps\le R/\Keps$), so $g\le w_s<c_6\eps$; and it is $\le\CmuA M^2\lambda^{2\gamma}\min(\lambda,c_6\eps)/\eps$.
With $\bar\rho:=\max(\eps,\min\rho)$ the integral piece is
$\int_{\bar\rho}^\infty\omega y^{-2}dy\le\CmuA M^2\lambda^{2\gamma}c_6h_1(c_6\bar\rho)$ ($z=c_6y$). By the second form of
Lemma~\ref{lem:dist} ($w_s\ge g$), $c_6\bar\rho\ge\min(g,K_RR)\ge\min(g,K_R\lambda)$, and $c_6\bar\rho\ge c_6\eps$; so
$h_1(c_6\bar\rho)\le\bar h(g):=h_1(\max(c_6\eps,\min(g,K_R\lambda)))$, non-increasing in $g$.

\emph{Time integral of \textup{I}} (in units of $v\CmuA M^2\lambda^{2\gamma}$). By Lemma~\ref{lem:gapint}(iii) at most $\kappa c_6+1$ legs
have $g<c_6\eps$, each lasting $\le\dL\eps$: the first piece contributes $\le(\kappa c_6+1)\dL\min(\lambda,c_6\eps)\le\dL(\kappa c_6+1)\lambda$.
By Lemma~\ref{lem:gapint}(i) ($g_0=0$, $h=c_6\bar h\,\mathbf 1\{g<K_f\lambda\}$), Lemma~\ref{lem:int} and $K_f\le K_R$, the second piece
contributes $\le c_6\dL\eps\,h_1(c_6\eps)+c_6\chig\int_0^{K_f\lambda}h_1\le\dL\lambda+c_6\chig\PsiI\lambda$. Hence
$\int_{t_n}^t\mathrm I\le v\DI M^2\lambda^{1+2\gamma}$.

\emph{Term \textup{II}.} By Lemma~\ref{lem:field}(iii), $|q(s)|\le|p(s)|/D(s)^2$, $p(s):=(x_s-P)\times e=(x_s-x_a)\times e+p_a$, and
$|(x_s-x_a)\times e|\le\int_a^s|T-e|\le\ConeA M\lambda^\gamma\max(\sigma_s,\eps/2)$ by \eqref{eq:Emu}. As $P\in[x_\beta,x_{\beta^+}]$,
$|p_a|\le\max_{a'\in\{\beta,\beta^+\}}|(x_a-x_{a'})\times e|$. Let $\Psi(y):=\int_{\{s\in(a,b):w_s<y\}}|\vsig|\,ds$; it vanishes for
$y\le g$, and for $y\ge\eps/2$, by \eqref{eq:Emu} and $|\vsig|\le|T-e|$, $\Psi(y)\le\ConeA M\lambda^\gamma\min(\lambda,y)$.

\emph{Short case: $(K_f+1)\lambda\le R$.} Then $w_s<(K_f+1)\lambda\le K_RR$, so $D(s)\ge\max(w_s,c_6\eps)/c_6$ by Lemma~\ref{lem:dist}.

(i) \emph{Offset $|p_a|$.} \emph{Case $g=0$} ($a=\beta$; $x_a$ is on the line). If $\beta^+\le b$,
$|(x_{\beta^+}-x_a)\times e|\le\ConeA M\lambda^\gamma\max(\beta^+-\beta,\eps/2)\le\ConeA M\lambda^\gamma\eps$. If $\beta^+>b$, the pair
$I_0:=(\beta,\beta^+)$ ($\lambda_{I_0}\in[\eps/\kappa,\eps]$) belongs to the family and contains $(a,b)$ with the same left end, so
$|(x_{\beta^+}-x_a)\times e|\le c(I_0)|e_{I_0}-e|\le\kappa\eps\,C_AM\eps^\gamma\le\kappa2^\gamma C_AM\lambda^\gamma\eps$. So $|p_a|\le k_0M\lambda^\gamma\eps$.
\emph{Case $g>0$} ($a\ge\beta^+$, so $g\ge\eps/\kappa>\eps/2$). For $a'\in\{\beta,\beta^+\}$ let $g':=a-a'\le g$. If $g'=0$ the term
vanishes. If $g'\ge\eps/2$, the pairs $(a',a)$, $(a',b)$, $(a,b)$ are in the family ($\lambda(a',b)\le g+\lambda<(K_f+1)\lambda\le R$), and
$|(x_a-x_{a'})\times e|\le c(a',a)(|e_{a'a}-e_{a'b}|+|e_{a'b}-e|)\le\kappa g(C_A+C_B)M(g+\lambda)^\gamma$. If $0<g'<\eps/2$ (then
$a'=\beta^+$), $J_0:=(a',a'+\eps/2)\subset(a',b)$ is in the family and contains $(a',a)$; by Cauchy--Schwarz, Lemma~\ref{lem:static}(b)
and $|e_{J_0}-e|\le|e_{J_0}-e_{a'b}|+|e_{a'b}-e|\le\CTL M(g'+\lambda)^\gamma$,
\[
|(x_a-x_{a'})\times e|\le l(J_0)^{1/2}\bigl(N(J_0)^2+c(J_0)|e_{J_0}-e|^2\bigr)^{1/2}
\le\tfrac12\kappa^{1/2}\eps(1+\kappa\CTL^2)^{1/2}M(g+\lambda)^\gamma
\]
($\eps/2\le\lambda$). Since $\eps\le\kappa g$ and $\frac12\kappa^{1/2}(1+\kappa\CTL^2)^{1/2}\le\CTL$ (i.e.\ $\kappa\le\CTL^2(4-\kappa^2)$,
true as $\kappa\le4/3$ and $\CTL\ge c_t(3)\ge3^{1/2}$), this is $\le\kappa\CTL M(g+\lambda)^\gamma g$. So $|p_a|\le\kappa\CTL M(g+\lambda)^\gamma g$
whenever $g>0$.

(ii) \emph{Split.} As $\max(\sigma_s,\eps/2)\le c_6D(s)$,
$\mathrm{II}\le v\ConeA M\lambda^\gamma c_6\int_a^b|\vsig|/D\,ds+v|p_a|\int_a^b|\vsig|/D^2\,ds=:\mathrm{II}_\alpha+\mathrm{II}_\beta$. Here
$1/D\le\varphi_1(w_s)$, $\varphi_1(w):=c_6/\max(w,c_6\eps)=\int_{\max(w,c_6\eps)}^\infty c_6y^{-2}dy$, and $1/D^2\le\varphi_1(w_s)^2=
\int_{\max(w_s,c_6\eps)}^\infty2c_6^2y^{-3}dy$. By Fubini, $\int|\vsig|\varphi_1(w_s)\,ds=\int_{c_6\eps}^\infty c_6y^{-2}\Psi(y)\,dy$, and
similarly for $\varphi_1^2$. Put $y_0:=\max(c_6\eps,g)$ and $h_2(g):=h_1(y_0)$.
\begin{itemize}
\item $\mathrm{II}_\alpha\le vc_6^2\ConeA^2M^2\lambda^{2\gamma}\int_{y_0}^\infty\min(\lambda,y)y^{-2}dy=vc_6^2\ConeA^2M^2\lambda^{2\gamma}h_2(g)$.
\item $g=0$: $\Psi(y)\le\ConeA M\lambda^\gamma y$ and $\int_{c_6\eps}^\infty2c_6^2y^{-2}dy=2c_6/\eps$, so
$\mathrm{II}_\beta\le2c_6\ConeA k_0vM^2\lambda^{2\gamma}$.
\item $g>0$: with (i), $h_3(y_0)\le h_3(g)$ and Lemma~\ref{lem:int}, for $x:=g/\lambda$,
\[
\begin{aligned}
\mathrm{II}_\beta&\le v|p_a|\,2c_6^2\ConeA M\lambda^\gamma h_3(y_0)\le2c_6^2\kappa\CTL\ConeA vM^2\lambda^\gamma(g+\lambda)^\gamma g\,h_3(g)\\
&=2c_6^2\kappa\CTL\ConeA vM^2\lambda^{2\gamma}(1+x)^\gamma G_3(x).
\end{aligned}
\]
Here $(1+x)^\gamma G_3(x)\le H(x)$,
$H(x):=2^\gamma(1-x/2)$ for $x\le1$, $H(x):=(1+x)^\gamma/(2x)$ for $x>1$; $H$ is continuous, non-increasing ($(1+x)^\gamma/x$ decreases, as
$\gamma<1$), and $H\le2^\gamma$.
\end{itemize}
\emph{Time integrals.} $\mathrm{II}_\alpha$: Lemma~\ref{lem:gapint}(i) ($g_0=0$, $h=h_2\mathbf 1\{g<K_f\lambda\}$) and Lemma~\ref{lem:int} give
$\le\dL\eps h_1(c_6\eps)+\chig\int_0^{K_f\lambda}h_1\le\Ta\lambda$. $\mathrm{II}_\beta$ with $g=0$: one leg, of duration $\le\dL\eps\le2\dL\lambda$.
$\mathrm{II}_\beta$ with $g>0$: Lemma~\ref{lem:gapint}(i) ($g_0=\eps/\kappa$, $h=H(\cdot/\lambda)\mathbf 1\{g<K_f\lambda\}$) gives
$\le2^\gamma\dL\eps+\chig\lambda\int_0^{K_f}H\le\Tb\lambda$, using $\int_0^12^\gamma(1-x/2)dx=\frac342^\gamma$ and, on $[1,K_f]$,
$(1+x)^\gamma\le(2x)^\gamma$, $\int_1^{K_f}(2x)^{\gamma-1}dx=2^{\gamma-1}(K_f^\gamma-1)/\gamma$. In total
$\int_{t_n}^t\mathrm{II}\le v\DIIs M^2\lambda^{1+2\gamma}$.

\emph{Long case: $(K_f+1)\lambda>R$.} Now $|p(s)|\le\rho_s\le D(s)$, so $|q|\le1/D$; and $w_s<(K_f+1)\lambda\le(K_f+1)R\le K_RR$, so again
$1/D\le\varphi_1(w_s)$. Hence $\mathrm{II}\le v\int|\vsig|\varphi_1(w_s)\,ds\le vc_6\ConeA M\lambda^\gamma h_2(g)$, and as for
$\mathrm{II}_\alpha$, $\int_{t_n}^t\mathrm{II}\le vc_6\ConeA\Ta M\lambda^{1+\gamma}$. Since $M\lambda^\gamma=m(\lambda/R)^\gamma>m(K_f+1)^{-\gamma}$,
$M\lambda^{1+\gamma}\le((K_f+1)^\gamma/m)M^2\lambda^{1+2\gamma}$, and $\int_{t_n}^t\mathrm{II}\le v\DIIl M^2\lambda^{1+2\gamma}$.

\emph{Conclusion.} $E(t)\le\frac12\fF^2M^2\lambda^{1+2\gamma}+v[\DI+\max(\DIIs,\DIIl)]M^2\lambda^{1+2\gamma}\le\frac12\bar\theta M^2\lambda^{1+2\gamma}$,
as $s_0\le\fF$ and $v\le\vz$.
\end{proof}

\begin{proposition}[Length]\label{prop:length}
With $\delta:=0$ for body pairs and $\delta:=\delta_t$ for tail pairs,
\[
l_t(a,b)\le(1+\delta)\exp\bigl(\Xfar(U_0)+B_5I_5\vz m^2\bigr)\lambda+\vz\kappa\TI\lambda\le(1+\delta)\LJ(U_0)\lambda<\kappa\lambda .
\]
\end{proposition}

\begin{proof}
$l(t_\circ)\le(1+\delta)\lambda$ ($l(\ttd)=\lambda$; (TW)), and $l'=v\int T\cdot AT\le v\int a^*(\rho)\,ds$ (Lemma~\ref{lem:excess}(i),
Lemma~\ref{lem:field}(ii)). The pair is far on an initial segment of $[t_\circ,t]$ and near afterwards.

\emph{Far, $g>\eta_1R$.} By Proposition~\ref{prop:far}, $\rho\ge\rho_m\ge\eps$ on the arc, so $(\log l)'\le v/\rho_m$, which is $\le v/(c_RR)$
while $g>K_RR$ (for a time $\le\chig U$) and $\le v/(c_\rho g)$ while $\eta_1R<g\le K_RR$. Lemma~\ref{lem:gapint}(i) ($g_0=\eta_1R$,
$h(y)=y^{-1}\mathbf 1\{y\le K_RR\}$) gives $\int d\tau/g\le\dL\eps/(\eta_1R)+\chig\log(K_R/\eta_1)\le\dL/(\Keps\eta_1)+\chig\log(K_R/\eta_1)$.
So $\log l$ increases by at most $\Xfar(U_0)$ here.

\emph{Far, $g\le\eta_1R$.} By Lemma~\ref{lem:excess}(v), $l'\le(2v/\rho_m)[l(|p_a|+l^{1/2}N)^2/\rho_m^2+N^2]$, with $|p_a|\le C_pMg^{1+\gamma}$,
$\rho_m\ge c_\rho g$ (Proposition~\ref{prop:far}), $l^{1/2}N\le\kappa^{1/2}M\lambda^{1+\gamma}\le\kappa^{1/2}K_f^{-1-\gamma}Mg^{1+\gamma}$ and
$N^2\le M^2\lambda^{1+2\gamma}\le lM^2(g/K_f)^{2\gamma}$. So $(\log l)'\le B_5vM^2g^{2\gamma-1}$. These times form a final segment of the
far period with gaps $\le\eta_1R$; as $\gamma>1/2$, $y\mapsto y^{2\gamma-1}$ is non-decreasing, and Lemma~\ref{lem:gapint}(ii)
($\bar g=\eta_1R$) gives $\int g^{2\gamma-1}d\tau\le\chig(\eta_1R+\eps/\kappa)^{2\gamma}/(2\gamma)\le I_5R^{2\gamma}$ ($\eps\le R/\Keps$). So
$\log l$ increases by at most $B_5I_5vM^2R^{2\gamma}=B_5I_5vm^2$ here.

\emph{Near.} $l'\le v[\tilde\omega(\eps)/\eps+\int_\eps^\infty\tilde\omega(y)y^{-2}dy]$ with $\tilde\omega(y):=|\{s\in(a,b):\rho_s<y\}|$
(arclength). As in Term~I of Proposition~\ref{prop:near}, with (J1) for $(a,a+\lambda_y)$ in place of \eqref{eq:Emu},
$\tilde\omega(y)\le\kappa\min(\lambda,c_6y)$, and $\tilde\omega(\eps)>0$ only if $g<c_6\eps$. So the increase while near is
$\le v\kappa\TI\lambda$.

Hence $l_t\le(1+\delta)\lambda\exp(\Xfar(U_0)+B_5I_5\vz m^2)+\vz\kappa\TI\lambda\le(1+\delta)\LJ(U_0)\lambda$, which is $<\kappa\lambda$ by
(P7) resp.\ (PT).
\end{proof}

\subsection{Completion of the proof of Theorem~\ref{thm:end}}\label{sec:endproof}

\eqref{eq:CA} was proved in Section~\ref{sec:end}. Let $t^*$ be as in Lemma~\ref{lem:cont}, i.e.\ as in the standing assumptions, and
$t\in[\ttd,t^*]$. Every pair of $\FamS(t)$ satisfies the bounds of (Imp): for a tail pair at $t=\tW$ by (TW); otherwise $E$ is bounded by
Proposition~\ref{prop:far} ($\frac12\fF^2\le\frac12\bar\theta$) or Proposition~\ref{prop:near}, and $l$ by Proposition~\ref{prop:length}. In
these proofs every family bound and every use of \eqref{eq:CA} concerns a pair allowed by the locality principle, so they are covered by
the assumption on $[\ttd,t^*]$. Lemma~\ref{lem:cont} gives the proviso, the stop, (a), and that $\tW$ is reached if $\We>0$.

(b) By (a), (J1) holds at every leg start, so Lemma~\ref{lem:legs} applies with $t'=\Tstop=t_J$: every non-final leg advances the base by
$[\eps/\kappa,\eps]$, there are at most $\kappa U/\eps$ of them, and each lasts at most $\eps/(1-v)$ (Lemma~\ref{lem:chase}(1)); so
$\Tstop\le\ttd+\kappa U/(1-v)$. The intervals $[a_j,a_{j+1}]$, $j\le J$, cover $[0,U]$ ($a_{J+1}=U$), and Lemma~\ref{lem:chase}(2) gives
the claim on true labels.
(c) For $I=(a,b)\in\FamS(t)$, the pairs of $\Fam(t)$ contained in $[a_{j(t)},\max(b,a_{j(t)+1})]$ are admissible
(Lemma~\ref{lem:local}(i) with $t'=t$; for $t=\tW$, $\FamS(\tW)=\Fam(\tW)$), so they satisfy (J1), (J2) by (a), and Lemma~\ref{lem:tilt}(i)
applies.
(d) is Lemma~\ref{lem:dist}.
(e) Let $\tilde P=P$ on $[0,t_j]$. By Lemma~\ref{lem:chase}(4) the working fronts agree on $[\ttd,t_j]$, hence so do the legs, the bases,
the families and $\tW$ (if $\tW\le t_j$) up to $t_j$. The statements (a), (c), (d) at a time $t\le t_j$ involve only the front at time $t$
and these data, so they hold for $\tilde P$. \qed

% sec5.tex -- writer W3.
\section{Heading and departure}\label{sec:dep}

At touch-down the working front starts at $P(\ttd)=\xz(0)$ with tangent close to $\nhat$, and its other end $\xz(U)$ lies just
behind $P$, at $P(\ttd)-z_0\nhat-y_0\nhat^\perp$ (Theorem~\ref{thm:touchdown}). This section shows that the first legs of the
chase run nearly in the direction $\nhat$ (Lemma~\ref{lem:heading}) and that, while they do, the arc near the far end recedes from
$P$ and stays almost straight (Theorem~\ref{thm:departure}). This provides the start data (TW) of the tail pairs in
Theorem~\ref{thm:end}.

Throughout, $d=1$, $0<v\le\vz$, $0<\eps\le\epsS$, and the parameter tuple $\pA$ and the numbers $\We$, $\UA$, $\UE$, $G_f$,
$\Dbar$, $z_0^{\max}$, $X_*$, $\zeta$, $\etaD$, $\etaN$, $\theta_0$, $\theta_y$, $\kD$ are those of \eqref{eq:asm} in
Section~\ref{sec:proof}; $\bar\omega$, $c_k$ and $\delta_0=0.251113\,R$ (Table~\ref{tab:asm}) are free parameters. All values are in Appendix~\ref{app:const}. The two
results below use only the inequalities among these numbers collected in (A3), (A4) of Section~\ref{sec:asmconst}.

\subsection{Heading of the first legs}

The \emph{stage-A chase} is a chase (Definition~\ref{def:chase}) on a working front, cut at the label $\UA$, for which the
conclusions (a)--(d) of Theorem~\ref{thm:end} hold with the tuple $\pA$ and excuse length $0$ on $[\ttd,t_j]$ for the leg starts
$t_j$ considered: exponent $\gD$, $M_{\mathrm A}=mR^{-\gD}$, tilt constant $\CT^{\mathrm A}$. With excuse
length $0$ every pair is good and $\FamS=\Fam$. We use the \emph{heading constants}
\begin{equation}\label{eq:dep-headconst}
\begin{gathered}
\bar\sigma:=\theta_{\mathrm v}G_f^{1-\mub},\qquad\theta_G:=\theta_{\mathrm v}G_f^{1-\mub-\gD},\qquad k_\sigma:=\frac{1+\bar\sigma/2}{1-\bar\sigma^2/2},\qquad
C_3:=2\CT^{\mathrm A},\\
c_\tau:=\frac1{1-m^2/2},\qquad\rho_*:=c_A(R-\delta_0-\epsS)-\delta_1\epsS,\qquad
K_{\mathrm{rot}}:=\frac{\kappa}{1-\vz}\Bigl[\frac{2R}{c_A\delta_0}+\frac{\We+\epsS}{\rho_*}\Bigr],\\
\omega_0:=\CT^{\mathrm A}m\Bigl(\frac{\delta_0}R\Bigr)^{\gD}+k_\sigma\bar\sigma+\vz K_{\mathrm{rot}},\qquad
\theta_{h1}:=\frac{\CT^{\mathrm A}m}{R^{\gD}}+2^{\gD}k_\sigma\theta_G+\frac{2^{\gD+2}\vz\kappa C_3m}{c_A(1-\vz)R^{\gD}},\\
\vartheta_h:=c_\tau^{\gD}\max\bigl(\theta_{h1},\,\omega_0\delta_0^{-\gD}\bigr).
\end{gathered}
\end{equation}
and the \emph{heading bound} $\hH(\toff):=\min\bigl(\vartheta_h(\toff+\eps)^{\gD},\,\omega_0\bigr)$, $\toff\ge0$.

\begin{lemma}[Heading]\label{lem:heading}
Consider the stage-A chase, and let $\nhat$ be a unit vector with $|T(u,\ttd)-\nhat|\le\theta_{\mathrm v}u^{1-\mub}$ for $u\le G_f$. Assume
\textup{(A3)}. Let $j$ be such that $a_j<\We$ and leg $j+1$ is non-final, with direction $\uhat_{j+1}$. Then for all $t$ in leg
$j+1$, with $\toff:=t-\ttd$,
\[
|\uhat_{j+1}-\nhat|\le\hH(\toff).
\]
\end{lemma}

\begin{proof}
For $s\ge0$ let $\psi_\star(s):=\min\bigl(1,s/(1-s^2/2)\bigr)$ for $0\le s\le1$ and $\psi_\star(s):=1$ for $s>1$; then
$\psi_\star(s)\le\min(1,2s)$, as $1-s^2/2\ge\frac12$ for $s\le1$. Put $\Lambda:=a_j+\eps<\We+\eps$, $\lambda_I:=\min(\Lambda,\delta_0)\in[\eps,\delta_0]$ (as
$\eps\le\epsS\le\delta_0$), $b:=a_j+\lambda_I$ and $I:=(a_j,b)$. Then $b<\We+\delta_0\le G_f\le\UA$, and $I$ is a family pair at all
times of legs $i+1\le j+1$ (its first point is $a_j\ge a_{i+1}$ or is the base, and $\eps/2\le\lambda_I\le R$).

\emph{Step 1 (time $t_j$).} $J:=(a_j,a_{j+1})$ has label length in $[\eps/\kappa,\eps]\subset[\eps/2,R]$ (Theorem~\ref{thm:end}(b),
$\kappa\le2$), so it is a family pair, $J\subset I$ with the same left end, and $e_J(t_j)=\uhat_{j+1}$ since $P(t_j)=x_{a_j}(t_j)$
(Definition~\ref{def:chase}). By Theorem~\ref{thm:end}(c), $|\uhat_{j+1}-e_I(t_j)|\le\CT^{\mathrm A}m(\lambda_I/R)^{\gD}$.

\emph{Step 2 (time $\ttd$).} Let $\sigma:=\sup_{u\in I}|T(u,\ttd)-\nhat|$. As $b\le G_f$ and $\gD\le1-\mub$,
$\sigma\le\theta_{\mathrm v}b^{1-\mub}\le\min(\bar\sigma,\theta_Gb^{\gD})$, and $b\le a_j+\Lambda\le2\Lambda$. At $\ttd$ the labels are
arclength, so $l(I)=\lambda_I$; from $c\,e_I-l\,\nhat=\int_I(T-\nhat)\,ds$ and $c\ge\int_IT\cdot\nhat\,ds\ge l(1-\sigma^2/2)$ we get
$c|e_I-\nhat|\le l\sigma+(l-c)\le l(\sigma+\sigma^2/2)$, hence
\[
|e_I(\ttd)-\nhat|\le\frac{\sigma(1+\sigma/2)}{1-\sigma^2/2}\le k_\sigma\sigma .
\]

\emph{Step 3 (rotation of $e_I$ on $[\ttd,t_j]$).} Write $a:=a_j$. From $c_Ie_I=x_b-x_a$ and $\partial_tx=vF(r)$, since $e_I'\perp e_I$,
$c_I|e_I'|=v|e_I\times(F(r_b)-F(r_a))|\le v\bigl(|e_I\times F(r_a)|+|e_I\times F(r_b)|\bigr)$, and
$|e\times F(r)|=f(\rho)|e\times\rhat|\le|e\times\rhat|$ (Lemma~\ref{lem:field}(i)). Let $t$ lie in leg $i+1$, $0\le i\le j-1$. By
Lemma~\ref{lem:chase}(1), $P=x_{a_{i+1}}-\varrho\uhat_{i+1}$ with $\varrho\ge0$, and $\uhat_{i+1}=e_{(a_i,a_{i+1})}(t)$ (the three points
are collinear in the order $x_{a_i},P,x_{a_{i+1}}$, and $c(a_i,a_{i+1};t)>0$ by (CA)). Put $W:=a_j-a_{i+1}\ge0$ and
$\kappa_W:=m\bigl((W+\lambda_I+\eps)/R\bigr)^{\gD}$.

\emph{($\alpha$) Close case, $W+\lambda_I+\eps\le R$.} Let $s:=C_3\kappa_W$. If $i\le j-2$, the pairs $K_1:=(a_{i+1},a_j)$,
$K_2:=(a_{i+1},b)$, $K_3:=(a_i,a_{i+1})$, $K_4:=(a_i,b)$ and $I$ belong to the family at $t$: their first points are $a_i$ or
$\ge a_{i+1}$, their label lengths are $\ge\eps/\kappa\ge\eps/2$ and $\le W+\lambda_I+\eps\le R$. Here $I\subset K_2\subset K_4$ share
the right end, $K_1\subset K_2$ and $K_3\subset K_4$ the left end, so by Theorem~\ref{thm:end}(c)
\[
|e_I-e_{K_1}|\le|e_I-e_{K_2}|+|e_{K_2}-e_{K_1}|\le2\CT^{\mathrm A}\kappa_W=s,\qquad
|e_I-\uhat_{i+1}|\le|e_I-e_{K_4}|+|e_{K_4}-e_{K_3}|\le s .
\]
Now $r_a=c(K_1)e_{K_1}+\varrho\uhat_{i+1}$ and $r_b=c(K_2)e_{K_2}+\varrho\uhat_{i+1}$ with non-negative coefficients; since
$|e\times e'|\le|e-e'|$ and $e\cdot e'=1-\frac12|e-e'|^2$, we get $|e_I\times r_a|\le s\,(c(K_1)+\varrho)$ and, if $s\le1$,
$\rho_a\ge r_a\cdot e_I\ge(1-s^2/2)(c(K_1)+\varrho)$; likewise for $b$ with $K_2$. If $i=j-1$, then $r_a=\varrho\uhat_j$,
$r_b=c_Ie_I+\varrho\uhat_j$, the pairs $(a_{j-1},a_j)$, $(a_{j-1},b)$ belong to the family (lengths $\le\eps+\lambda_I\le R$), and
$|e_I-\uhat_j|\le|e_I-e_{(a_{j-1},b)}|+|e_{(a_{j-1},b)}-e_{(a_{j-1},a_j)}|\le s$; the same bounds follow. In all cases
$|e_I\times\rhat_a|,|e_I\times\rhat_b|\le\psi_\star(s)$ (for $s>1$ trivially). As $c_I\ge c_A\lambda_I$ by (CA) ($\lambda_I\le R\le K_RR$),
\[
|e_I'|\le\frac{2v\,\psi_\star(C_3\kappa_W)}{c_A\lambda_I}.
\]

\emph{($\beta$) Far case, $W+\lambda_I+\eps>R$.} Then $a_j-a_i\ge W>R-\delta_0-\eps$ and $b\ge a_j$, so Theorem~\ref{thm:end}(d) (base
$a_i$; $K_R\ge1$) gives $\rho_a,\rho_b\ge c_A(R-\delta_0-\eps)-\delta_1\eps\ge\rho_*\ge\epsS\ge\eps$. Hence $F(r)=\rhat$ at $a$ and
$b$; with $p:=e_I\times r_a=e_I\times r_b$ (as $r_b-r_a=c_Ie_I$) and $|p|\le\min(\rho_a,\rho_b)$,
$|e_I\times(\rhat_b-\rhat_a)|=|p|\,|\rho_b-\rho_a|/(\rho_a\rho_b)\le c_I/\max(\rho_a,\rho_b)$. So $|e_I'|\le v/\rho_*$.

\emph{Durations.} Each non-final leg lasts at most $\eps/(1-v)$ (Lemma~\ref{lem:chase}(1)) and advances the base by at least
$\eps/\kappa$; so legs $1,\dots,j$ last at most $\kappa a_j/(1-v)\le\kappa\Lambda/(1-v)$ in total. The close-case legs have
$a_{i+1}\in[a_j-(R-\lambda_I-\eps),a_j]$, so there are at most $\kappa(R-\lambda_I-\eps)/\eps+1$ of them, and they last at most
$\kappa R/(1-v)$ ($\kappa\ge1$, $\lambda_I\ge\eps$). Thus the close-case time is $\le\kappa\min(\Lambda,R)/(1-v)$, and the far-case
time is $\le\kappa(\We+\eps)/(1-v)$. Moreover $W+\lambda_I+\eps\le a_j+\lambda_I+\eps\le2\Lambda$, so in the close case
$\kappa_W\le m\min(1,2\Lambda/R)^{\gD}$.

\emph{Rotation bound.} If $\Lambda\le\delta_0$ ($\le R/2$), then $\lambda_I=\Lambda$, $W+\lambda_I+\eps\le2\Lambda\le R$, only the
close case occurs, and with $\psi_\star(s)\le2s$,
\[
|e_I(t_j)-e_I(\ttd)|\le\frac{4vC_3m(2\Lambda/R)^{\gD}}{c_A\Lambda}\cdot\frac{\kappa\Lambda}{1-v}
=\frac{2^{\gD+2}v\kappa C_3m}{c_A(1-v)R^{\gD}}\,\Lambda^{\gD}.
\]
In general, with $\psi_\star\le1$ and $\min(\Lambda,R)/\lambda_I\le R/\delta_0$,
$|e_I(t_j)-e_I(\ttd)|\le\frac{v\kappa}{1-v}\bigl[\frac{2R}{c_A\delta_0}+\frac{\We+\eps}{\rho_*}\bigr]\le\vz K_{\mathrm{rot}}$.

\emph{Step 4 (sum).} By Steps 1--3, $|\uhat_{j+1}-\nhat|\le\CT^{\mathrm A}m(\lambda_I/R)^{\gD}+k_\sigma\sigma+|e_I(t_j)-e_I(\ttd)|$.
(i) Always, with $\lambda_I\le\delta_0$ and $\sigma\le\bar\sigma$, this is $\le\omega_0$. (ii) If $\Lambda\le\delta_0$, with
$\sigma\le\theta_G(2\Lambda)^{\gD}$ and $v\le\vz$ it is $\le\theta_{h1}\Lambda^{\gD}$. (iii) If $\Lambda>\delta_0$, it is
$\le\omega_0\le\omega_0\delta_0^{-\gD}\Lambda^{\gD}$. So $|\uhat_{j+1}-\nhat|\le\max(\theta_{h1},\omega_0\delta_0^{-\gD})\Lambda^{\gD}$.

\emph{Time.} Leg $i+1$ lasts at least $c(a_i,a_{i+1};t_i)$, since $\rho'\ge-1$ in Lemma~\ref{lem:chase}(1). By (J2) for the family pair
$(a_i,a_{i+1})$, whose length $\lambda$ is at most $\eps\le R$, and $l\ge\lambda$ (Lemma~\ref{lem:flow}(a)),
$c=l-E\ge\lambda-\frac12m^2R^{-2\gD}\lambda^{1+2\gD}\ge\lambda(1-m^2/2)$. Summing, $t_j-\ttd\ge(1-m^2/2)a_j$, so for
$\toff\ge t_j-\ttd$: $\Lambda\le c_\tau\toff+\eps\le c_\tau(\toff+\eps)$ ($c_\tau\ge1$). Hence
$|\uhat_{j+1}-\nhat|\le\vartheta_h(\toff+\eps)^{\gD}$, and $\le\omega_0$ by (i).
\end{proof}

\subsection{Departure}

\noindent\textbf{Setting.} $P$ is a $1$-Lipschitz path, and $x(s,t)$, $s\in[0,\UE]$, is a virtual flow of $P$ from a curve
$x^\circ$ with $|\partial_sx^\circ|\equiv1$ at time $\ttd$ (the \emph{departure arc}); its label $0$ is the end near $P$. $T$ is the
unit tangent in the direction of increasing $s$, $w:=|\partial_sx|$, and $l,c,N$ are as in Definition~\ref{def:excess} for pairs
$0\le a<b\le\UE$, $\lambda=b-a$. Let $\uhat$ be a unit vector, $0<D\le\Dbar$, $\toff:=t-\ttd\in[0,D]$, and for a label $s$:
\[
r_s=x_s-P=z_s\uhat+y_s\uhat^\perp,\qquad X_s(\toff):=s+z_0+\toff .
\]
Hypotheses:
\begin{itemize}
\item[(D0)] $x(0,\ttd)-P(\ttd)=z_0\uhat+y_0\uhat^\perp$ with $\eps/\zeta\le z_0\le z_0^{\max}$ and $|y_0|\le\theta_yz_0^{1+\gD}$;
\item[(D1)] $|T(s,\ttd)-\uhat|\le\theta_0(s+z_0)^{\gD}$ for $s\in[0,\UE]$;
\item[(D2)] for a.e.\ $\toff\in[0,D]$: $|P'|=1$, $P'\cdot\uhat<0$ and $|\uhat\times P'|\le\hH(\toff)$;
\item[(D3)] $\gD<1$, $\bar\omega\le\frac12$, $\zeta>0$, $\etaD\le\gD/2$, $c_k\ge1$, $\omega_0\le\bar\omega$, $\theta_0X_*^{\gD}\le\bar\omega$,
$\bar\omega_y:=\theta_yX_*^{\gD}\le\bar\omega$, $\omega_T:=\Theta_0X_*^{\gD}\le\frac12$.
\end{itemize}
The \emph{departure constants} are $\theta_{0y}:=\max(\theta_0,\theta_y)$, $x_c:=(\omega_0/\vartheta_h)^{1/\gD}$,
$L_h:=\frac1{2\gD}+\log^+\frac{\Dbar+\epsS}{x_c}$, and
\begin{equation}\label{eq:dep-const}
\begin{gathered}
\Theta_0:=\theta_0+\frac{\vz\theta_{0y}}{\zeta^2}+\frac{\vz\vartheta_h}{\zeta^2(1-\etaD)(\gD-\etaD)},\qquad
\Theta_b:=\Theta_0+\frac{\theta_{0y}}\zeta+\frac{\vartheta_h}{\zeta(1-\etaD)},\\
K_w:=\frac{(\Theta_0+\theta_{0y}/\zeta)^2X_*^{2\gD}}{\gD\zeta}+\frac{2\omega_0^2L_h}{\zeta^3(1-\etaD)^2},\qquad
\delta_t:=e^{\vz K_w}-1,\\
\Theta_{\mathrm p}:=\frac{(1+\bar\omega_y+\vz)(1+1/\kD)^{\gD}\Theta_0}{1-\omega_T^2/2}+\theta_{0y}+\frac{\vartheta_h}{1-\etaD},\\
c_n:=\frac{(1+\delta_t)^{1/2}\kD^{\gD}}{\gD\zeta},\qquad
c_\beta:=\frac{2(1+\delta_t)^{3/2}\kD^{\gD-1}}{\pi\zeta^3(1-\gD)} .
\end{gathered}
\end{equation}
They are numbers: the actual $v\le\vz$ and $\eps\le\epsS$ enter only through the hypotheses.

\begin{theorem}[Departure]\label{thm:departure}
Assume \textup{(D0)--(D3)}. For $\toff\in[0,D]$, $t=\ttd+\toff$, and every label $s\in[0,\UE]$:
\begin{enumerate}
\item[(i)] $z_s\ge\zeta X_s(\toff)$ and $\rho_s\ge\zeta X_s(\toff)\ge\eps$; hence $a^*(\rho_s)\le1/(\zeta X_s)$, and
\[
\exp\Bigl(v\int_{\toff'}^{\toff}a^*(\rho_s)\Bigr)\le\Bigl(\frac{X_s(\toff)}{X_s(\toff')}\Bigr)^{\etaD}\qquad(0\le\toff'\le\toff);
\]
\item[(ii)] $|y_s(\toff)|\le X_s(\toff)\bigl[\theta_{0y}X_s(0)^{\gD}+\hH(\toff)/(1-\etaD)\bigr]$;
\item[(iii)] $|T(s,t)-\uhat|\le\Theta_0X_s(\toff)^{\gD}\le\omega_T$;
\item[(iv)] $1\le w\le1+\delta_t$; hence $\lambda\le l_t(a,b)\le(1+\delta_t)\lambda$ for all pairs;
\item[(v)] for every pair with $\lambda\ge\eps/2$: $c(a,b;t)\ge(1-\omega_T^2/2)\,l_t(a,b)$ and
\[
N(a,b;t)\le\Bigl(1+\frac{\toff}{\kD\lambda}\Bigr)^{\etaN}\Bigl[N(a,b;\ttd)+v\,(c_n\Theta_b+c_\beta\Theta_{\mathrm p})\,\lambda^{1/2+\gD}\Bigr].
\]
\end{enumerate}
\end{theorem}

\begin{proof}
Write $\gamma:=\gD$, $X:=X_s$, $X_0:=X_s(0)$. Since $s\le\UE$, $z_0\le z_0^{\max}$ and $\toff\le\Dbar$, we have $X\le X_*$. The
quantities $\zeta,\etaD,\etaN$ are computed with $\vz$; we use $v\le\vz$ without further comment.

\textbf{(i)} By (D2), $-P'\cdot\uhat=(1-(\uhat\times P')^2)^{1/2}\ge(1-\bar\omega^2)^{1/2}$, as $\hH\le\omega_0\le\bar\omega$. Hence
$z_s'=vf(\rho_s)z_s/\rho_s-P'\cdot\uhat\ge-v+(1-\bar\omega^2)^{1/2}\ge\zeta$ a.e. At $\ttd$,
$z_s=z_0+\int_0^sT\cdot\uhat\,ds'$ with $T\cdot\uhat=1-\frac12|T-\uhat|^2\ge1-\bar\omega^2/2\ge\zeta$, because
$\theta_0(s'+z_0)^\gamma\le\theta_0X_*^\gamma\le\bar\omega$. So $z_s\ge z_0+\zeta(s+\toff)\ge\zeta X\ge\zeta z_0\ge\eps$, $\rho_s\ge z_s$,
$a^*(\rho_s)=1/\rho_s\le1/(\zeta X)$, and $\int_{\toff'}^{\toff}v/(\zeta X)\le\etaD\log(X(\toff)/X(\toff'))$.

\textbf{(ii)} As $\rho_s\ge\eps$, $F(r_s)=r_s/\rho_s$ and $y_s'=(v/\rho_s)y_s-\uhat\times P'$; so $|y_s|'\le va^*(\rho_s)|y_s|+\hH$ a.e.,
and by (i)
$|y_s(\toff)|\le(X(\toff)/X_0)^{\etaD}|y_s(0)|+\int_0^{\toff}(X(\toff)/X(\toff'))^{\etaD}\hH(\toff')\,d\toff'$. At $\ttd$,
$y_s(0)=y_0+\int_0^s\uhat\times(T-\uhat)\,ds'$, so by (D0), (D1)
$|y_s(0)|\le\theta_0sX_0^\gamma+\theta_yz_0^{1+\gamma}\le\theta_{0y}(s+z_0)X_0^\gamma=\theta_{0y}X_0^{1+\gamma}$. In the integral, $X(\toff')\ge\toff'+z_0$ and $\hH$ is
non-decreasing, so it is at most $X(\toff)^{\etaD}\hH(\toff)(\toff+z_0)^{1-\etaD}/(1-\etaD)\le X(\toff)\hH(\toff)/(1-\etaD)$; and
$(X/X_0)^{\etaD}X_0^{1+\gamma}\le XX_0^\gamma$. This gives (ii) and, for $0\le\toff'\le\toff$,
\begin{equation}\label{eq:dep-y}
|y_s(\toff')|\le\Bigl(\frac{X(\toff')}{X_0}\Bigr)^{\etaD}\theta_{0y}X_0^{1+\gamma}
+\frac{X(\toff')^{\etaD}(\toff'+z_0)^{1-\etaD}\hH(\toff')}{1-\etaD}.
\end{equation}

\textbf{(iii)} By Lemma~\ref{lem:flow}(b), $T_t=v(AT-(T\cdot AT)T)$, so $\mathfrak d:=1-T\cdot\uhat=\frac12\delta^2$, $\delta:=|T-\uhat|$,
satisfies $\partial_t\mathfrak d=v[(T-\uhat)\cdot AT-(T\cdot AT)\mathfrak d]\le v(T-\uhat)\cdot AT$. Write $T-\uhat=:\mathbf d$; then
$\mathbf d\cdot AT=\mathbf d\cdot A\mathbf d+\mathbf d\cdot A\uhat$, $\mathbf d\cdot A\mathbf d\le a^*\delta^2$, and $\mathbf d\cdot A\uhat=(\uhat\cdot A\uhat)(\mathbf d\cdot\uhat)+(\uhat^\perp\cdot A\uhat)(\mathbf d\cdot\uhat^\perp)
\le\delta|q_{\uhat}|$, where $q_{\uhat}:=\uhat^\perp\cdot A\uhat$, since $\mathbf d\cdot\uhat=-\mathfrak d\le0$ and $A\ge0$. By
Lemma~\ref{lem:field}(iii), $|q_{\uhat}|\le|r_s\times\uhat|/\rho_s^2=|y_s|/\rho_s^2\le|y_s|/(\zeta X)^2$. With
$\delta_\varpi:=(\delta^2+\varpi^2)^{1/2}$ we get $\delta_\varpi'\le va^*\delta_\varpi+v|q_{\uhat}|$; by variation of constants, (i), and
$\varpi\to0$,
\[
\delta(\toff)\le\Bigl(\frac X{X_0}\Bigr)^{\etaD}\delta(0)+\frac v{\zeta^2}\int_0^{\toff}\Bigl(\frac{X(\toff)}{X(\toff')}\Bigr)^{\etaD}
\frac{|y_s(\toff')|}{X(\toff')^2}\,d\toff'.
\]
Insert \eqref{eq:dep-y}. The first part contributes at most
$\frac v{\zeta^2}\theta_{0y}X_0^{1+\gamma}(X/X_0)^{\etaD}\int_0^\infty(X_0+\toff')^{-2}d\toff'\le\frac{v\theta_{0y}}{\zeta^2}X^\gamma$. The
second, using $X(\toff')\ge\toff'+z_0$, $\hH(\toff')\le\vartheta_h(\toff'+\eps)^\gamma\le\vartheta_h(\toff'+z_0)^\gamma$ and $\gamma>\etaD$,
at most
\[
\frac{v\vartheta_hX^{\etaD}}{\zeta^2(1-\etaD)}\int_0^{\toff}(\toff'+z_0)^{\gamma-1-\etaD}d\toff'
\le\frac{v\vartheta_h}{\zeta^2(1-\etaD)(\gamma-\etaD)}X^{\etaD}(\toff+z_0)^{\gamma-\etaD}\le\frac{v\vartheta_hX^\gamma}{\zeta^2(1-\etaD)(\gamma-\etaD)}.
\]
With $\delta(0)\le\theta_0X_0^\gamma$ by (D1) and $(X/X_0)^{\etaD}X_0^\gamma\le X^\gamma$: $\delta\le\Theta_0X^\gamma\le\Theta_0X_*^\gamma=\omega_T$.

\textbf{(iv)} By Lemma~\ref{lem:flow}(b), $w_t/w=v\sin^2\!\phi/\rho$ ($\rho\ge\eps$), and
$|\sin\phi|=|T\times\rhat|\le\delta+|\uhat\times\rhat|=\delta+|y|/\rho\le(\Theta_0+\theta_{0y}/\zeta)X^\gamma+\hH/(\zeta(1-\etaD))$ by (ii),
(iii) and $\rho\ge\zeta X$. With $(\alpha+\beta)^2\le2\alpha^2+2\beta^2$, $\rho\ge\zeta X$ and $X(\toff')\ge\toff'+\eps$,
\[
\int_0^{\toff}\frac{\sin^2\phi}\rho\,d\toff'\le\frac2\zeta\Bigl(\Theta_0+\frac{\theta_{0y}}\zeta\Bigr)^2\int_0^{\toff}X^{2\gamma-1}d\toff'
+\frac2{\zeta^3(1-\etaD)^2}\int_0^D\frac{\hH(\toff')^2}{\toff'+\eps}\,d\toff'\le K_w,
\]
because $\int_0^{\toff}X^{2\gamma-1}\le X_*^{2\gamma}/(2\gamma)$ and
$\int_0^D\min(\vartheta_h^2(\toff'+\eps)^{2\gamma},\omega_0^2)\frac{d\toff'}{\toff'+\eps}\le\int_0^{x_c}\vartheta_h^2y^{2\gamma-1}dy
+\int_{x_c}^{\Dbar+\epsS}\frac{\omega_0^2}y\,dy\le\omega_0^2L_h$ (the second integral is absent if $\Dbar+\epsS\le x_c$). So
$w\le e^{vK_w}\le1+\delta_t$; $w\ge1$ by Lemma~\ref{lem:flow}(a).

\textbf{(v)} \emph{Chord:} $c\ge(x_b-x_a)\cdot\uhat=\int_a^bT\cdot\uhat\,ds\ge(1-\omega_T^2/2)\,l$ by (iii). Put $\tilde g:=a+z_0$
and $\toff_0:=\max(0,\kD\lambda-\tilde g)$; for labels $s\in[a,b]$, $X_s\ge\tilde g+\toff\ge\toff+\eps$.

\emph{Phase $\toff\le\toff_0$.} By Lemma~\ref{lem:excess}(iv), $E'\le vN\|AT\|_{L^2(ds)}$; with $N_\varpi:=(N^2+\varpi^2)^{1/2}$,
$N_\varpi'\le v\|AT\|_{L^2(ds)}$, and we let $\varpi\to0$. On the arc $\rho\ge\eps$, so $A=\Pi/\rho$ and, as in (iv) with
$\hH\le\vartheta_h(\toff+\eps)^\gamma\le\vartheta_hX^\gamma$,
$|AT|=|\sin\phi|/\rho\le\Theta_bX_s^{\gamma-1}/\zeta\le\Theta_b(\tilde g+\toff)^{\gamma-1}/\zeta$. By (iv),
$\|AT\|_{L^2(ds)}\le(1+\delta_t)^{1/2}\lambda^{1/2}\Theta_b(\tilde g+\toff)^{\gamma-1}/\zeta$, hence for $\toff\le\toff_0$
\[
N(\toff)\le N(0)+v(1+\delta_t)^{1/2}\frac{\Theta_b}\zeta\lambda^{1/2}\frac{(\tilde g+\toff_0)^\gamma}\gamma\le N(0)+vc_n\Theta_b\lambda^{1/2+\gamma},
\]
since $\tilde g+\toff_0=\kD\lambda$ if $\toff_0>0$ (if $\toff_0=0$ there is nothing to prove).

\emph{Phase $\toff\ge\toff_0$.} The arc $[a,b]$ lies in $\{x:(x-P)\cdot\uhat\ge\min_sz_s\}$, so by (i), (iv) and
$\kD\ge4(1+\delta_t)/\zeta$ its arc distance satisfies
$\rhoarc\ge\zeta(\tilde g+\toff)\ge\zeta\kD\lambda\ge4(1+\delta_t)\lambda\ge4l$, and $\rhoarc\ge\zeta z_0\ge\eps$. Lemma~\ref{lem:duhamel}
with $K=4$ applies on $[\toff_0,\toff]$, with $\alpha\le\aD(4)v/\rhoarc\le\etaN/(\tilde g+\toff')$ and
$\beta=2vl^{3/2}|p_a|/(\pi\rhoarc^3)$.
Here $p_a=r_a\times e_{ab}=z_a(\uhat\times e_{ab})+y_a(\uhat^\perp\times e_{ab})$, so $|p_a|\le\rho_a|\uhat\times e_{ab}|+|y_a|$, and:
$|\uhat\times e_{ab}|\le(l/c)\sup_{[a,b]}\delta\le\Theta_0X_b^\gamma/(1-\omega_T^2/2)$ with
$X_b=\tilde g+\lambda+\toff'\le(1+1/\kD)(\tilde g+\toff')$; $\rho_a(\toff')\le a+z_0+|y_0|+(1+v)\toff'\le(1+\bar\omega_y+v)(\tilde g+\toff')$,
as $|y_0|\le\theta_yz_0^{1+\gamma}\le\bar\omega_yz_0$; and $|y_a|\le(\tilde g+\toff')^{1+\gamma}(\theta_{0y}+\vartheta_h/(1-\etaD))$ by (ii)
($X_a=\tilde g+\toff'$, $\hH\le\vartheta_h(\tilde g+\toff')^\gamma$). Hence $|p_a|\le\Theta_{\mathrm p}(\tilde g+\toff')^{1+\gamma}$ and
$\beta\le2v(1+\delta_t)^{3/2}\Theta_{\mathrm p}\lambda^{3/2}(\tilde g+\toff')^{\gamma-2}/(\pi\zeta^3)$. As
$\int_{\toff_0}^\infty(\tilde g+\toff')^{\gamma-2}d\toff'\le(\kD\lambda)^{\gamma-1}/(1-\gamma)$, Lemma~\ref{lem:duhamel} gives
\[
N(\toff)\le\Bigl(\frac{\tilde g+\toff}{\tilde g+\toff_0}\Bigr)^{\etaN}\bigl[N(\toff_0)+vc_\beta\Theta_{\mathrm p}\lambda^{1/2+\gamma}\bigr],
\]
and $(\tilde g+\toff)/(\tilde g+\toff_0)\le1+\toff/(\kD\lambda)$ (if $\toff_0=0$ because $\tilde g\ge\kD\lambda$; otherwise
$\tilde g+\toff_0=\kD\lambda$). Combined with the first phase this is (v).
\end{proof}

% sec6.tex -- writer W4. Straight runs and design T.
\section{Straight runs and design T}\label{sec:designT}

Design T consists of two straight runs. On a straight run the true candidates obey two exact integrals, and so do the first and
second derivatives of the front in the label; this section collects these exact facts (valid for every $v\in(0,1)$) and defines
design T.

\emph{Heading frame.} While $P$ moves with constant unit velocity (its \emph{heading}), we use coordinates in which $P$ is at the
origin and the heading is $1\in\C$; a candidate is at $r=\rho e^{i\beta}=z+iy$, and for $0<\beta<\pi$ we put $\tau:=\tan(\beta/2)$.

\begin{proposition}\label{prop:runs}
On a straight run, a true candidate with $\rho>0$ satisfies
\[
\rho'=v-\cos\beta,\qquad\rho\,\beta'=\sin\beta;\qquad h:=\rho(\cos\beta+v)\ \text{ satisfies }\ h'=-(1-v^2).
\]
If $0<\beta<\pi$, then $\beta$ stays in $(0,\pi)$ and $C:=\rho\,\sg(\beta)$ is constant, where
\[
\sg(\beta):=\frac{\sin\beta}{\tan^v(\beta/2)}=\frac{2\tau^{1-v}}{1+\tau^2},\qquad\frac{\sg'}{\sg}=\frac{\cos\beta-v}{\sin\beta}.
\]
The function $\sk:=(\cos\beta+v)/\sg$ is strictly decreasing on $(0,\pi)$, from $+\infty$ to $-\infty$; $\sg$ is increasing on
$(0,\arccos v]$ and decreasing on $[\arccos v,\pi)$, and
\[
1=\sg(\pi/2)\le\gmax:=\sg(\arccos v)=(1+v)^{(1+v)/2}(1-v)^{(1-v)/2}\le e^{v^2}.
\]
\end{proposition}
\begin{proof}
$r'=v\rhat-1$; the components along $\rhat$ and $\rhat^\perp$ give the first two equations. Then
$h'=(\cos\beta+v)(v-\cos\beta)-\rho\sin\beta\,\beta'=v^2-\cos^2\beta-\sin^2\beta$. Since $\beta'=0$ at $\beta\in\{0,\pi\}$, $\beta$
cannot leave $(0,\pi)$. On $(0,\pi)$, $\log\sg=\log\sin\beta-v\log\tan(\beta/2)$ gives $\sg'/\sg=\cot\beta-v/\sin\beta$, and, as $\beta'/\sin\beta=1/\rho$,
$\rho'/\rho=(v-\cos\beta)\beta'/\sin\beta=-(\sg'/\sg)(\beta)\,\beta'$; so $C$ is constant. With $\sin\beta=2\tau/(1+\tau^2)$ we get the $\tau$-form of $\sg$, so $\sg\to0$ at
both ends of $(0,\pi)$. Next,
$\sk'=[-\sin\beta-(\cos\beta+v)(\cos\beta-v)/\sin\beta]/\sg=-(1-v^2)/(\sg\sin\beta)<0$, and $\cos\beta+v\to1+v$, resp.\ $-(1-v)$,
at $0$, resp.\ $\pi$, which gives the limits of $\sk$. The monotonicity of $\sg$ follows from the sign of $\cos\beta-v$. At
$\cos\beta=v$: $\sin\beta=(1-v^2)^{1/2}$ and $\tan(\beta/2)=((1-v)/(1+v))^{1/2}$, which gives the closed form. Finally
$\log\gmax=\frac12[(1+v)\log(1+v)+(1-v)\log(1-v)]=\sum_{n\ge1}\frac{v^{2n}}{2n(2n-1)}\le v^2\sum_{n\ge1}\frac1{2n(2n-1)}=v^2\log2\le v^2$.
\end{proof}

For a $1$-Lipschitz path $P$ and a time $t$, the \emph{true front} is the curve $\theta\mapsto S_\theta(t)$ on the labels whose true
candidates exist on $[0,t]$. The field $x\mapsto(x-P(t))/|x-P(t)|$ is $C^\infty$ in $x\ne P(t)$ and continuous in $t$; hence near
every label whose candidate has had $\rho>0$ on $[0,t]$ the front is $C^\infty$ in $\theta$, and $q:=\partial_\theta S_\theta$
satisfies $\partial_tq=v(\Pi/\rho)q$, $\Pi=I-\rhat\rhat^{\mathsf T}$. We write $w=|q|$, $T=q/w=e^{i\psi}$, $\phi=\arg r-\psi$ and
$k=\partial_s\psi=\partial_\theta\psi/w$ (so $k=1$ for the unit circle), as for virtual fronts (Definition~\ref{def:flow}).

\begin{lemma}\label{lem:truefront}
At every label and time with $\rho>0$ on $[0,t]$, the true front satisfies
\[
\frac{w_t}w=\frac{v\sin^2\!\phi}\rho,\qquad\psi_t=-\frac v\rho\sin\phi\cos\phi,\qquad
k_t=\frac v\rho(1-3\sin^2\!\phi)\,k+\frac v{\rho^2}\sin\phi\,(3\cos^2\!\phi-1),
\]
and $w\ge1$. If the true or the virtual candidate $\theta$ has $\rho\ge\eps$ on $[0,t]$, then $x_\theta=S_\theta$ on $[0,t]$; if this
holds for all labels of an open set, the true and the virtual front coincide there with all their label derivatives.
\end{lemma}
\begin{proof}
Let $n:=\rhat^\perp$; then $T=\cos\phi\,\rhat-\sin\phi\,n$, $\Pi T=-\sin\phi\,n$, $T^\perp=\cos\phi\,n+\sin\phi\,\rhat$. From
$\partial_t(wT)=v\Pi(wT)/\rho$: $w_t=vw\,T\cdot\Pi T/\rho=vw\sin^2\!\phi/\rho$, and
\[
T_t=\frac v\rho(\Pi T-(T\cdot\Pi T)T)=-\frac v\rho[\sin\phi\cos^2\!\phi\,n+\sin^2\!\phi\cos\phi\,\rhat]=-\frac v\rho\sin\phi\cos\phi\,T^\perp .
\]
At $t=0$, $w=|\partial_\theta e^{i\theta}|=1$, and $w$ is non-decreasing. For the curvature: $\partial_s\rho=\cos\phi$,
$\partial_s\rhat=\Pi T/\rho=-(\sin\phi/\rho)n$, so $\partial_s\phi=-\sin\phi/\rho-k$. Since $k=\partial_\theta\psi/w$ and the front is
smooth, $k_t=\partial_s(\psi_t)-k\,w_t/w$, and
\[
\partial_s\Bigl(-\frac v{2\rho}\sin2\phi\Bigr)=\frac v{\rho^2}\sin\phi\cos^2\!\phi+\frac v\rho\cos2\phi\Bigl(k+\frac{\sin\phi}\rho\Bigr);
\]
so $k_t=\frac v\rho k(\cos2\phi-\sin^2\!\phi)+\frac v{\rho^2}\sin\phi(\cos^2\!\phi+\cos2\phi)$, which is the claim. Finally, $F(x)=x/|x|$
for $|x|\ge\eps$ (Lemma~\ref{lem:field}(i)): a true candidate with $\rho\ge\eps$ on $[0,t]$ solves the virtual equation, whose
solutions are unique (Lemma~\ref{lem:flow}(a)); a virtual candidate with $\rho\ge\eps$ solves the true equation, whose solutions are
unique while $\rho>0$.
\end{proof}

From now on in this section, $P$ moves with constant unit velocity on a time interval, coordinates are those of the heading frame,
and candidates have $0<\beta<\pi$.

\begin{lemma}\label{lem:coord}
\begin{enumerate}
\item[(a)] $\rho=\frac C2(\tau^{v-1}+\tau^{v+1})$, $y=\rho\sin\beta=C\tau^v$, $z=\rho\cos\beta=\frac C2(\tau^{v-1}-\tau^{v+1})$, and
$h=\frac C2[(1+v)\tau^{v-1}-(1-v)\tau^{1+v}]$.
\item[(b)] $\tau\mapsto(1+v)\tau^{v-1}-(1-v)\tau^{1+v}$ is a strictly decreasing bijection $(0,\infty)\to\R$. Hence $\tau$ is strictly
increasing in time, and the state at any time of the run is determined by $(C,h)$.
\item[(c)] $\min_\tau\rho=C/\gmax$, attained at $\tau^2=(1-v)/(1+v)$, i.e.\ $\cos\beta=v$; $\rho$ is decreasing in time while
$\cos\beta>v$.
\item[(d)] If $h>0$: $\tau\le((1+v)C/(2h))^{1/(1-v)}$. If $h<0$: $\tau\ge(2|h|/((1-v)C))^{1/(1+v)}$, and
\[
\tau^{1+v}=\frac{2|h|}{(1-v)C}\Bigl[1+\frac{(1+v)C\tau^{v-1}}{2|h|}\Bigr].
\]
\end{enumerate}
\end{lemma}
\begin{proof}
(a) $\sin\beta=2\tau/(1+\tau^2)$, $\cos\beta=(1-\tau^2)/(1+\tau^2)$, $\rho=C/\sg(\beta)$ and $h=z+v\rho$. (b) The derivative is
$-(1-v^2)(\tau^{v-2}+\tau^v)<0$, the limits are $\pm\infty$, and $h'<0$. (c) $\rho'=v-\cos\beta$ (Proposition~\ref{prop:runs}), and
$\frac d{d\tau}(\tau^{v-1}+\tau^{v+1})=\tau^{v-2}((1+v)\tau^2-(1-v))$. (d) By (a), $(1+v)\tau^{v-1}=2h/C+(1-v)\tau^{1+v}\ge2h/C$, resp.\
$(1-v)\tau^{1+v}=2|h|/C+(1+v)\tau^{v-1}\ge2|h|/C$.
\end{proof}

\begin{lemma}[transport]\label{lem:transport}
Let $\theta\mapsto S_\theta$ be a $C^2$ family of true candidates with $0<\beta<\pi$ on the run, $q:=\partial_\theta S_\theta$
(heading frame), $w=|q|$, $q=we^{i\psi}$, $\phi=\beta-\psi$, and
\[
\Xi:=q+v\,e^{i\beta}\,\bar q .
\]
\begin{enumerate}
\item[(a)] $\operatorname{Re}\Xi=\partial_\theta h=:\sa$ and $\operatorname{Im}\Xi=\tau^v\partial_\theta C=:\tau^v\sfc$, where
$\sa(\theta)$, $\sfc(\theta)$ do not depend on time; $q=(\Xi-ve^{i\beta}\bar\Xi)/(1-v^2)$ and $(1-v)w\le|\Xi|\le(1+v)w$.
\item[(b)] Let $Y_r:=\sfc\tau^v/\sa$ (where $\sa\ne0$), $N_\psi:=(1+v\cos\beta)Y_r-v\sin\beta$, $D_\psi:=1-v\cos\beta-v\sin\beta\,Y_r$.
Then $(1-v^2)q=\sa(D_\psi+iN_\psi)$. If $\sa>0$ and $D_\psi>0$: $|\psi|<\pi/2$, $\tan\psi=N_\psi/D_\psi$, $|\psi|\le|\tan\psi|$ and
$w=\sa|D_\psi+iN_\psi|/(1-v^2)$.
\item[(c)] $J:=\sa\,\partial_\theta\sfc-\sfc\,\partial_\theta\sa$ does not depend on time, and
\[
k=\frac{\tau^vJ}{(1-v^2)\,w^3}-\TS,\qquad
\TS:=\frac{v\sin\phi}{(1-v^2)\,\rho}\Bigl[\frac{\rho\,\sa\,\partial_\theta\log C}{w^2}-v-\cos(\phi-\psi)\Bigr].
\]
\item[(d)] For two times $t'<t''$ of the run (data $'$, $''$): $\tau''\ge\tau'$,
\[
k''=\Bigl(\frac{\tau''}{\tau'}\Bigr)^v\Bigl(\frac{w'}{w''}\Bigr)^3(k'+\TS')-\TS'',\qquad\text{so}\qquad
|k''|\le\Bigl(\frac{\tau''}{\tau'}\Bigr)^v(|k'|+|\TS'|)+|\TS''|,
\]
and $w''\le(\tau''/\tau')^v|\Xi'|/(1-v)\le\frac{1+v}{1-v}(\tau''/\tau')^vw'$.
\item[(e)] At every time, $|\TS|\le\dfrac{v|\sin\phi|}{(1-v^2)\rho}\Bigl[\dfrac{\rho|\sa||\partial_\theta\log C|}{w^2}+1+v\Bigr]$
and $|\TS|\le\dfrac{v(3+v)}{2(1-v)y}$.
\end{enumerate}
\end{lemma}
\begin{proof}
(a) Write $q=X+iY$. As $P$ does not depend on $\theta$: $\partial_\theta z=X$, $\partial_\theta y=Y$,
$\partial_\theta\rho=X\cos\beta+Y\sin\beta$, $\rho\,\partial_\theta\beta=-X\sin\beta+Y\cos\beta=-w\sin\phi$. Since
$ve^{i\beta}\bar q=v(X\cos\beta+Y\sin\beta)+iv(X\sin\beta-Y\cos\beta)$,
\[
\partial_\theta h=X+v\,\partial_\theta\rho=\operatorname{Re}\Xi .
\]
By Lemma~\ref{lem:coord}(a), $C=y\tau^{-v}$, and $\partial_\theta\log\tau=\partial_\theta\beta/\sin\beta$; so
$\tau^v\partial_\theta C=Y-v\rho\,\partial_\theta\beta=Y(1-v\cos\beta)+vX\sin\beta=\operatorname{Im}\Xi$. Since $C$ is constant and
$h+(1-v^2)t$ is constant (Proposition~\ref{prop:runs}), $\sa$ and $\sfc$ do not depend on time. The map
$L:x\mapsto x+vR_\beta x$, $R_\beta x:=e^{i\beta}\bar x$, is $I+vR_\beta$ with $R_\beta$ a reflection; so $L^{-1}=(I-vR_\beta)/(1-v^2)$
and the singular values of $L$ are $1\pm v$.

(b) With $\Xi=\sa(1+iY_r)$: $(1-v^2)q=\Xi-ve^{i\beta}\bar\Xi$ and
$e^{i\beta}(1-iY_r)=\cos\beta+Y_r\sin\beta+i(\sin\beta-Y_r\cos\beta)$ give $(1-v^2)q=\sa(D_\psi+iN_\psi)$. If $\sa,D_\psi>0$, then
$\operatorname{Re}q>0$, so $|\psi|<\pi/2$ and $|\psi|\le|\tan\psi|$.

(c) $J$ is built from the time-independent $\sa,\sfc$. Put $V:=\Xi=\sa+i\sfc\tau^v=Lq$. Then
$\partial_\theta V=\partial_\theta\sa+i(\tau^v\partial_\theta\sfc+v\sfc\tau^v\partial_\theta\beta/\sin\beta)$,
$\partial_\beta(L^{-1})=-vR'_\beta/(1-v^2)$ with $R'_\beta x:=ie^{i\beta}\bar x$, and
$\partial_\theta q=L^{-1}\partial_\theta V+\partial_\theta\beta\,\partial_\beta(L^{-1})V$. For a real $2\times2$ matrix $M$,
$(Ma)\times(Mb)=\det M\,(a\times b)$, and $\det L^{-1}=1/(1-v^2)$; hence
\[
kw^3=q\times\partial_\theta q=\frac{V\times\partial_\theta V}{1-v^2}-\frac{v\,\partial_\theta\beta}{1-v^2}\,q\times R'_\beta(Lq).
\]
Here $V\times\partial_\theta V=\tau^vJ+v\sa\sfc\tau^v\partial_\theta\beta/\sin\beta$ and $\sfc\tau^v/\sin\beta=\rho\,\partial_\theta\log C$
(Lemma~\ref{lem:coord}(a): $\tau^v/\sin\beta=\tau^{v-1}(1+\tau^2)/2=\rho/C$). Further $Lq=w(e^{i\psi}+ve^{i\phi})$,
$R'_\beta(Lq)=iw(e^{i\phi}+ve^{i\psi})$, $q\times R'_\beta(Lq)=w^2(\cos(\phi-\psi)+v)$, and $\partial_\theta\beta=-w\sin\phi/\rho$.
Inserting gives
$kw^3=\frac{\tau^vJ}{1-v^2}-\frac{vw\sin\phi}{(1-v^2)\rho}[\rho\sa\,\partial_\theta\log C-w^2(v+\cos(\phi-\psi))]$.

(d) Eliminate $J$ between the formulas of (c) at $t'$ and $t''$. $\tau$ increases in time (Lemma~\ref{lem:coord}(b)) and $w$ is
non-decreasing (Lemma~\ref{lem:truefront}), so $w'/w''\le1$. $\Xi''=\sa+i\tau''^v\sfc$ gives $|\Xi''|\le(\tau''/\tau')^v|\Xi'|$; then use
(a).

(e) The first bound is (c) with $|v+\cos(\phi-\psi)|\le1+v$. For the second: $\partial_\theta\log C=\sfc/C=\operatorname{Im}\Xi/y$ and
$\sa=\operatorname{Re}\Xi$, so $\rho|\sa\,\partial_\theta\log C|/w^2\le\rho|\Xi|^2/(2yw^2)\le(1+v)^2\rho/(2y)$; with $|\sin\phi|\le1$ and
$\rho\ge y$: $|\TS|\le\frac v{1-v}[\frac{1+v}{2y}+\frac1\rho]\le\frac{v(3+v)}{2(1-v)y}$.
\end{proof}

\begin{definition}[design T]\label{def:T}
Let $t_1:=1/(1-v)+1$ $(=2+\mu)$.
\begin{itemize}
\item \emph{Run 1.} $P(t):=t$ for $0\le t\le t_1$.
\item \emph{Run 2.} Let $\alpha_0:=\arccos(v+1/t_1)$ and write $S_{\alpha_0}(t_1)-P(t_1)=\rho_0e^{i\beta^*}$ ($\rho_0>0$), $\nhat:=e^{i\beta^*}$.
Then $P(t):=P(t_1)+(t-t_1)\nhat$ for $t_1\le t\le\ttd$, where
\[
\ttd:=t_1+\frac{\rho_0-\eps}{1-v}+\eps\int_0^1\frac{d\sigma}{1-v\sigma(2-\sigma)} .
\]
\end{itemize}
Further $\ta:=t_1+(\rho_0-\eps)/(1-v)$, $\tc:=t_1+\rho_0/(1-v)$, $X_0:=P(\ttd)$ and $\omega(t):=\tc-t$.
\end{definition}

Here $v+1/t_1=v+(1-v)/(2-v)\in(\frac12,1)$, and $\rho_0>0$ because in run 1 the candidate $\alpha_0\in(0,\pi/2)$ has
$C=\sg(\alpha_0)>0$, so $\rho\ge C/\gmax$ (Lemma~\ref{lem:coord}(c)); $\rho_0>\eps$ by Lemma~\ref{lem:tang}(b), so $t_1<\ta<\ttd<\tc$.
$P$ is $1$-Lipschitz, $P(0)=0$, and $P$ depends only on $(v,\eps)$: $\nhat$ is computed from the state of the true candidate $\alpha_0$
at $t_1$, which is determined by run 1.

\emph{Run-1 data.} For $\alpha\in(0,\pi)$ let $(\rho_1,\beta_1)=(\rho_1(\alpha),\beta_1(\alpha))$ be the polar coordinates of
$S_\alpha(t_1)-P(t_1)$ in the heading frame of run 1, $\tau_1:=\tan(\beta_1/2)$, $\tau_\alpha:=\tan(\alpha/2)$ and
$y(\alpha,t_1):=\operatorname{Im}S_\alpha(t_1)$. In run 1 the lower half is the mirror image of the upper one,
$S_{2\pi-\alpha}=\overline{S_\alpha}$. At $t=0$, $\rho=1$, $\beta=\alpha$; hence, in run 1,
\begin{gather*}
C=\sg(\alpha),\qquad h(t)=\cos\alpha+v-(1-v^2)t,\qquad h_1:=h(t_1)=-(1-v^2)-(1-\cos\alpha)<0,\\
\sa=-\sin\alpha,\qquad\sfc=(\cos\alpha-v)\tau_\alpha^{-v},
\end{gather*}
using $(1-v^2)t_1=(1+v)(2-v)$ and $\sg'=\sg(\cos\alpha-v)/\sin\alpha$.

\emph{Run-2 data.} For $\theta\in(0,2\pi)\setminus\{\alpha_0\}$ let $(\sr(\theta),\sbb(\theta))$ be the polar coordinates of
$S_\theta(t_1)-P(t_1)$ about $\nhat$ (the heading frame of run 2), $\tau_{\sbb}:=\tan(\sbb/2)$, and, for run 2,
\[
C:=\sr\,\sg(\sbb),\qquad\sa:=\partial_\theta[\sr(\cos\sbb+v)],\qquad\sfc:=\partial_\theta C,\qquad\sA:=\partial_\theta\sr .
\]
From now on $C$, $h$, $\sa$, $\sfc$ refer to run 2 unless said otherwise. The \emph{window} is $I_\eta:=\{\theta:|\theta-\alpha_0|\le\eta\}$,
$\eta=0.11639$ (Section~\ref{sec:touch}), and $\xi:=|\theta-\alpha_0|$. In the frame of run 2 at time $t$ we write $S_\theta(t)-P(t)=(z+iy)\nhat=\rho e^{i\beta}\nhat$
and $e^{i\psi}\nhat$ for the unit tangent. Since $(1-v^2)(\tc-t_1)=(1+v)\rho_0$, every label whose true candidate has $\rho>0$ on run 2 (including $\alpha_0$, for $t<\tc$) has in run 2
\begin{equation}\label{eq:h2}
h(t)=\sr(\cos\sbb+v)-(1-v^2)(t-t_1)=(1+v)(\sr-\rho_0)-\sr(1-\cos\sbb)+(1-v^2)\,\omega(t).
\end{equation}
Whenever $h(t)>0$, Lemma~\ref{lem:coord}(d) gives $\tau\le((1+v)C/(2h))^{1+\mu}$, hence $y=C\tau^v\le C^{1+\mu}((1+v)/(2h))^\mu$, and
$C\le2\sr\tau_{\sbb}^{1-v}$ gives $C^{1+\mu}\le(2\sr)^{1+\mu}\tau_{\sbb}$; therefore
\begin{equation}\label{eq:tau2}
\tau\le\tau_{\sbb}\Bigl(\frac{(1+v)\sr}h\Bigr)^{1+\mu},\qquad y\le2\sr\,\tau_{\sbb}\Bigl(\frac{(1+v)\sr}h\Bigr)^{\mu}\qquad(h>0).
\end{equation}

% sec7.tex -- only inputs the two halves (W4: sec7a, W5: sec7b).
% sec7a.tex -- writer W4. Section 7 heading, 7.1 (constants), 7.2 (run 1, tangency, forward arc, virtual labels).
\section{The front at touch-down}\label{sec:touch}

In this section $P$ is design T (Definition~\ref{def:T}) and nothing after $\ttd$ is considered. The fixed data are
\[
\vbar=5.8\cdot10^{-3},\quad\eta=0.11639,\quad\epsbar=10^{-3},\quad\xi_1=1.16171,\quad\lambda_0=0.0337254,\quad G_{\mathrm T}=0.030573,
\]
and throughout $0<v\le\vbar$, $0<\eps\le\epsbar$. We put $\mub:=\vbar/(1-\vbar)$, $\muT:=3v/(1+v)$, $\muTb:=3\vbar/(1+\vbar)$; so
$\mu\le\mub$ and $\muT\le\muTb$. The notation of Section~\ref{sec:designT} is used throughout ($t_1$, $\alpha_0$, $\rho_0$, $\beta^*$,
$\nhat$, $\ta$, $\ttd$, $\tc$, $X_0$, $\omega$; run-1 data $\rho_1,\beta_1,\tau_1,\tau_\alpha,y(\alpha,t_1),h_1$; run-2 data
$\sr,\sbb,\tau_{\sbb},C,\sa,\sfc,\sA$; $I_\eta$, $\xi$).

\subsection{The constants of the touch-down}\label{sec:tdconst}

In the next definition every formula is evaluated with $v:=\vbar$, $\mu:=\mub$, $\muT:=\muTb$, $\eps:=\epsbar$; Lemma~\ref{lem:mono}
explains why this is legitimate. The numerical values are listed in Table~\ref{tab:td}.

\begin{definition}[touch-down constants]\label{def:tdconst}
{\small
\begin{itemize}\setlength\itemsep{2pt}
\item \emph{Basic.} $\bar t_1:=2+\mub$; $\gb:=(1-\vbar^2)^{1/2}((1+\vbar)/(1-\vbar))^{\vbar/2}$ $(=\gmax$ at $v=\vbar)$;
$\underline\alpha:=\arccos(\frac12+\vbar)$; $\Delta_{\mathrm a}:=\pi/3-\underline\alpha$; $d_1:=\vbar\bar t_1+\mub$;
$e_\rho:=\vbar\bar t_1+\Delta_{\mathrm a}+\mub$; $T_2:=\bar t_1+(\sqrt3+e_\rho)/(1-\vbar)$; $d_{\mathrm s}:=\vbar T_2$;
$d_P:=\max(\mub,d_{\mathrm s}+\Delta_{\mathrm a})$.
\item \emph{Run 1.} $\rho_{\mathrm{mid}}:=\sin\frac12$; for $\rho>0$: $W_1(\rho):=e^{\vbar\bar t_1/\rho}$, $\omega_1(\rho):=\vbar\bar t_1/(2\rho)$,
$\vk(\rho):=(2/\rho+1/\rho^2)\vbar\bar t_1e^{2\vbar\bar t_1/\rho}$; $k_{\mathrm{arm}}:=(1+\vbar)/(\cos\frac12-\vbar)+2\vbar/(1-\vbar^2)$;
$d_{\mathrm{arm}}:=(1+\vbar)/((1-\vbar)(\cos\frac12-\vbar))+\vbar/(1-\vbar^2)$;
$e_1:=(2-\cos\frac12)(1+\tan^2\frac14)+(1+\vbar)\tan(\frac14)^{1-\vbar}$, $\mathcal E_1:=\max(1,e_1/(1-\vbar))$,
$W_{\mathrm{arm}}:=\mathcal E_1^{\vbar}/(1-\vbar)$; $k_{\mathrm{all}}:=\max(k_{\mathrm{arm}},1+\vk(\rho_{\mathrm{mid}}))$,
$W_{\mathrm{all}}:=\max(W_{\mathrm{arm}},W_1(\rho_{\mathrm{mid}})(\pi/2)^{2\vbar})$.
\item \emph{Window.} $\theta_L:=\underline\alpha-\eta$, $\theta_R:=\pi/3+\eta$; $R_0(\theta):=(5-4\cos\theta)^{1/2}=|e^{i\theta}-2|$,
$R_L:=R_0(\theta_L)$, $R_R:=R_0(\theta_R)$; $\sr_{\min}:=R_L-d_1$, $\sr_{\max}:=R_R+d_1$; $\rho_{\mathrm w}:=\sin\theta_L$,
$W_{\mathrm w}:=W_1(\rho_{\mathrm w})$, $\omega_{\mathrm w}:=\omega_1(\rho_{\mathrm w})$, $\vk_{\mathrm w}:=\vk(\rho_{\mathrm w})$;
$\iota_{\mathrm w}:=1+\vbar\bar t_1/\sin\theta_L$; $q_{\min}:=\sin(\underline\alpha+\eta/2)/R_R$, $q_{\max}:=\sin(\pi/3-\eta/2)/R_L$;
$\ell_{\min}:=4q_{\min}/(1+d_1/R_L)$, $\ell_{\max}:=2\bar t_1\iota_{\mathrm w}q_{\max}/(1-d_1/R_L)$;
$\Lambda_{\min}:=4\sin(\underline\alpha-\eta/2)/\sr_{\max}^2$, $\Lambda_{\max}:=2\bar t_1\iota_{\mathrm w}\sin(\pi/3+\eta/2)/\sr_{\min}^2$;
$m_0:=\max(|1-2\cos\theta_L|,|1-2\cos\theta_R|)/R_L$, $\omega_{\mathrm p}:=\arcsin(d_1/R_L)$,
$\angle_{\mathrm w}:=\arcsin m_0+\omega_{\mathrm w}+\omega_{\mathrm p}$, $\sA_{\min}:=\cos\angle_{\mathrm w}$, $\sA_{\max}:=W_{\mathrm w}$;
$b_{\mathrm{lo}}:=\Lambda_{\min}(1-\eta^2/48)/4$, $b_{\mathrm{hi}}:=\Lambda_{\max}/4$, $\bar b:=b_{\mathrm{hi}}\eta^2$,
$e_\tau:=\tan(\bar b/2)/(\bar b/2)-1$, $\beta_2:=b_{\mathrm{hi}}(1+e_\tau)/2$;
$C_{\mathrm{hi}}:=2\sr_{\max}\beta_2^{1-\vbar}$, $C_{\mathrm{lo}}:=2\sr_{\min}(b_{\mathrm{lo}}/2)/(1+\tan^2(\bar b/2))$;
$\mathrm{aux}:=\ell_{\max}\bar b\eta/2$, $\sa_{\min}:=\sA_{\min}\cos\bar b-\mathrm{aux}$, $\sa_{\max}:=\sA_{\max}(1+\vbar)+\mathrm{aux}$;
$L_C:=2\Lambda_{\max}/\Lambda_{\min}+\eta\sA_{\max}/\sr_{\min}$,
$L_{\mathrm{lo}}:=2(\cos\bar b-\vbar)(\Lambda_{\min}/\Lambda_{\max})(1-\eta^2/32)$;
$S_1:=\ell_{\max}(\sr_{\max}\sa_{\max}L_C+(1+\vbar)\eta)/(2(1-\vbar^2)\sr_{\min})$.
\item \emph{Forward arc.} $h_{\mathrm f}:=\sA_{\min}-\sr_{\max}b_{\mathrm{hi}}^2\eta^3/2$, $\rho_{\mathrm c}:=h_{\mathrm f}/(1+\vbar)$,
$Q:=(1+\vbar)\sr_{\max}/h_{\mathrm f}$, $y_{\mathrm f}:=2\sr_{\max}\beta_2$;
$Y_1:=(\ell_{\max}/2+\sA_{\max}b_{\mathrm{hi}}(1+e_\tau)\eta)/\sa_{\min}$, $\bar Y:=Y_1\eta^{1-\mub}Q^{\mub}$,
$\theta_{\mathrm f}:=\max((1+\vbar)Y_1,\vbar y_{\mathrm f}/\rho_{\mathrm c})/(1-\vbar-\vbar\bar Y)$,
$W_{\mathrm f}:=\sa_{\max}[(1+\vbar+\vbar\bar Y)^2+((1+\vbar)\bar Y+\vbar)^2]^{1/2}/(1-\vbar^2)$;
$\delta_{\mathrm{fw}}:=\vbar/(3(1-\vbar)^2\xi_1)$, $\tau_{\mathrm f}:=\beta_2Q^{1+\mub}$, $\bar\tau_{\mathrm f}:=\tau_{\mathrm f}\eta^{1-\mub}$,
$\rho_{\mathrm u}:=((1+\vbar)\sA_{\max}+\delta_{\mathrm{fw}})/(1-2\bar\tau_{\mathrm f}^2)$,
$S_2:=(2\tau_{\mathrm f}+\theta_{\mathrm f}Q^{\mub})(\rho_{\mathrm u}\sa_{\max}L_C+1+\vbar)/((1-\vbar^2)\rho_{\mathrm c})$,
$k_{\mathrm f}:=Q^{\mub}(1+\vk_{\mathrm w}+S_1\vbar)+S_2\vbar$.
\item \emph{Virtual labels.} $\theta_{\mathrm v}:=Q^{\mub}e^{\mub}(\theta_{\mathrm f}+\mub e^{2\mub}y_{\mathrm f}\xi_1)$,
$W_{\mathrm v}:=W_{\mathrm f}e^{2\mub}$, $u_{\mathrm v}:=W_{\mathrm v}\xi_1$.
\item \emph{Tail.} $\xi_-:=C_{\mathrm{hi}}^{-1/2}$, $\xi_+:=(\gb/C_{\mathrm{lo}})^{1/(2-2\vbar)}$, $\bar\xi_{\mathrm e}:=\xi_+\epsbar^{1/2}$;
$\delta_{\mathrm{tl}}:=(1+\vbar)\vbar/(3(1-\vbar)\xi_-)$, $h_{\mathrm t}:=\sA_{\min}-\delta_{\mathrm{tl}}$,
$h_{\mathrm T}:=(1+\vbar)\sA_{\max}+\sr_{\max}b_{\mathrm{hi}}^2\eta^3/2$; $m_\tau:=(2h_{\mathrm t}/C_{\mathrm{hi}})\eta^{-(1-2\vbar)}$,
$\tau_{\mathrm t}:=m_\tau^{1/(1+\vbar)}$, $e_{\mathrm t}:=(1+\vbar)/(m_\tau\tau_{\mathrm t}^{1-\vbar})$;
$y_{\mathrm c}:=2(1+e_{\mathrm t})h_{\mathrm T}/(1-\vbar)$, $y_{\mathrm T}:=C_{\mathrm{hi}}^{1/(1+\vbar)}\max(1,y_{\mathrm c})^{\vbar}$,
$T_{\mathrm t}:=(\max(1,y_{\mathrm c}/C_{\mathrm{lo}})\cdot2/b_{\mathrm{lo}})^{\vbar}$;
$Y_2:=T_{\mathrm t}\ell_{\max}/(2\sa_{\min})$, $\bar Y_2:=Y_2\eta^{1-\muTb}$,
$\theta_{\mathrm t}:=Y_2+\vbar y_{\mathrm T}/h_{\mathrm t}$, $z_{\mathrm T}:=(h_{\mathrm T}+\vbar y_{\mathrm T}\eta^{1-\muTb})/(1-\vbar)$,
$\rho_{\mathrm T}:=z_{\mathrm T}(1+(y_{\mathrm T}\eta^{1-\muTb}/h_{\mathrm t})^2)^{1/2}$;
$W_{\mathrm t}:=\sa_{\max}[(1+\vbar+\vbar\bar Y_2)^2+((1+\vbar)\bar Y_2+\vbar)^2]^{1/2}/(1-\vbar^2)$;
$S_3:=(y_{\mathrm T}/h_{\mathrm t}+\theta_{\mathrm t})(\rho_{\mathrm T}\sa_{\max}L_C+1+\vbar)/((1-\vbar^2)h_{\mathrm t})$,
$k_{\mathrm t}:=T_{\mathrm t}(1+\vk_{\mathrm w}+S_1\vbar)+S_3\vbar$; $\xi_2:=1/(3(1-\vbar)\sA_{\min})$;
$h_{\mathrm a}:=1-\sA_{\max}/2-\sr_{\max}b_{\mathrm{hi}}^2\epsbar^3/32$, $Q_{\mathrm a}:=(1+\vbar)\sr_{\max}/h_{\mathrm a}$,
$\tau_{\mathrm a}:=\beta_2Q_{\mathrm a}^{1+\mub}\epsbar^{1-\mub}/4$.
\item \emph{Bulk and hole.} $e_\beta:=\arcsin(e_\rho/\sqrt3)$, $b_{\mathrm l}:=b_{\mathrm{lo}}\eta^2$, $y_{\mathrm b}:=\sin b_{\mathrm l}$,
$H_{\mathrm b}:=3+d_1+(1+\vbar)(\sqrt3+e_\rho)/(1-\vbar)$, $T_{\mathrm b}:=2H_{\mathrm b}/(y_{\mathrm b}\tan(b_{\mathrm l}/2))$,
$S_{\mathrm b}:=(3+\vbar)/(2(1-\vbar)y_{\mathrm b})$, $k_{\mathrm b}:=T_{\mathrm b}^{\vbar}(k_{\mathrm{all}}+S_{\mathrm b}\vbar)+S_{\mathrm b}\vbar$,
$W_{\mathrm b}:=((1+\vbar)/(1-\vbar))T_{\mathrm b}^{\vbar}$; $r_{\mathrm B}:=0.005+\vbar\bar t_1$,
$\rho_{\mathrm h}:=\frac12-d_P-r_{\mathrm B}-\vbar T_2$.
\item \emph{Assembly.} $W_x:=\max(W_{\mathrm v},W_{\mathrm t},W_{\mathrm b}W_{\mathrm{all}})$,
$U_{\max}:=(W_{\mathrm v}+W_{\mathrm t})\eta+W_{\mathrm b}W_{\mathrm{all}}(2\pi-2\eta)/(1-2\vbar)$;
$\hat L:=(8/\lambda_0)^{2\vbar}/(1-2\vbar)$, $\Delta_0:=\lambda_0/(W_x\hat L)$, $c_{\mathrm{sep}}:=2\sin(\Delta_0/2)$,
$\kc:=\min(\cos(\theta_{\mathrm v}\Delta_0^{1-\mub}),h_{\mathrm t}/W_{\mathrm t})$;
$k_*:=\max(k_{\mathrm f}W_{\mathrm v}^{\mub},k_{\mathrm t}W_{\mathrm t}^{\muTb},k_{\mathrm b})$, $c_{\mathrm{far}}:=k_*/\sqrt{12}$,
$c_{\mathrm{near}}:=\theta_{\mathrm v}((12u_{\mathrm v}^2+6u_{\mathrm v}+1)/(3-2\mub))^{1/2}$,
$c_{\mathrm{tail}}:=\theta_{\mathrm t}(19/(3-2\muTb))^{1/2}$, $c_N:=\max(c_{\mathrm{far}},c_{\mathrm{near}},c_{\mathrm{tail}})$;
$c_y:=y_{\mathrm T}/h_{\mathrm t}^2$; $c_{T4}:=\theta_{\mathrm t}/h_{\mathrm t}$.
\end{itemize}}
\end{definition}

\noindent\textbf{Side conditions.}\namedlabel{hyp:SC}{(SC)} {\small
\hyptag{SC1} $\sqrt3-e_\rho>\epsbar$. \hyptag{SC2} $\cos\frac12>\vbar$. \hyptag{SC3} $\theta_L\ge\frac12$. \hyptag{SC4} $\theta_R\le\pi/2$.
\hyptag{SC5} $\eta\le\pi/3$. \hyptag{SC6} $\cos\theta_L<4/5$. \hyptag{SC7} $\cot(\underline\alpha-\eta/2)/2<s(\theta_R)$, where
$s(\theta):=2\sin\theta/(5-4\cos\theta)$. \hyptag{SC8} $d_1<R_L$. \hyptag{SC9} $\angle_{\mathrm w}<\pi/2$. \hyptag{SC10} $\bar b<\pi/2$.
\hyptag{SC11} $\beta_2<1$. \hyptag{SC12} $b_{\mathrm{lo}}<2$. \hyptag{SC13} $\sa_{\min}>0$. \hyptag{SC14} $\cos\bar b>\vbar$.
\hyptag{SC15} $L_{\mathrm{lo}}>\eta\sA_{\max}/\sr_{\min}$. \hyptag{SC16} $0<h_{\mathrm f}\le1-\vbar^2$. \hyptag{SC17} $\rho_{\mathrm c}\xi_1>1$.
\hyptag{SC18} $Q\ge1$. \hyptag{SC19} $1-\vbar-\vbar\bar Y>0$. \hyptag{SC20} $1-2\bar\tau_{\mathrm f}^2>\vbar$. \hyptag{SC21} $C_{\mathrm{hi}}<1$.
\hyptag{SC22} $C_{\mathrm{lo}}\eta^2>\epsbar\gb$. \hyptag{SC23} $\gb>C_{\mathrm{lo}}$. \hyptag{SC24} $\bar\xi_{\mathrm e}<\eta$.
\hyptag{SC25} $h_{\mathrm t}>0$. \hyptag{SC26} $m_\tau>1$. \hyptag{SC27} $\xi_2\vbar\le\frac12$. \hyptag{SC28} $\xi_-\ge\frac12$.
\hyptag{SC29} $h_{\mathrm a}>0$. \hyptag{SC30} $1-2\tau_{\mathrm a}^2>0$. \hyptag{SC31} $\sA_{\min}>4\tau_{\mathrm a}^2$.
\hyptag{SC32} $b_{\mathrm l}\le\pi/6-e_\beta$. \hyptag{SC33} $\pi/3+2e_\beta\le\pi-b_{\mathrm l}$. \hyptag{SC34} $T_{\mathrm b}\ge1$.
\hyptag{SC35} $\vbar\le\frac1{10}$. \hyptag{SC36} $\rho_{\mathrm h}>0$. \hyptag{SC37} $2\gb\epsbar<\theta_L$ and $\gb\epsbar\le1$.
\hyptag{SC38} $2-2d_{\mathrm s}\ge1$. \hyptag{SC39} $2\lambda_0+\min(\bar\xi_{\mathrm e},\lambda_0)\le\eta$. \hyptag{SC40}
$2u_{\mathrm v}\epsbar\le\eta$. \hyptag{SC41} $W_x\hat L\le2$. \hyptag{SC42} $\Delta_0\le\eta$.
\hyptag{SC43} $c_{\mathrm{sep}}>2\epsbar$. \hyptag{SC44} $\theta_{\mathrm v}\Delta_0^{1-\mub}\le\pi/2$.
\hyptag{SC45} $k_*\lambda_0^{1-\muTb}\le\pi$. \hyptag{SC46} $u_{\mathrm v}\ge\frac12$.
\hyptag{SC47} $G_{\mathrm T}>0$ and $\bar\xi_{\mathrm e}+G_{\mathrm T}\le\eta$. \hyptag{SC48} $z_{\mathrm T}\bar\xi_{\mathrm e}+\epsbar<1$ and
$G_{\mathrm T}+z_{\mathrm T}\bar\xi_{\mathrm e}\le1$.}

\medskip
Each side condition is one inequality between explicitly computable numbers; all hold (checked in Appendix~\ref{app:const}).

\begin{lemma}[monotonicity]\label{lem:mono}
Assume \hyptag{SC1}--\hyptag{SC48}, and let $0<v\le\vbar$ and $0<\eps\le\epsbar$. Every constant of
Definition~\ref{def:tdconst} other than $\kc$ is the value at $(\vbar,\epsbar)$ of an
expression $E(v,\eps)$ (replace $\vbar,\mub,\muTb,\epsbar$ by $v,\mu,\muT,\eps$) which is non-decreasing in $v\in(0,\vbar]$ and in
$\eps\in(0,\epsbar]$ (an \emph{upper} constant) or non-increasing in both (a \emph{lower} constant). The lower constants are
$\underline\alpha$, $\rho_{\mathrm w}$, $\sr_{\min}$, $\ell_{\min}$, $\Lambda_{\min}$, $\sA_{\min}$, $b_{\mathrm{lo}}$, $C_{\mathrm{lo}}$,
$\sa_{\min}$, $L_{\mathrm{lo}}$, $h_{\mathrm f}$, $\rho_{\mathrm c}$, $\xi_-$, $h_{\mathrm t}$, $m_\tau$, $\tau_{\mathrm t}$, $h_{\mathrm a}$,
$b_{\mathrm l}$, $y_{\mathrm b}$, $\rho_{\mathrm h}$, $\Delta_0$, $c_{\mathrm{sep}}$, $\theta_L$, $R_L$, $q_{\min}$; all others except $\kc$
are upper constants.
Consequently, wherever a proof uses a constant as an upper (lower) bound for a quantity formed with the actual $v,\mu,\muT,\eps$, the
constant is $\ge$ ($\le$) the value of that expression; and every side condition other than \hyptag{SC44} implies the same
inequality at $(v,\eps)$.
\end{lemma}
\begin{proof}
On $(0,\vbar]\subset(0,\frac1{10}]$ the following are increasing in $v$: $v$, $\mu$, $\muT$, $1+v$, $1/(1-v)$, $1/(1-v^2)$,
$t_1=2+\mu$, $vt_1$ and $\gmax(v)$ ($d\log\gmax/dv=\operatorname{artanh}v>0$, from the closed form in Proposition~\ref{prop:runs});
$1-v$, $\cos\frac12-v$, $\cos\bar b-v$ and $\arccos(\frac12+v)$ are decreasing. Powers are handled by three rules:
(1) $X^{p(v)}$ with $X\ge1$ and $p$ non-decreasing is non-decreasing: $Q^\mu$, $Q^{1+\mu}$, $Q_{\mathrm a}^{1+\mu}$, $\mathcal E_1^v$,
$(\pi/2)^{2v}$, $T_{\mathrm b}^v$, $(8/\lambda_0)^{2v}$, $W_{\mathrm v}^\mu$, $W_{\mathrm t}^{\muT}$,
$(\max(1,y_{\mathrm c}/C_{\mathrm{lo}})2/b_{\mathrm{lo}})^v$ (base $\ge1$ by \hyptag{SC12}), $\max(1,y_{\mathrm c})^v$,
$(\gb/C_{\mathrm{lo}})^{1/(2-2v)}$ (base $>1$ by \hyptag{SC23}); here $Q\ge1$ is \hyptag{SC18}, $Q_{\mathrm a}\ge(1+v)\sr_{\max}\ge1$ as
$h_{\mathrm a}\le1$ and $R_0\ge1$, and $W_{\mathrm v},W_{\mathrm t}\ge\sa_{\max}\ge\sA_{\max}=W_{\mathrm w}\ge1$;
(2) $X^{p(v)}$ with $0<X\le1$ and $p$ non-increasing is non-decreasing: $\beta_2^{1-v}$ (\hyptag{SC11}), $C_{\mathrm{hi}}^{1/(1+v)}$
(\hyptag{SC21}), $\eta^{1-\mu}$, $\eta^{1-\muT}$, $\epsbar^{1-\mu}$, $\tan(\frac14)^{1-v}$, $\lambda_0^{1-\muT}$;
(3) $X^{-p(v)}$ with $0<X<1$ and $p$ non-increasing is non-increasing: $\eta^{-(1-2v)}$ in $m_\tau$; then $m_\tau^{1/(1+v)}$ with
$m_\tau>1$ is a lower constant by \hyptag{SC26}, and $\tau^{v-1}\le\tau_{\mathrm t}^{v-1}\le\tau_{\mathrm t}^{\vbar-1}$ for
$\tau\ge\tau_{\mathrm t}\ge1$ (in $e_{\mathrm t}$). The constant $\kc$ and \hyptag{SC44} enter the proofs only as fixed numbers
(in (T2), $\xi<\Delta_0$ refers to the fixed $\Delta_0$), so no monotonicity is needed for $\kc$ or $\Delta_0^{1-\mu}$. Every constant is built from these blocks, from the fixed data and from earlier
constants by sums, products, quotients of positive quantities, $\max$, $\min$, and the monotone functions $\exp$, $\sqrt\cdot$,
$\arcsin$ on $[0,1]$, $\sin$ and $\tan$ on $[0,\pi/2)$, $\cos$ on $[0,\pi/2]$ and $\cot$ on $(0,\pi/2)$; the arguments stay in these
ranges by \hyptag{SC1}, \hyptag{SC8}--\hyptag{SC10} and $m_0\le1$ (checked in Appendix~\ref{app:const}). Going through Definition~\ref{def:tdconst} in order, each operation combines upper
into upper and lower into lower constants in the direction stated (upper minus lower is upper, upper over lower is upper, etc.).
Three further steps use side conditions: $C_{\mathrm{lo}}$ uses $(b_{\mathrm{lo}}\xi^2/2)^{1-v}\ge(b_{\mathrm{lo}}/2)\xi^{2-2v}$
(\hyptag{SC12}); $\xi_-=C_{\mathrm{hi}}^{-1/2}\le C_{\mathrm{hi}}^{-1/(2-2v)}$ (\hyptag{SC21}); $\bar Y$ and $\theta_{\mathrm f}Q^{\mub}$
use $Q\ge1$. In $\eps$: the constants containing $\epsbar$ are $\bar\xi_{\mathrm e}$, $h_{\mathrm a}$, $Q_{\mathrm a}$, $\tau_{\mathrm a}$,
all monotone in the stated direction. Each side condition other than \hyptag{SC12}, the second part of \hyptag{SC16}, \hyptag{SC18}, \hyptag{SC23}, \hyptag{SC32}--\hyptag{SC34}, \hyptag{SC42} and \hyptag{SC46} compares an upper with a lower expression (or a number) in the direction that is preserved when $(\vbar,\epsbar)$ is replaced by $(v,\eps)$. The exceptions hold for every $v\in(0,\vbar]$ and $\eps\in(0,\epsbar]$ directly. Since $R_0\ge1$ we have $\sr_{\max}\ge1$, hence $b_{\mathrm{lo}}\le\Lambda_{\min}/4\le1$ \hyptag{SC12}, $b_{\mathrm l}\le\eta^2$, and $C_{\mathrm{lo}}\le\sr_{\min}b_{\mathrm{lo}}\le1/\sr_{\max}\le1/R_R<1\le\gb$ \hyptag{SC23}. Since $h_{\mathrm f}\le\sA_{\min}\le1$, $Q\ge\sr_{\max}\ge1$ \hyptag{SC18}. The bound $\eta^2+2e_\beta<\pi/6$ gives \hyptag{SC32} and \hyptag{SC33}. Since $y_{\mathrm b}\tan(b_{\mathrm l}/2)<1<H_{\mathrm b}$, $T_{\mathrm b}>1$ \hyptag{SC34}. Since $W_x,\hat L\ge1$, $\Delta_0\le\lambda_0<\eta$ \hyptag{SC42}. Since $W_{\mathrm v}\ge1$, $u_{\mathrm v}\ge\xi_1>\frac12$ \hyptag{SC46}. Finally, the lower constant $h_{\mathrm f}$ is at most its value $0.99152$ at $v=0$, which is $<1-\vbar^2\le1-v^2$ (second part of \hyptag{SC16}).
\end{proof}

From now on Lemma~\ref{lem:mono} is used tacitly: ``by $v\le\vbar$'' means that a $v$-expression is bounded by its constant.

\subsection{Run 1, tangency and the forward arc}\label{sec:fwd}

\begin{lemma}[run 1]\label{lem:run1}
\begin{enumerate}
\item[(a)] Let $\as\in(0,\arccos v)$ be given by $\sg(\as)=\eps\gmax$. A true candidate $\alpha\in(-\pi,\pi]$ is seen during
$[0,t_1]$ iff $|\alpha|\le\as$. For $\alpha\in[\as,2\pi-\as]$, $\rho_\alpha\ge\eps$ on $[0,t_1]$, so virtual $=$ true on $[0,t_1]$.
Moreover $\as\le2\gb\eps<\theta_L$, and every candidate, true or virtual, satisfies $|S(t)-e^{i\alpha}|\le vt$.
\item[(b)] \emph{Middle.} For $\alpha\in[\frac12,2\pi-\frac12]$, $\rho_\alpha\ge\rho_{\mathrm{mid}}$ on $[0,t_1]$; for
$\alpha\in[\frac12,\frac\pi2]$ (and the mirror labels), $\rho_\alpha\ge\sin\alpha$ on $[0,t_1]$. If $\rho_\alpha\ge\rho$ on $[0,t_1]$,
then at $t_1$: $1\le w_1\le W_1(\rho)$, $|\arg\partial_\alpha S_\alpha(t_1)-(\alpha+\pi/2)|\le\omega_1(\rho)$, $|k_1-1|\le\vk(\rho)$.
\item[(c)] \emph{Arm.} For $\alpha\in(0,\frac12]$: $\rho_1\ge1+v$, $\tau_1>1$, $y(\alpha,t_1)\ge\sin\alpha$; $|k_1|\le k_{\mathrm{arm}}$;
the tangent at $t_1$ is $e^{i(\pi/2+\delta)}$ with $0\le\delta\le\tan\delta\le d_{\mathrm{arm}}\,y(\alpha,t_1)$;
$w_1\le W_{\mathrm{arm}}(2/\alpha)^{2v}$; and $y(\alpha,t_1)\le2\mathcal E_1\tau_\alpha^{(1-v)/(1+v)}$.
\item[(d)] For all $\alpha\in(0,2\pi)$: $|k_1|\le k_{\mathrm{all}}$, $1\le w_1\le W_{\mathrm{all}}(2/\circab{\alpha})^{2v}$ and
$\rho_1\ge1+v$, where $\circab\alpha:=\min(\alpha,2\pi-\alpha)$.
\item[(e)] For every $\alpha\in(0,2\pi)$ the true candidate $S_\alpha$ exists on $[0,\ttd]$ (it does not meet $P$). The candidate
$\alpha=0$ is caught at $1/(1-v)<t_1$.
\end{enumerate}
In (b)--(d) the true front is meant, also at the labels $0<\circab\alpha<\as$, which have been seen but keep moving by the true law.
\end{lemma}
\begin{proof}
(a) By Lemma~\ref{lem:coord}(b),(c), $\tau$ increases from $\tau_\alpha$ to $\tau_1$, and $h_1<0$ means $\cos\beta_1<-v$, so
$\tau_1^2>(1+v)/(1-v)>(1-v)/(1+v)$. If $0<\alpha\le\arccos v$ the minimum point of $\rho$ is passed, and
$\min_{[0,t_1]}\rho=\sg(\alpha)/\gmax$, which is $\le\eps$ iff $\alpha\le\as$ ($\sg$ is increasing on $(0,\arccos v]$); if
$\alpha\in(\arccos v,\pi)$, $\rho$ is increasing and $\min\rho=1$. The label $\pi$ has $\beta\equiv\pi$ and $\rho\ge1$; the lower half is
the mirror image; $\alpha=0$ is caught (see (e)). Virtual $=$ true by Lemma~\ref{lem:truefront}. With $\tau_{\mathrm s}:=\tan(\as/2)\le1$:
$\tau_{\mathrm s}^{1-v}=\eps\gmax(1+\tau_{\mathrm s}^2)/2\le\eps\gmax$, so $\as\le2\tau_{\mathrm s}\le2(\eps\gmax)^{1/(1-v)}\le2\eps\gmax\le2\gb\eps$,
because $\eps\gmax\le\gb\epsbar\le1$; and $2\gb\eps<\theta_L$ (\hyptag{SC37}). Every candidate has speed $\le v$ (true: by definition;
virtual: $|F|\le1$, Lemma~\ref{lem:field}(i)).

(b) For $\alpha\in[\frac12,\frac\pi2]$: $y=C\tau^v$ is non-decreasing, so $\rho\ge y\ge\sin\alpha\ge\sin\frac12$. For
$\alpha\in[\frac\pi2,\pi]$: $X:=\operatorname{Re}S_\alpha$ has $X'=v(X-t)/\rho$ and $(X-t)'\le v-1<0$, $X(0)=\cos\alpha\le0$; so $X-t\le X\le\cos\alpha\le0$,
and with $y\ge\sin\alpha$, $\rho^2\ge\cos^2\alpha+\sin^2\alpha=1$. Mirror for the lower half. By Lemma~\ref{lem:truefront}, $(\log w)_t\in[0,v/\rho]$,
so $1\le w_1\le e^{vt_1/\rho}$; $|\psi_t|\le v/(2\rho)$, and the tangent at $t=0$ is $e^{i(\alpha+\pi/2)}$; and, as
$|1-3\sin^2\!\phi|\le2$ and $|\sin\phi(3\cos^2\!\phi-1)|\le1$ (the maximum of $|s(2-3s^2)|$ on $[-1,1]$ is $1$),
$\frac d{dt}|k-1|\le(2v/\rho)|k-1|+2v/\rho+v/\rho^2$, whence by Gronwall ($k(0)=1$) $|k_1-1|\le(2/\rho+1/\rho^2)vt_1e^{2vt_1/\rho}$. Since
$vt_1\le\vbar\bar t_1$, these are the bounds claimed.

(c) $h_1=-(1-v^2)-(1-\cos\alpha)<0$, so $\cos\beta_1<-v$, $\beta_1>\pi/2$, $\tau_1>1$; and $|h_1|=\rho_1|\cos\beta_1+v|\le\rho_1(1-v)$ gives
$\rho_1\ge1+v$. $y=C\tau^v$ is non-decreasing and $y(0)=\sin\alpha$.

\emph{Curvature.} Apply Lemma~\ref{lem:transport}(c) at $t=0$, where $w=1$, $k=1$, $\rho=1$, $\psi=\alpha+\pi/2$, $\phi=-\pi/2$ and
$\sa\,\partial_\alpha\log C=-\sin\alpha\cdot(\cos\alpha-v)/\sin\alpha$: there
$\TS=-\frac v{1-v^2}[-(\cos\alpha-v)-v-\cos(-\pi-\alpha)]=0$, so $\tau_\alpha^vJ=1-v^2$ and
\[
k_1=\frac{\mathcal T}{w_1^3}-\TS_1,\qquad\mathcal T:=\Bigl(\frac{\tau_1}{\tau_\alpha}\Bigr)^v\ge1 .
\]
By Lemma~\ref{lem:transport}(a), $(1+v)w_1\ge\operatorname{Im}\Xi=\tau_1^v\sfc=(\cos\alpha-v)\mathcal T>0$ (\hyptag{SC2}), and $w_1\ge1$, so
$0\le\mathcal T/w_1^3\le(1+v)/(\cos\alpha-v)\le(1+v)/(\cos\frac12-v)$. In $\TS_1$, $|\sa\,\partial_\alpha\log C|\le1$, so by
Lemma~\ref{lem:transport}(e) $|\TS_1|\le\frac v{(1-v^2)\rho_1}[\rho_1+1+v]\le\frac{2v}{1-v^2}$. Hence $|k_1|\le k_{\mathrm{arm}}$.

\emph{Tilt.} By Lemma~\ref{lem:transport}(a), with $\sa=-\sin\alpha<0$ and $\sfc\tau_1^v=(\cos\alpha-v)\mathcal T>0$,
\[
(1-v^2)\,\partial_\alpha S_\alpha(t_1)=\sa[(1-v\cos\beta_1)-iv\sin\beta_1]+\sfc\tau_1^v[-v\sin\beta_1+i(1+v\cos\beta_1)]
\]
has negative real part and positive imaginary part; so the tangent is $e^{i(\pi/2+\delta)}$ with $\delta\in[0,\pi/2)$ and
\[
\tan\delta=\frac{\sin\alpha(1-v\cos\beta_1)+v\sin\beta_1(\cos\alpha-v)\mathcal T}{v\sin\alpha\sin\beta_1+(1+v\cos\beta_1)(\cos\alpha-v)\mathcal T}
\le\frac{(1+v)\sin\alpha}{(1-v)(\cos\alpha-v)\mathcal T}+\frac{v\sin\beta_1}{1-v};
\]
with $\mathcal T\ge1$, $\sin\alpha\le y$, $\sin\beta_1=y/\rho_1\le y/(1+v)$, $\cos\alpha\ge\cos\frac12$: $\tan\delta\le d_{\mathrm{arm}}y$.

\emph{Stretch.} $|\Xi|^2=\sin^2\!\alpha+(\cos\alpha-v)^2\mathcal T^2\le\mathcal T^2$, so $w_1\le\mathcal T/(1-v)$. By
Lemma~\ref{lem:coord}(d) and $\tau_1>1$: $(1-v)\tau_1^{1+v}\le2|h_1|/C+1+v$, with $|h_1|\le2-\cos\frac12$,
$C=\sg(\alpha)\ge2\tau_\alpha^{1-v}/(1+\tan^2\frac14)$ and $1+v\le(1+v)\tan(\frac14)^{1-v}\tau_\alpha^{-(1-v)}$ ($\tau_\alpha\le\tan\frac14$).
So $\tau_1^{1+v}\le\mathcal E_1\tau_\alpha^{-(1-v)}$, $\tau_1/\tau_\alpha\le\mathcal E_1^{1/(1+v)}\tau_\alpha^{-2/(1+v)}$ and
$\mathcal T\le\mathcal E_1^{v}(2/\alpha)^{2v}$ ($\mathcal E_1\ge1$, $\tau_\alpha\ge\alpha/2$, $2/\alpha>1$). Hence
$w_1\le W_{\mathrm{arm}}(2/\alpha)^{2v}$. Finally $y=C\tau_1^v\le2\tau_\alpha^{1-v}(\mathcal E_1\tau_\alpha^{-(1-v)})^{v/(1+v)}\le
2\mathcal E_1\tau_\alpha^{(1-v)/(1+v)}$.

(d) If $\circab\alpha\ge\frac12$, (b) with $\rho=\rho_{\mathrm{mid}}$ gives $|k_1|\le1+\vk(\rho_{\mathrm{mid}})$ and
$w_1\le W_1(\rho_{\mathrm{mid}})(\pi/2)^{2v}(2/\pi)^{2v}\le W_{\mathrm{all}}(2/\circab\alpha)^{2v}$. If $\circab\alpha<\frac12$, use (c) and
the mirror symmetry. $\rho_1\ge1+v$: the argument of (c) uses only $h_1\le-(1-v^2)$, which holds for all $\alpha$ (for $\alpha=\pi$,
$\beta\equiv\pi$ and $\rho_1=1+(1+v)t_1$).

(e) By Lemma~\ref{lem:coord}(c), on a straight run $\rho\ge C/\gmax>0$ whenever $0<\beta<\pi$ and $C>0$. In run 1,
$C=\sg(\alpha)>0$ for $\alpha\in(0,\pi)$ (mirror for $(\pi,2\pi)$), and the label $\pi$ has $\rho\ge1$. In run 2,
$\sbb\in(0,\pi)$ for $\alpha\ne\alpha_0$ (Lemma~\ref{lem:tang}(a); for $\alpha=\pi$, $\sbb=\pi-\beta^*$), so $C=\sr\sg(\sbb)>0$; and the
label $\alpha_0$ has $\sbb=0$, $\rho=\rho_0-(1-v)(t-t_1)>0$ for $t<\tc$, while $\ttd<\tc$ (Lemma~\ref{lem:tang}(b), whose proof does not
use (e)). The label $0$ lies on the axis ahead of $P$: $\rho(0)=1$, $\rho'=-(1-v)$; it is caught, a fortiori seen, at $1/(1-v)$.
\end{proof}

\begin{lemma}[tangency; data at $t_1$]\label{lem:tang}
\begin{enumerate}
\item[(a)] For $\theta\in(0,\pi)$: $\rho_1\,\partial_\theta\beta_1=\sin\beta_1\cdot t_1(\cos\alpha_0-\cos\theta)/\sin\theta$. Hence $\beta_1$ has its
unique minimum on $(0,\pi)$ at $\alpha_0$; there $\sbb=\beta_1-\beta^*\in[0,\pi/2)$; on the lower half $\sbb\in(0,\pi)$; and
\[
\partial_\theta\sbb=\Lambda(\theta)\sin\frac{\theta-\alpha_0}2,\qquad
\Lambda(\theta):=\frac{2t_1\sin\beta_1\sin((\theta+\alpha_0)/2)}{\rho_1\sin\theta}>0\qquad(0<\theta<\pi).
\]
So $\sbb$ has a double zero at $\alpha_0$: the true front at $t_1$ is tangent at $S_{\alpha_0}(t_1)$ to the ray of run 2.
\item[(b)] $\alpha_0\in[\underline\alpha,\pi/3]$; $|\rho_0-\sqrt3|\le e_\rho$; $\rho_0>\eps$; $\ttd\le\tc\le T_2$,
$0<\tc-\ttd\le v\eps/(3(1-v)^2)$ and $\ttd-\ta\le\eps/(1-v)$; every candidate, true or virtual, satisfies
$|x_\theta(t)-e^{i\theta}|\le vt\le d_{\mathrm s}$ for $t\le\ttd$; $P([t_1,\ttd])$ lies within $d_P$ of the segment $[2,e^{i\pi/3}]$; the
virtual candidate $\alpha_0$ stays on the ray of run 2, $\rho_{\alpha_0}(\ta)=\eps$ and $P(\ttd)=x_{\alpha_0}(\ttd)$.
\item[(c)] $I_\eta\subset[\theta_L,\theta_R]\subset[\frac12,\frac\pi2]$, and on $I_\eta$:
\begin{gather*}
\sr_{\min}\le\sr\le\sr_{\max},\quad\Lambda_{\min}\le\Lambda\le\Lambda_{\max},\quad\ell_{\min}\le\sr\Lambda\le\ell_{\max},\quad
\sA_{\min}\le\sA\le\sA_{\max},\\
1\le w_1\le W_{\mathrm w},\quad|k_1-1|\le\vk_{\mathrm w},\quad b_{\mathrm{lo}}\xi^2\le\sbb\le b_{\mathrm{hi}}\xi^2\le\bar b,\quad
\sbb/2\le\tau_{\sbb}\le\beta_2\xi^2,\\
C_{\mathrm{lo}}\xi^{2-2v}\le C\le C_{\mathrm{hi}}\xi^{2-2v},\quad\sa_{\min}\le\sa\le\sa_{\max},\quad|\partial_\theta\log C|\le L_C/\xi;
\end{gather*}
for $\theta<\alpha_0$: $\partial_\theta\log C\le-(L_{\mathrm{lo}}-\eta\sA_{\max}/\sr_{\min})/\xi<0$, so $C$ is strictly monotone in $\xi$ on
that side; and the quantity $\TS$ of Lemma~\ref{lem:transport} at $t_1$, in the frame of run 2, satisfies $|\TS|\le S_1v$.
\item[(d)] For $\theta\in(0,\pi)\setminus\{\alpha_0\}$, at $t_1$ in the frame of run 2,
\[
\tau_{\sbb}^v\,\partial_\theta C=\frac{2\tau_{\sbb}\sA}{1+\tau_{\sbb}^2}+\sr\,\partial_\theta\sbb\,(\cos\sbb-v).
\]
Consequently, on $I_\eta$: for $\theta>\alpha_0$, $0<\tau_{\sbb}^v\partial_\theta C\le\xi(\ell_{\max}/2+\sA_{\max}b_{\mathrm{hi}}(1+e_\tau)\eta)$;
for $\theta<\alpha_0$, $\partial_\theta C<0$ and $|\tau_{\sbb}^v\partial_\theta C|\le\sr|\partial_\theta\sbb|\le\ell_{\max}\xi/2$.
\end{enumerate}
\end{lemma}
\begin{proof}
(a) From $C=\rho\sg(\beta)$, $h=\rho(\cos\beta+v)$ and $\sg'/\sg=(\cos\beta-v)/\sin\beta$, on any straight run
\[
(\cos\beta+v)\,\partial_\theta C-\sg\,\partial_\theta h=\rho\,\partial_\theta\beta\cdot\sg(1-v^2)/\sin\beta .
\]
In run 1 at $t_1$: $\partial_\theta C=\sg(\theta)(\cos\theta-v)/\sin\theta$, $\partial_\theta h=-\sin\theta$, $\sg(\theta)/\sg(\beta_1)=\rho_1$ and
$\rho_1(\cos\beta_1+v)=h_1$; so $(1-v^2)\rho_1\partial_\theta\beta_1=\sin\beta_1[h_1(\cos\theta-v)+\sin^2\!\theta]/\sin\theta$, and
$h_1(\cos\theta-v)+\sin^2\!\theta=(1-v^2)(1-t_1(\cos\theta-v))=(1-v^2)t_1(\cos\alpha_0-\cos\theta)$, as $t_1(\cos\alpha_0-v)=1$. Since
$h_1<0$, $\beta_1\in(\arccos(-v),\pi)$, so $\beta^*>\pi/2$ and on the upper half $\sbb=\beta_1-\beta^*\in[0,\pi/2)$. The mirror point
of $S_\alpha$ has polar angle $-\beta_1(\alpha)$ about the axis of run 1, i.e.\ $\sbb=2\pi-\beta_1(\alpha)-\beta^*\in(\pi-\beta^*,2\pi-2\beta^*]\subset(0,\pi)$.
The formula for $\partial_\theta\sbb$ is $\cos\alpha_0-\cos\theta=2\sin\frac{\theta+\alpha_0}2\sin\frac{\theta-\alpha_0}2$. In the frame of run 2,
$\partial_\theta(S_\theta-P)(t_1)=(\sA+i\sr\,\partial_\theta\sbb)e^{i\sbb}$, which is real at $\alpha_0$.

(b) $t_1=(2-v)/(1-v)$, so $\cos\alpha_0-\frac12=v(3-2v)/(2(2-v))\in[0,v]$ and $\alpha_0\in[\underline\alpha,\pi/3]$. By Lemma~\ref{lem:run1}(a),
$|S_{\alpha_0}(t_1)-P(t_1)-(e^{i\pi/3}-2)|\le vt_1+|e^{i\alpha_0}-e^{i\pi/3}|+|t_1-2|\le vt_1+\Delta_{\mathrm a}+\mu\le e_\rho$, and
$|e^{i\pi/3}-2|=\sqrt3$; so $|\rho_0-\sqrt3|\le e_\rho$, $\rho_0\ge\sqrt3-e_\rho>\eps$ (\hyptag{SC1}), and $\tc=t_1+\rho_0/(1-v)\le T_2$.
The true candidate $\alpha_0$ has $\sbb=0$: it stays on the ray, at distance $\rho_0-(1-v)(t-t_1)$ (Proposition~\ref{prop:runs}), which
is $\eps$ at $\ta$; up to $\ta$ it is the virtual one (Lemma~\ref{lem:truefront}). For $t\ge\ta$, $x=P+\sigma\nhat$ solves
$x'=vF(x-P)$ iff $\sigma'=-1+vf(\sigma)$, $\sigma(\ta)=\eps$; by uniqueness this is the virtual candidate $\alpha_0$, and $\sigma$ reaches
$0$ after the time $\int_0^\eps d\sigma/(1-vf(\sigma))=\eps\int_0^1d\sigma/(1-v\sigma(2-\sigma))=\ttd-\ta\le\eps/(1-v)$. Hence
$P(\ttd)=x_{\alpha_0}(\ttd)$, and, as $f<1$ on $[0,\eps)$ and $\int_0^\eps(1-f)=\eps/3$,
$\tc-\ttd=\int_0^\eps\frac{v(1-f)}{(1-v)(1-vf)}\,d\sigma\in(0,v\eps/(3(1-v)^2)]$. Positions: speed $\le v$ and $t\le\ttd\le T_2$. Finally
$P(t_1)=2+\mu$ is within $\mu$ of the segment, $P(\ttd)=x_{\alpha_0}(\ttd)$ is within $vT_2+\Delta_{\mathrm a}\le d_{\mathrm s}+\Delta_{\mathrm a}$
of $e^{i\pi/3}$, and the distance to a convex set is convex along the segment $P([t_1,\ttd])$.

(c) $I_\eta\subset[\theta_L,\theta_R]$ by (b), $\subset[\frac12,\frac\pi2]$ by \hyptag{SC3}, \hyptag{SC4}. Window labels are $\ge\theta_L>\as$,
and $\rho_\theta\ge\sin\theta\ge\rho_{\mathrm w}$ on $[0,t_1]$, so Lemma~\ref{lem:run1}(b) with $\rho=\rho_{\mathrm w}$ gives the bounds on
$w_1$, $k_1$, and $\partial_\theta S_\theta(t_1)=w_1e^{i\psi_1}$ with $|\psi_1-(\theta+\pi/2)|\le\omega_{\mathrm w}$.

\emph{$\sr$.} $\mathsf p:=S_\theta(t_1)-P(t_1)=e^{i\theta}-2+\mathsf d$ with $|\mathsf d|\le vt_1+\mu\le d_1$, and $R_0$ is increasing on $(0,\pi)$.

\emph{$\Lambda$.} $\sin\beta_1/\rho_1=\operatorname{Im}\mathsf p/\sr^2$ with $\operatorname{Im}\mathsf p=\operatorname{Im}S_\theta(t_1)\in[\sin\theta,\sin\theta+vt_1]$
($y$ non-decreasing in run 1, speed $\le v$); so $\Lambda=2t_1(\operatorname{Im}\mathsf p/\sin\theta)\sin((\theta+\alpha_0)/2)/\sr^2$ with
$1\le\operatorname{Im}\mathsf p/\sin\theta\le\iota_{\mathrm w}$, $t_1\in[2,\bar t_1]$ and
$(\theta+\alpha_0)/2\in[\underline\alpha-\eta/2,\pi/3+\eta/2]\subset(0,\pi/2]$ (\hyptag{SC3}, \hyptag{SC5}), where $\sin$ is increasing.

\emph{$\sr\Lambda$} $=2t_1(\operatorname{Im}\mathsf p/\sin\theta)\,\tilde q(\theta)\,(R_0/\sr)$ with $\tilde q:=\sin((\theta+\alpha_0)/2)/R_0$. Here
$\frac d{d\theta}\log\tilde q=\frac12\cot\frac{\theta+\alpha_0}2-s(\theta)$, $s=R_0'/R_0$; $s'$ has the sign of $5\cos\theta-4<0$
(\hyptag{SC6}), so $s\ge s(\theta_R)$, while $\frac12\cot\frac{\theta+\alpha_0}2\le\frac12\cot(\underline\alpha-\eta/2)$; by \hyptag{SC7},
$\tilde q$ is decreasing on $I_\eta$. Hence $\tilde q\in[\tilde q(\alpha_0+\eta),\tilde q(\alpha_0-\eta)]\subset[q_{\min},q_{\max}]$, and
$R_0/\sr\in[1/(1+d_1/R_L),1/(1-d_1/R_L)]$ (\hyptag{SC8}).

\emph{$\sA$} $=\partial_\theta S_\theta(t_1)\cdot\hat{\mathsf p}=w_1\cos\angle(e^{i\psi_1},\hat{\mathsf p})$. With $\mathsf p_0:=e^{i\theta}-2$:
$ie^{i\theta}\cdot\mathsf p_0=2\sin\theta>0$ and $ie^{i\theta}\times\mathsf p_0=2\cos\theta-1$, so
$\angle(ie^{i\theta},\mathsf p_0)=\arcsin(|2\cos\theta-1|/R_0)\le\arcsin m_0$; and $\angle(\mathsf p_0,\mathsf p)\le\arcsin(|\mathsf d|/|\mathsf p_0|)\le\omega_{\mathrm p}$.
So $\angle(e^{i\psi_1},\hat{\mathsf p})\le\angle_{\mathrm w}<\pi/2$ (\hyptag{SC9}) and $\sA\in[\cos\angle_{\mathrm w},W_{\mathrm w}]$.

\emph{$\sbb$} $(\alpha_0\pm\xi)=\int_0^\xi\Lambda(\alpha_0\pm\xi')\sin(\xi'/2)\,d\xi'\in[\Lambda_{\min},\Lambda_{\max}]\cdot4\sin^2(\xi/4)$, and
$(\xi^2/4)(1-\xi^2/48)\le4\sin^2(\xi/4)\le\xi^2/4$. Since $\tan s/s$ is increasing on $(0,\pi/2)$ and $\sbb\le\bar b<\pi/2$ (\hyptag{SC10}):
$\sbb/2\le\tau_{\sbb}\le(\sbb/2)(1+e_\tau)\le\beta_2\xi^2$.

\emph{$C$} $=2\sr\tau_{\sbb}^{1-v}/(1+\tau_{\sbb}^2)$: $C\le2\sr_{\max}(\beta_2\xi^2)^{1-v}\le C_{\mathrm{hi}}\xi^{2-2v}$ (\hyptag{SC11}), and
$C\ge2\sr_{\min}(b_{\mathrm{lo}}\xi^2/2)^{1-v}/(1+\tan^2(\bar b/2))\ge C_{\mathrm{lo}}\xi^{2-2v}$ (\hyptag{SC12}).

\emph{$\sa$} $=\sA(\cos\sbb+v)-\sr\sin\sbb\,\partial_\theta\sbb$, and $|\sr\sin\sbb\,\partial_\theta\sbb|\le\bar b\,\ell_{\max}\xi/2\le\mathrm{aux}$, as
$|\partial_\theta\sbb|\le\Lambda\sin(\xi/2)\le\Lambda\xi/2$.

\emph{$\partial_\theta\log C$} $=\sA/\sr+(\cos\sbb-v)\,\partial_\theta\sbb/\sin\sbb$. Now
\[
\frac{|\partial_\theta\sbb|}{\sbb}=\frac{\Lambda(\theta)\sin(\xi/2)}{\int_0^\xi\Lambda\sin(\xi'/2)d\xi'}\in
\Bigl[\frac{\Lambda_{\min}}{\Lambda_{\max}},\frac{\Lambda_{\max}}{\Lambda_{\min}}\Bigr]\cdot\frac12\cot\frac\xi4,
\]
and for $0<y\le1$, $1/y-y/2\le\cot y\le1/y$ (the function $y\cos y-\sin y+\frac{y^2}2\sin y$ vanishes at $0$ and has derivative
$\frac{y^2}2\cos y\ge0$), so $(2/\xi)(1-\xi^2/32)\le\frac12\cot(\xi/4)\le2/\xi$. Moreover
$(\cos\sbb-v)\sbb/\sin\sbb\in[\cos\bar b-v,\sbb\cot\sbb]\subset[\cos\bar b-v,1]$, and $\cos\bar b-v>0$ (\hyptag{SC14}). Hence
$(\cos\sbb-v)|\partial_\theta\sbb|/\sin\sbb\in[L_{\mathrm{lo}}/\xi,\,2\Lambda_{\max}/(\Lambda_{\min}\xi)]$, while
$0<\sA/\sr\le\sA_{\max}/\sr_{\min}\le(\eta\sA_{\max}/\sr_{\min})/\xi$. For $\theta>\alpha_0$ both terms are positive and
$\partial_\theta\log C\le L_C/\xi$; for $\theta<\alpha_0$, $\partial_\theta\log C\le\sA_{\max}/\sr_{\min}-L_{\mathrm{lo}}/\xi\le
-(L_{\mathrm{lo}}-\eta\sA_{\max}/\sr_{\min})/\xi<0$ (\hyptag{SC15}), and $\ge-2\Lambda_{\max}/(\Lambda_{\min}\xi)\ge-L_C/\xi$.

\emph{$\TS$ at $t_1$.} The tangent is the direction of $(\sA+i\sr\,\partial_\theta\sbb)e^{i\sbb}$, so $|\sin\phi|=\sr|\partial_\theta\sbb|/w_1\le\ell_{\max}\xi/2$;
by Lemma~\ref{lem:transport}(e) with $\rho=\sr\ge\sr_{\min}$ and $w_1\ge1$:
$|\TS|\le\frac{v\ell_{\max}\xi}{2(1-v^2)\sr_{\min}}[\sr_{\max}\sa_{\max}L_C/\xi+1+v]\le S_1v$ ($\xi\le\eta$).

(d) $\partial_\theta C=\sg(\sbb)\sA+\sr\sg'(\sbb)\partial_\theta\sbb$ with $\sg'=\sg(\cos\beta-v)/\sin\beta$, and $\tau^v\sg(\beta)=\sin\beta=2\tau/(1+\tau^2)$
(Lemma~\ref{lem:coord}(a)); this is the identity. On the forward side both terms are positive,
$2\tau_{\sbb}\sA/(1+\tau_{\sbb}^2)\le2\beta_2\xi^2\sA_{\max}=\sA_{\max}b_{\mathrm{hi}}(1+e_\tau)\xi^2\le\sA_{\max}b_{\mathrm{hi}}(1+e_\tau)\eta\,\xi$, and
$0<\sr\,\partial_\theta\sbb(\cos\sbb-v)\le\sr\Lambda\xi/2\le\ell_{\max}\xi/2$. On the tail side $\partial_\theta\log C<0$ by (c), so
$|\tau_{\sbb}^v\partial_\theta C|=\sr|\partial_\theta\sbb|(\cos\sbb-v)-2\tau_{\sbb}\sA/(1+\tau_{\sbb}^2)\le\sr|\partial_\theta\sbb|$.
\end{proof}

\begin{lemma}[forward arc]\label{lem:fwd}
Let $\theta=\alpha_0+\xi$, $0<\xi\le\eta$. For every $t\in[t_1,\tc)$, in the frame of run 2:
\begin{enumerate}
\item[(a)] $h(t)\ge h_{\mathrm f}\xi+(1-v^2)\omega(t)\ge h_{\mathrm f}(\xi+\omega(t))>0$ and $\rho\ge h/(1+v)\ge\rho_{\mathrm c}(\xi+\omega)$;
in particular $\rho>0$ and $0<\beta<\pi$.
\item[(b)] $\tau\le\tau_{\sbb}((1+v)\sr/h)^{1+\mu}$, $(\tau/\tau_{\sbb})^v\le(Q/(\xi+\omega))^\mu$, and
$0<y\le y_{\mathrm f}\xi^2(Q/(\xi+\omega))^\mu$.
\item[(c)] $0<Y_r\le Y_1\xi(Q/(\xi+\omega))^\mu\le\bar Y$; $|\psi|\le|\tan\psi|\le\theta_{\mathrm f}\xi(Q/(\xi+\omega))^\mu\le
\theta_{\mathrm f}Q^{\mub}\xi^{1-\mu}$; $1\le w\le W_{\mathrm f}$.
\item[(d)] At $\ttd$: $\tau\le\tau_{\mathrm f}\xi^{1-\mu}\le\bar\tau_{\mathrm f}$, and $\rho(t)\ge\rho(\ttd)\ge\rho_{\mathrm c}\xi$ for all
$t\in[t_1,\ttd]$; $\rho\ge\eps$ on $[0,t_1]$. Hence every label with $\rho_{\mathrm c}\xi>\eps$, in particular every
$\xi\ge\xi_1\eps$ and a neighbourhood of $\xi=\xi_1\eps$ (\hyptag{SC17}), has $\rho>\eps$ on $[t_1,\ttd]$: it is true up to $\ttd$,
virtual $=$ true there, and the front is $C^\infty$ near it. For $\xi_1\eps\le\xi\le\eta$ moreover $\rho(\ttd)\le\rho_{\mathrm u}\xi$
and $|k|\le k_{\mathrm f}\xi^{-\mu}$ at $\ttd$.
\item[(e)] The true front at every $t\in[t_1,\tc)$ is tangent to $\nhat$ at the candidate $\alpha_0$.
\end{enumerate}
\end{lemma}
\begin{proof}
(a) $\sr-\rho_0=\int_{\alpha_0}^\theta\sA\ge\sA_{\min}\xi$ and $\sr(1-\cos\sbb)\le\sr_{\max}\sbb^2/2\le(\sr_{\max}b_{\mathrm{hi}}^2\eta^3/2)\xi$
(Lemma~\ref{lem:tang}(c)); insert into \eqref{eq:h2} and use $h_{\mathrm f}\le1-\vbar^2\le1-v^2$ (\hyptag{SC16}). $h=\rho(\cos\beta+v)\le(1+v)\rho$.
$\rho>0$ and Proposition~\ref{prop:runs} give $0<\beta<\pi$ (as $0<\sbb<\pi$ at $t_1$).

(b) \eqref{eq:tau2}, $v(1+\mu)=\mu$, $(1+v)\sr/h\le Q/(\xi+\omega)$ by (a), $\tau_{\sbb}\le\beta_2\xi^2$ and $\sr\le\sr_{\max}$.

(c) By Lemma~\ref{lem:transport}(b) and Lemma~\ref{lem:tang}(d), $Y_r=(\tau/\tau_{\sbb})^v\tau_{\sbb}^v\partial_\theta C/\sa\in(0,Y_1\xi(Q/(\xi+\omega))^\mu]$
($\sa\ge\sa_{\min}>0$, \hyptag{SC13}), and $Y_1\xi(Q/(\xi+\omega))^\mu\le Y_1Q^\mu\xi^{1-\mu}\le\bar Y$ (\hyptag{SC18}). So
$N_\psi\in[-v\sin\beta,(1+v)Y_r]$ and $D_\psi\ge1-v-v\bar Y>0$ (\hyptag{SC19}); hence $|\psi|\le|\tan\psi|$. With
$\sin\beta=y/\rho\le(y_{\mathrm f}/\rho_{\mathrm c})\xi(Q/(\xi+\omega))^\mu$:
$|\tan\psi|\le\max((1+v)Y_1,vy_{\mathrm f}/\rho_{\mathrm c})\,\xi(Q/(\xi+\omega))^\mu/(1-v-v\bar Y)\le\theta_{\mathrm f}\xi(Q/(\xi+\omega))^\mu$.
$w=\sa|D_\psi+iN_\psi|/(1-v^2)\le W_{\mathrm f}$, since $|D_\psi|\le1+v+v\bar Y$, $|N_\psi|\le(1+v)\bar Y+v$; $w\ge1$ by Lemma~\ref{lem:truefront}.

(d) By (b) with $h\ge h_{\mathrm f}\xi$: $\tau(\ttd)\le\beta_2\xi^2(Q/\xi)^{1+\mu}\le\tau_{\mathrm f}\xi^{1-\mu}\le\bar\tau_{\mathrm f}$, so
$\cos\beta\ge1-2\bar\tau_{\mathrm f}^2>v$ (\hyptag{SC20}). As $\tau$ increases in time, $\cos\beta>v$ on $[t_1,\ttd]$ and $\rho$ is
decreasing there (Lemma~\ref{lem:coord}(c)): $\rho(t)\ge\rho(\ttd)\ge\rho_{\mathrm c}\xi$. In run 1, $\rho\ge\eps$ since
$\theta\ge\theta_L>\as$ (Lemma~\ref{lem:run1}(a)). Virtual $=$ true by Lemma~\ref{lem:truefront}. For $\xi\ge\xi_1\eps$:
$\omega(\ttd)\le v\eps/(3(1-v)^2)\le\delta_{\mathrm{fw}}\xi$, so by \eqref{eq:h2} $h\le(1+v)\sA_{\max}\xi+\delta_{\mathrm{fw}}\xi$, and
$\rho=h/(\cos\beta+v)\le h/(1-2\bar\tau_{\mathrm f}^2)\le\rho_{\mathrm u}\xi$.

\emph{Curvature.} Lemma~\ref{lem:transport}(d) with $t'=t_1$, $t''=\ttd$: $(\tau/\tau_{\sbb})^v\le(Q/\xi)^\mu\le Q^{\mub}\xi^{-\mu}$,
$|k_1|\le1+\vk_{\mathrm w}$ and $|\TS'|\le S_1v$ (Lemma~\ref{lem:tang}(c)). For $\TS''$, Lemma~\ref{lem:transport}(e) with
$|\sin\phi|=|\sin(\beta-\psi)|\le\sin\beta+|\psi|\le2\tau+|\psi|\le(2\tau_{\mathrm f}+\theta_{\mathrm f}Q^{\mub})\xi^{1-\mu}$,
$\rho\ge\rho_{\mathrm c}\xi$ and $\rho|\sa||\partial_\theta\log C|/w^2\le\rho_{\mathrm u}\xi\,\sa_{\max}L_C/\xi$ gives $|\TS''|\le S_2v\xi^{-\mu}$.
So $|k|\le Q^{\mub}\xi^{-\mu}(1+\vk_{\mathrm w}+S_1v)+S_2v\xi^{-\mu}\le k_{\mathrm f}\xi^{-\mu}$.

(e) For $t<\tc$ the true front is smooth near the label $\alpha_0$ ($\rho_{\alpha_0}(t)>0$), and by (c) $\psi\to0$ as $\xi\downarrow0$ with
$\omega(t)>0$ fixed.
\end{proof}

\begin{lemma}[virtual labels next to $X_0$]\label{lem:virt}
At $\ttd$, for every $\theta=\alpha_0+\xi$, $0\le\xi\le\eta$, the virtual front has tangent $e^{i\psi}\nhat$ with
$|\psi(\theta,\ttd)|\le\theta_{\mathrm v}\xi^{1-\mu}$, and $1\le w\le W_{\mathrm v}$; $\psi(\alpha_0,\ttd)=0$. For $\xi\ge\xi_1\eps$ the
virtual front is the true one.
\end{lemma}
\begin{proof}
For $\xi_1\eps\le\xi\le\eta$: Lemma~\ref{lem:fwd}(c),(d) at $\ttd$, as $\theta_{\mathrm f}Q^{\mub}\le\theta_{\mathrm v}$ and
$W_{\mathrm f}\le W_{\mathrm v}$. Let $0<\xi<\xi_1\eps$; then $\xi<\eta$, since $\xi_1\eps\le u_{\mathrm v}\epsbar\le\eta/2$ (\hyptag{SC40},
$W_{\mathrm v}\ge1$). On $[0,t_1]$ the label is true with $\rho\ge\eps$ (Lemma~\ref{lem:run1}(a)). On $[t_1,\ta]$, by \eqref{eq:h2},
$h_\theta(t)-h_{\alpha_0}(t)=(1+v)(\sr-\rho_0)-\sr(1-\cos\sbb)\ge h_{\mathrm f}\xi>0$ (proof of Lemma~\ref{lem:fwd}(a)) and
$h_{\alpha_0}(t)=(1+v)\rho_{\alpha_0}(t)\ge(1+v)\eps$ (the candidate $\alpha_0$ is on the ray), so $(1+v)\rho_\theta\ge h_\theta>(1+v)\eps$.
Hence virtual $=$ true on $[0,\ta]$ for all labels in $(\alpha_0,\alpha_0+\xi_1\eps)$ (Lemma~\ref{lem:truefront}), and by
Lemma~\ref{lem:fwd}(b),(c) at $\ta$ ($\omega(\ta)=\eps/(1-v)\ge\eps$):
$|\psi(\ta)|\le\theta_{\mathrm f}\xi(Q/\xi)^\mu$, $0\le y(\ta)\le y_{\mathrm f}\xi^2(Q/\xi)^\mu$, $w(\ta)\le W_{\mathrm f}$. On
$[\ta,\ttd]$, of length $I\le\eps/(1-v)$ (so $vI/\eps\le\mu$), the virtual field acts. With $r=x_\theta-P=(z+iy)\nhat$:
$r'=vF(r)-\nhat$, so $y'=v(f(\rho)/\rho)y$ with $0\le f(\rho)/\rho\le2/\eps$, hence $0\le y(t)\le y(\ta)e^{2\mu}$. By
Lemma~\ref{lem:flow}(b), $|\psi_t|\le v\rho|\sin\phi|/\max(\rho,\eps)^2$, where $\rho|\sin\phi|=|T\times r|=|y\cos\psi-z\sin\psi|\le y+\rho|\psi|$
and $\rho/\max(\rho,\eps)^2\le1/\eps$; so $|\psi|'\le(v/\eps)|\psi|+vy/\eps^2$ almost everywhere. By Gronwall,
\[
|\psi(\ttd)|\le e^{\mu}\Bigl[|\psi(\ta)|+\frac{vI}{\eps^2}e^{2\mu}y(\ta)\Bigr]\le e^\mu\bigl[\theta_{\mathrm f}+\mu e^{2\mu}y_{\mathrm f}\,\xi/\eps\bigr]
Q^\mu\xi^{1-\mu}\le\theta_{\mathrm v}\xi^{1-\mu}
\]
($\xi/\eps<\xi_1$). By Lemma~\ref{lem:flow}(b) and Lemma~\ref{lem:field}(ii), $w_t\le va^*w\le(2v/\eps)w$, so
$w(\ttd)\le w(\ta)e^{2\mu}\le W_{\mathrm v}$; $w\ge1$ by Lemma~\ref{lem:flow}(a). For $\xi=0$: the virtual front is $C^1$ and coincides
at $\ta$ with the true front on $(\alpha_0,\alpha_0+\xi_1\eps)$, so $w(\alpha_0,\ta)\le W_{\mathrm f}$ and its tangent at $\alpha_0$ is the
limit of the true tangents, $\nhat$ (Lemma~\ref{lem:fwd}(e)); thus $\psi(\ta)=0=y(\ta)$, and the same inequalities give $\psi(\alpha_0,\ttd)=0$
and $w(\alpha_0,\ttd)\le W_{\mathrm v}$.
\end{proof}

% sec7b.tex -- writer W5. Sections 7.3-7.5: tail, bulk and hole, the curve x^0, the touch-down theorem.
% Standing: 0<v<=\vbar, 0<eps<=\epsbar, the constants of Definition 7.1, Lemma 7.2 used tacitly; mu=v/(1-v), mu_T=3v/(1+v).

\subsection{The tail and the end}\label{sec:tail}

Throughout Sections~\ref{sec:tail}--\ref{sec:x0}, $0<v\le\vbar$, $0<\eps\le\epsbar$, the constants are those of
Definition~\ref{def:tdconst}, and Lemma~\ref{lem:mono} is used tacitly: a constant bounds the corresponding expression formed with
the actual $v$, $\mu=v/(1-v)\le\mub$, $\muT=3v/(1+v)\le\muTb$ and $\eps$. The \emph{tail} consists of the labels
$\theta=\alpha_0-\xi$, $0<\xi\le\eta$, behind the point of tangency; there $\theta\ge\underline\alpha-\eta=\theta_L$.

\begin{lemma}\label{lem:tail}
\begin{enumerate}
\item[(a)] There is a unique $\be=\alpha_0-\xe$ with $\xe\in(0,\eta)$ and $C(\be)=\eps\gmax$, and
\[
\xi_-\eps^{1/(2-2v)}\le\xe\le\xi_+\eps^{1/(2-2v)}\le\bar\xi_{\mathrm e},\qquad\xe\ge\xi_-\eps\ge\eps/2 .
\]
\item[(b)] Every true candidate $\theta\in[\be,\alpha_0]$ is seen by $\ttd$.
\item[(c)] Let $\theta=\alpha_0-\xi$ with $\xe\le\xi\le\eta$. Then $S_\theta$ exists on $[0,\ttd]$, $\rho_\theta>0$ there, and in the
frame of run~2 at $\ttd$:
$h<0$ and $h_{\mathrm t}\xi\le|h|\le h_{\mathrm T}\xi$; $\tau\ge\tau_{\mathrm t}>1$; $-z_{\mathrm T}\xi\le z\le-h_{\mathrm t}\xi$;
$0<y\le y_{\mathrm T}\xi^{2-\muT}$; $\rho\le\rho_{\mathrm T}\xi$; $(\tau/\tau_{\sbb})^v\le T_{\mathrm t}\xi^{-\muT}$;
$-\theta_{\mathrm t}\xi^{1-\muT}\le\psi<0$; $1\le w\le W_{\mathrm t}$; and $|k|\le k_{\mathrm t}\xi^{-\muT}$.
\item[(d)] The virtual and the true candidate $\be$ coincide on $[0,\ttd]$, and $\xend:=x_{\be}(\ttd)$ satisfies
$\xend-X_0=z_{\mathrm E}\nhat+y_{\mathrm E}\nhat^\perp$ with $-z_{\mathrm T}\xe\le z_{\mathrm E}\le-h_{\mathrm t}\xe$ and
$0<y_{\mathrm E}\le y_{\mathrm T}\xe^{2-\muT}$.
\end{enumerate}
\end{lemma}

\begin{proof}
Tail labels lie in the window, where the bounds of Lemma~\ref{lem:tang}(c) hold; in particular $\xi\mapsto C(\alpha_0-\xi)$ is
continuous and strictly increasing on $(0,\eta]$, $C_{\mathrm{lo}}\xi^{2-2v}\le C\le C_{\mathrm{hi}}\xi^{2-2v}$, and
$\sbb\in(0,\pi)$. In run~1 every tail label has $\rho\ge\eps$, since $\theta\ge\theta_L>\as$ (Lemma~\ref{lem:run1}(a)).

(a) $C(\alpha_0-\xi)\le C_{\mathrm{hi}}\xi^{2-2v}\to0$ as $\xi\to0$, and $C(\alpha_0-\eta)\ge C_{\mathrm{lo}}\eta^2>\epsbar\gb\ge\eps\gmax$
\hyptag{SC22}. So $\be$ exists and is unique, and $C_{\mathrm{lo}}\xe^{2-2v}\le\eps\gmax\le C_{\mathrm{hi}}\xe^{2-2v}$. As
$C_{\mathrm{hi}}<1$ \hyptag{SC21} and $\gmax\ge1$, $\xe\ge(\eps/C_{\mathrm{hi}})^{1/(2-2v)}\ge\xi_-\eps^{1/(2-2v)}$; as $\gb/C_{\mathrm{lo}}>1$
\hyptag{SC23}, $\xe\le(\gb/C_{\mathrm{lo}})^{1/(2-2v)}\eps^{1/(2-2v)}\le\xi_+\eps^{1/(2-2v)}$. Finally
$\eps\le\eps^{1/(2-2v)}\le\eps^{1/2}\le\epsbar^{1/2}$, $\xi_-\ge\frac12$ \hyptag{SC28} and $\bar\xi_{\mathrm e}<\eta$ \hyptag{SC24}.

(c) \emph{$h$.} $\rho_0-\sr=\int_\theta^{\alpha_0}\sA\in[\sA_{\min}\xi,\sA_{\max}\xi]$, $0\le\sr(1-\cos\sbb)\le\sr_{\max}\sbb^2/2\le
\sr_{\max}b_{\mathrm{hi}}^2\eta^3\xi/2$, and $0\le(1-v^2)\omega(\ttd)\le(1+v)v\eps/(3(1-v))\le\delta_{\mathrm{tl}}\xi$, by
Lemma~\ref{lem:tang}(b) and $\eps\le\xi/\xi_-$ (from (a)). By \eqref{eq:h2}, $h\le-(\sA_{\min}-\delta_{\mathrm{tl}})\xi=-h_{\mathrm t}\xi<0$
\hyptag{SC25} and $|h|\le h_{\mathrm T}\xi$. In run~2, $\rho\ge C/\gmax>0$ (Lemma~\ref{lem:coord}(c)), and $\rho\ge\eps$ in run~1; so
$S_\theta$ exists on $[0,\ttd]$ with $\rho>0$, and $0<\beta<\pi$ throughout run~2.

\emph{$\tau$.} By Lemma~\ref{lem:coord}(d), $\tau^{1+v}\ge2|h|/((1-v)C)\ge2h_{\mathrm t}\xi/(C_{\mathrm{hi}}\xi^{2-2v})\ge
(2h_{\mathrm t}/C_{\mathrm{hi}})\eta^{-(1-2v)}\ge m_\tau>1$ \hyptag{SC26}; so $\tau\ge\tau_{\mathrm t}>1$. Then
$(1+v)C\tau^{v-1}/(2|h|)\le(1+v)\tau_{\mathrm t}^{\vbar-1}/m_\tau\le e_{\mathrm t}$, and Lemma~\ref{lem:coord}(d) gives
$2|h|/C\le(1-v)\tau^{1+v}\le(1+e_{\mathrm t})2|h|/C$.

\emph{$y$.} $y=C\tau^v=C^{1/(1+v)}(C\tau^{1+v})^{v/(1+v)}$, and $C\tau^{1+v}\le y_{\mathrm c}\xi$ by the bound on $\tau$; with
$C\le C_{\mathrm{hi}}\xi^{2-2v}$,
\[
y\le C_{\mathrm{hi}}^{1/(1+v)}\xi^{(2-2v)/(1+v)}(y_{\mathrm c}\xi)^{v/(1+v)}\le y_{\mathrm T}\xi^{(2-v)/(1+v)},\qquad\frac{2-v}{1+v}=2-\muT .
\]

\emph{$\tau/\tau_{\sbb}$.} $\tau^{1+v}\le y_{\mathrm c}\xi/(C_{\mathrm{lo}}\xi^{2-2v})$ and $\tau_{\sbb}\ge b_{\mathrm{lo}}\xi^2/2$, so
$\tau/\tau_{\sbb}\le\max(1,y_{\mathrm c}/C_{\mathrm{lo}})(2/b_{\mathrm{lo}})\xi^{-3/(1+v)}$ and $(\tau/\tau_{\sbb})^v\le T_{\mathrm t}\xi^{-\muT}$.

\emph{$\psi$.} By Lemma~\ref{lem:tang}(d), $\sfc=\partial_\theta C<0$ and $|\tau_{\sbb}^v\sfc|\le\ell_{\max}\xi/2$; with
$\sa\ge\sa_{\min}>0$ \hyptag{SC13}, $Y_r=(\tau/\tau_{\sbb})^v\tau_{\sbb}^v\sfc/\sa$ (Lemma~\ref{lem:transport}(b)) satisfies $Y_r<0$ and
$|Y_r|\le Y_2\xi^{1-\muT}\le\bar Y_2$. As $\cos\beta<-v$ ($h<0$), $N_\psi<0$, $|N_\psi|\le|Y_r|+v\sin\beta$ and $D_\psi\ge1$; so
$\psi\in(-\pi/2,0)$ and $|\psi|\le|\tan\psi|\le|Y_r|+v\sin\beta$. Further $\rho\ge|h|\ge h_{\mathrm t}\xi$ and
$\sin\beta=y/\rho\le y_{\mathrm T}\xi^{1-\muT}/h_{\mathrm t}$, whence $|\psi|\le\theta_{\mathrm t}\xi^{1-\muT}$.

\emph{$z$, $\rho$.} $z=h-v\rho<0$ and $|z|\ge|h|\ge h_{\mathrm t}\xi$; $\rho\le|z|+y$ gives
$(1-v)|z|\le|h|+vy\le(h_{\mathrm T}+vy_{\mathrm T}\eta^{1-\muT})\xi$, so $|z|\le z_{\mathrm T}\xi$, and
$\rho\le|z|(1+(y/|z|)^2)^{1/2}\le\rho_{\mathrm T}\xi$.

\emph{$w$, $k$.} $w=\sa|D_\psi+iN_\psi|/(1-v^2)\le W_{\mathrm t}$, as $|D_\psi|\le1+v+v\bar Y_2$ and $|N_\psi|\le(1+v)\bar Y_2+v$; and
$w\ge1$ (Lemma~\ref{lem:truefront}). By Lemma~\ref{lem:transport}(d) with $t'=t_1$, $t''=\ttd$ and the bounds $|k_1|\le1+\vk_{\mathrm w}$,
$|\TS(t_1)|\le S_1v$ of Lemma~\ref{lem:tang}(c): $|k|\le(\tau/\tau_{\sbb})^v(1+\vk_{\mathrm w}+S_1v)+|\TS(\ttd)|$. In
Lemma~\ref{lem:transport}(e), $|\sin\phi|\le\sin\beta+|\psi|\le(y_{\mathrm T}/h_{\mathrm t}+\theta_{\mathrm t})\xi^{1-\muT}$, $\rho\ge h_{\mathrm t}\xi$
and $\rho|\sa||\partial_\theta\log C|/w^2\le\rho_{\mathrm T}\xi\cdot\sa_{\max}L_C/\xi$; so $|\TS(\ttd)|\le S_3v\xi^{-\muT}$ and
$|k|\le k_{\mathrm t}\xi^{-\muT}$.

(d) In run~2 the label $\be$ has $\rho\ge C(\be)/\gmax=\eps$, and $\rho\ge\eps$ in run~1; by Lemma~\ref{lem:truefront} its virtual and
true candidates coincide on $[0,\ttd]$. The bounds are (c) at $\xi=\xe$.

(b) \emph{The label $\alpha_0$} has $\rho_{\alpha_0}(\ta)=\eps$ (Lemma~\ref{lem:tang}(b)). \emph{Labels with $\xi_2v\eps\le\xi\le\xe$.}
By monotonicity $C(\alpha_0-\xi)\le C(\be)=\eps\gmax$, and by \eqref{eq:h2}, $h(\ttd)\le-(1+v)\sA_{\min}\xi+(1+v)v\eps/(3(1-v))\le0$;
so $\cos\beta(\ttd)\le-v$, while at $t_1$, $\cos\beta=\cos\sbb\ge\cos\bar b>v$ \hyptag{SC14}. By Lemma~\ref{lem:coord}(b),(c), $\beta$
passes $\arccos v$ during $(t_1,\ttd)$, where $\rho$ attains its minimum $C/\gmax\le\eps$: the candidate is seen.
\emph{Labels with $0<\xi\le\eps/2$.} At $\ta$, $(1-v^2)\omega(\ta)=(1+v)\eps$, so by \eqref{eq:h2}
$h(\ta)=(1+v)\eps-(1+v)(\rho_0-\sr)-\sr(1-\cos\sbb)$. Here $\sA_{\max}=W_{\mathrm w}=e^{\vbar\bar t_1/\sin\theta_L}<2$ ($\theta_L\ge\frac12$
\hyptag{SC3}, $\vbar\le\frac1{10}$ \hyptag{SC35}, $\bar t_1<3$), so $\sA_{\max}\xi\le\eps$ and
$(1+v)(\eps-\sA_{\max}\xi)\ge(1-\sA_{\max}/2)\eps$; and $\sr(1-\cos\sbb)\le\sr_{\max}b_{\mathrm{hi}}^2\xi^4/2\le
\sr_{\max}b_{\mathrm{hi}}^2\epsbar^3\eps/32$. Hence $h_{\mathrm a}\eps\le h(\ta)\le(1+v)(\eps-\sA_{\min}\xi)$, with $h_{\mathrm a}>0$
\hyptag{SC29}. By \eqref{eq:tau2}, $\tau(\ta)\le\beta_2\xi^2(Q_{\mathrm a}/\eps)^{1+\mu}\le\beta_2Q_{\mathrm a}^{1+\mu}\eps^{1-\mu}/4\le
\tau_{\mathrm a}$, and $\tau(\ta)^2\le\tau_{\mathrm a}\beta_2Q_{\mathrm a}^{1+\mu}\xi^2\eps^{-1-\mu}$. Then
$\cos\beta+v\ge1+v-2\tau^2>0$ \hyptag{SC30}, and
\[
\rho(\ta)=\frac h{\cos\beta+v}\le\frac{(1+v)(\eps-\sA_{\min}\xi)}{1+v-2\tau^2}<\eps
\iff2\tau^2\eps<(1+v)\sA_{\min}\xi ,
\]
which holds since $2\tau^2\eps\le2\tau_{\mathrm a}\beta_2Q_{\mathrm a}^{1+\mu}\xi^2\eps^{-\mu}\le\tau_{\mathrm a}\beta_2Q_{\mathrm a}^{1+\mu}
\eps^{1-\mu}\xi\le4\tau_{\mathrm a}^2\xi<\sA_{\min}\xi$ ($\xi\le\eps/2$, \hyptag{SC31}). So the candidate is seen at $\ta$. Since
$\xi_2v\eps\le\eps/2$ \hyptag{SC27} and $\xe\ge\eps/2$, the three ranges cover $[\be,\alpha_0]$.
\end{proof}

\subsection{The bulk and the hole}\label{sec:bulk}

By Lemma~\ref{lem:run1}(e) every true candidate $S_\theta$, $\theta\in(0,2\pi)$, exists on $[0,\ttd]$ with $\rho_\theta>0$, and
$S_0$ is caught at $1/(1-v)<t_1$; as there, $\circab\theta=\min(\theta,2\pi-\theta)$.

\begin{lemma}\label{lem:bulk}
\begin{enumerate}
\item[(a)] \emph{(Bulk.)} Let $\theta\in(0,2\pi)\setminus I_\eta$. In the frame of run~2, $\sbb(\theta)\in[b_{\mathrm l},\pi/3+2e_\beta]$
and $y(\theta,t)\ge y_{\mathrm b}$ for $t\in[t_1,\ttd]$; at $\ttd$ the true front is $C^\infty$ near $\theta$, with $|k|\le k_{\mathrm b}$
and $w\le W_{\mathrm b}W_{\mathrm{all}}(2/\circab\theta)^{2v}$.
\item[(b)] \emph{(Hole.)} $Z_1:=\lim_{\alpha\to0\pm}S_\alpha(t_1)$ exists and equals $P(t_1)-(1+v)$, and
$\Gamma_1:=\{S_\alpha(t_1):0<|\alpha|\le0.005\}\cup\{Z_1\}$ is a $C^1$ arc of finite length with tangent $i$ at $Z_1$. Let $\Phi$ be
the flow map from $t_1$ to $\ttd$ of the field $x\mapsto v(x-P(t))/|x-P(t)|$. On the closed disc $\bar B(1,r_{\mathrm B})$, $\Phi$ is a
$C^\infty$ diffeomorphism onto its image; $\Gamma_1\subset\bar B(1,r_{\mathrm B})$; $Z_{\mathrm s}:=\Phi(Z_1)=\lim_{\alpha\to0\pm}S_\alpha(\ttd)$;
and $\Gamma_{\mathrm s}:=\Phi(\Gamma_1)$ is a $C^1$ arc through $Z_{\mathrm s}$, equal to the true front off $\alpha=0$, whose unit
tangent is $k_{\mathrm b}$-Lipschitz in arclength.
\item[(c)] \emph{(Chords.)} Let $\theta_a,\theta_b\in[0,2\pi)$ with $|e^{i\theta_b}-e^{i\theta_a}|\ge2\eps$, and let $X_\bullet(\ttd)$
($\bullet=a,b$) be the position at $\ttd$ of the virtual candidate $\theta_\bullet$, or of the true candidate $\theta_\bullet\ne0$, or
$Z_{\mathrm s}$ if $\theta_\bullet=0$. Then $|X_b(\ttd)-X_a(\ttd)|\ge|e^{i\theta_b}-e^{i\theta_a}|$.
\end{enumerate}
\end{lemma}

\begin{proof}
(a) \emph{$\sbb$.} By Lemma~\ref{lem:tang}(b), $\alpha_0\in[\underline\alpha,\pi/3]$, and candidates move with speed $\le v$, so
$|S_{\alpha_0}(t_1)-P(t_1)-(e^{i\pi/3}-2)|\le vt_1+|e^{i\alpha_0}-e^{i\pi/3}|+|t_1-2|\le e_\rho<\sqrt3=|e^{i\pi/3}-2|$; as
$\arg(e^{i\pi/3}-2)=5\pi/6$, $|\beta^*-5\pi/6|\le e_\beta$ and $\pi-\beta^*\in[\pi/6-e_\beta,\pi/6+e_\beta]$. In the upper half
($0<\theta<\pi$), $\beta_1$ is monotone on each side of $\alpha_0$ (Lemma~\ref{lem:tang}(a)), so for $\theta\notin I_\eta$,
$\sbb=\beta_1-\beta^*\ge\min(\sbb(\alpha_0-\eta),\sbb(\alpha_0+\eta))\ge b_{\mathrm{lo}}\eta^2=b_{\mathrm l}$, and $\sbb<\pi-\beta^*$; for
$\theta=\pi$, $\sbb=\pi-\beta^*$; in the lower half, $\sbb\in(\pi-\beta^*,2\pi-2\beta^*]$. Hence $\sbb\in[b_{\mathrm l},\pi/3+2e_\beta]$
(using $\pi-\beta^*\ge b_{\mathrm l}$ \hyptag{SC32}), and $\sin\sbb\ge\sin b_{\mathrm l}$ by \hyptag{SC33}.

\emph{$y$.} $\sr=\rho_1\ge1+v\ge1$ (Lemma~\ref{lem:run1}(d)), so $y(t_1)=\sr\sin\sbb\ge y_{\mathrm b}$; $y=C\tau^v$ is non-decreasing in
run~2, so $y\ge y_{\mathrm b}$ on $[t_1,\ttd]$; in particular $\rho\ge y_{\mathrm b}$ and $0<\beta<\pi$ there.

\emph{$\tau$.} $\tau(t_1)=\tau_{\sbb}\ge\tan(b_{\mathrm l}/2)$. At $\ttd$, $\tau=(\rho-z)/y\le2\rho/y_{\mathrm b}$, and
$\rho(\ttd)\le\sr+(1+v)(\ttd-t_1)\le|e^{i\theta}-2|+d_1+(1+v)(\sqrt3+e_\rho)/(1-v)\le H_{\mathrm b}$ (as $|S_\theta(t_1)-P(t_1)-(e^{i\theta}-2)|\le
vt_1+\mu\le d_1$, $\ttd-t_1\le\rho_0/(1-v)$ and $\rho_0\le\sqrt3+e_\rho$). So $1\le\tau(\ttd)/\tau_{\sbb}\le T_{\mathrm b}$ and
$(\tau(\ttd)/\tau_{\sbb})^v\le T_{\mathrm b}^{\vbar}$ \hyptag{SC34}.

\emph{$k$, $w$.} By Lemma~\ref{lem:transport}(e) with $y\ge y_{\mathrm b}$, $|\TS|\le S_{\mathrm b}v$ at $t_1$ and at $\ttd$. By
Lemma~\ref{lem:transport}(d) and Lemma~\ref{lem:run1}(d) ($|k_1|\le k_{\mathrm{all}}$, $w_1\le W_{\mathrm{all}}(2/\circab\theta)^{2v}$),
$|k(\ttd)|\le T_{\mathrm b}^{\vbar}(k_{\mathrm{all}}+S_{\mathrm b}\vbar)+S_{\mathrm b}\vbar=k_{\mathrm b}$ and
$w(\ttd)\le\frac{1+v}{1-v}T_{\mathrm b}^{\vbar}w_1\le W_{\mathrm b}W_{\mathrm{all}}(2/\circab\theta)^{2v}$. The front is $C^\infty$ near $\theta$
because $\rho>0$ on $[0,\ttd]$.

(b) \emph{$\Gamma_1$.} Let $0<\alpha\le\frac12$. By Lemma~\ref{lem:run1}(c), $y(\alpha,t_1)\le2\mathcal E_1\tan(\alpha/2)^{(1-v)/(1+v)}\to0$.
In the frame of run~1, $z=h_1-v\rho_1<0$ and $\rho_1\le|z|+y$, with $|h_1|=(1-v^2)+(1-\cos\alpha)$; so
$|h_1|\le(1-v)|z|\le|h_1|+vy$, hence $0\le|z|-(1+v)\le(1-\cos\alpha+vy)/(1-v)\to0$. Thus $S_\alpha(t_1)\to P(t_1)-(1+v)$, and
from below by the mirror symmetry of run~1. The tangent at $\alpha$ is $e^{i(\pi/2+\delta)}$ with $0\le\delta\le\arctan(d_{\mathrm{arm}}y)\to0$
(Lemma~\ref{lem:run1}(c)), and at the mirror label $e^{i(\pi/2-\delta)}$; both tend to $i$. The stretch $w_1\le W_{\mathrm{arm}}(2/|\alpha|)^{2v}$
is integrable ($2v<1$). Hence the arclength parametrisation of $\Gamma_1$ is continuous, $C^\infty$ with unit speed off $Z_1$, and its
derivative extends continuously (by $i$) to $Z_1$; so $\Gamma_1$ is $C^1$.

\emph{$\Phi$.} If $X(t_1)\in\bar B(1,r_{\mathrm B})$ and $t\in[t_1,\ttd]$, then $|X(t)-1|\le r_{\mathrm B}+vT_2$, while $P(t)$ lies within
$d_P$ of the segment $[2,e^{i\pi/3}]$ (Lemma~\ref{lem:tang}(b)), whose distance from $1$ is $\frac12$ (the foot of the perpendicular,
$(5+i\sqrt3)/4$, is its midpoint). So $|X(t)-P(t)|\ge\rho_{\mathrm h}>0$ \hyptag{SC36}: along these trajectories the field is $C^\infty$
with bounded derivatives, and $\Phi$ is a $C^\infty$ diffeomorphism of $\bar B(1,r_{\mathrm B})$ onto its image. For $|\alpha|\le0.005$,
$|S_\alpha(t_1)-1|\le|e^{i\alpha}-1|+vt_1\le r_{\mathrm B}$, and $Z_1$ is a limit of such points; so $\Gamma_1\subset\bar B(1,r_{\mathrm B})$.
Hence $\Gamma_{\mathrm s}$ is a $C^1$ arc, $\Phi(Z_1)=\lim\Phi(S_\alpha(t_1))=\lim S_\alpha(\ttd)$, and off $\alpha=0$ it is the true
front, whose curvature is bounded by $k_{\mathrm b}$ by (a), since $I_\eta\subset[\frac12,\frac\pi2]$ (Lemma~\ref{lem:tang}(c)) contains no
label with $|\alpha|\le0.005$. The tangent angle of $\Gamma_{\mathrm s}$ is continuous, with derivative bounded by $k_{\mathrm b}$ off one
point; so it is $k_{\mathrm b}$-Lipschitz.

(c) Virtual candidates solve $X'=vF(X-P)$ and true candidates $X'=v(X-P)/|X-P|$ without meeting $P$ on $[0,\ttd]$, all starting at
$e^{i\theta}$ at time $0$; so this is Lemma~\ref{lem:sep}(b) on $[0,\ttd]$. If $X_a(\ttd)=Z_{\mathrm s}$ and $X_b$ starts at
$e^{i\theta}$ with $|1-e^{i\theta}|\ge2\eps$, then $|e^{i\alpha}-e^{i\theta}|\ge|1-e^{i\theta}|$ for all small $\alpha$ of one sign (if
$\theta\ne\pi$), resp.\ $|e^{i\alpha}-e^{i\theta}|>2\eps$ for all small $\alpha$ (if $\theta=\pi$). (The case split is needed because
Lemma~\ref{lem:sep}(b) requires the initial distance $\ge2\eps$ for each approximating $\alpha$; the limit gives the bound.)
Lemma~\ref{lem:sep}(b) applied to $S_\alpha$
and $X_b$, and $\alpha\to0$ with~(b), give the claim.
\end{proof}

\subsection{The curve at touch-down}\label{sec:x0}

\begin{definition}\label{def:x0}
For $\alpha\in[\alpha_0,2\pi+\be]$ (labels mod $2\pi$) let
\[
\Ups(\alpha):=\begin{cases}
x_\alpha(\ttd)\ \text{(virtual candidate)},&\alpha\in[\alpha_0,\alpha_0+\xi_1\eps],\\
S_\alpha(\ttd)\ \text{(true candidate)},&\alpha\in[\alpha_0+\xi_1\eps,2\pi+\be],\ \alpha\ne2\pi,\\
Z_{\mathrm s}\ \text{(Lemma~\ref{lem:bulk}(b))},&\alpha=2\pi .
\end{cases}
\]
Let $\xz:[0,U]\to\R^2$ be $\Ups$ reparametrised by arclength from $\Ups(\alpha_0)=X_0$, let $\Theta(u)$ be the label of $\xz(u)$, and
$u':=U-u$. The \emph{forward zone} is $\{u:\Theta(u)\in[\alpha_0,\alpha_0+\eta]\}$, with $\xi:=\Theta-\alpha_0$; the \emph{tail zone} is
$\{u:\Theta(u)\in[2\pi+\alpha_0-\eta,2\pi+\be]\}$, with $\xi:=2\pi+\alpha_0-\Theta\in[\xe,\eta]$; the rest is the \emph{bulk}. The
\emph{ghost set} is
\[
\Ghost:=\{u:\Theta(u)=2\pi\}\cup\{u:\Theta(u)\in[\alpha_0,\alpha_0+\xi_1\eps)\text{ and }S_{\Theta(u)}\text{ was within $\eps$ of $P$ at
some time}\le\ttd\}.
\]
\end{definition}

That $\Ups$ is well defined and that $\xz$ exists are part of Lemma~\ref{lem:len}. The two definitions of $\Ups$ agree near
$\alpha_0+\xi_1\eps$ (Lemma~\ref{lem:fwd}(d)); the hole labels $0<|\alpha-2\pi|<\as$ are true (Lemma~\ref{lem:run1}(e)). We have
$0\in\Ghost$, since $\Theta(0)=\alpha_0$ and the true candidate $\alpha_0$ coincides with the virtual one up to $\ta$, where
$\rho_{\alpha_0}(\ta)=\eps$ (Lemma~\ref{lem:tang}(b)); and $U\notin\Ghost$, since $\Theta(U)=2\pi+\be$. Thus $\xz(U)=\xend$.

\begin{lemma}\label{lem:len}
$\xz$ is $C^1$ with $|\partial_u\xz|\equiv1$, and $1\le U\le U_{\max}$. On the forward zone $\xi\le u\le W_{\mathrm v}\xi$ and
$|T-\nhat|\le|\psi|\le\theta_{\mathrm v}\xi^{1-\mu}$; on the tail zone $\xi-\xe\le u'\le W_{\mathrm t}(\xi-\xe)$ and
$-\theta_{\mathrm t}\xi^{1-\muT}\le\psi<0$; here $T=e^{i\psi}\nhat$. On $\{u\ge u_{\mathrm v}\eps\}$ the tangent angle of $\xz$ is locally
Lipschitz in $u$, with almost everywhere derivative $k$, where $|k|\le k_{\mathrm f}\xi^{-\mu}$ on the forward zone,
$|k|\le k_{\mathrm t}\xi^{-\muT}$ on the tail zone and $|k|\le k_{\mathrm b}$ on the bulk (across the label $2\pi$ as well). For labels
$\Theta_a<\Theta_b$ with $\Delta:=\Theta_b-\Theta_a$ and $\lambda:=u_b-u_a$:
\[
\lambda=l(a,b)\le W_x\,\Delta\,(4/\Delta)^{2v}/(1-2v),\qquad\text{and}\qquad\lambda>\lambda_0\ \Longrightarrow\ \Delta>\Delta_0 .
\]
\end{lemma}

\begin{proof}
\emph{Regularity.} $\Ups$ is $C^1$ in the label, with $|\partial_\alpha\Ups|\ge1$, on $[\alpha_0,2\pi)$ and on $(2\pi,2\pi+\be]$: on
$[\alpha_0,\alpha_0+\xi_1\eps]$ it is a virtual front (Lemma~\ref{lem:flow}(a)); on the rest it is the true front, which is $C^\infty$
with $w\ge1$ by Lemmas~\ref{lem:fwd}(d), \ref{lem:tail}(c), \ref{lem:bulk}(a) and~\ref{lem:truefront}; the two coincide on a
neighbourhood of $\alpha_0+\xi_1\eps$ (Lemma~\ref{lem:fwd}(d)). Near the label $2\pi$, $\Ups$ is the arc $\Gamma_{\mathrm s}$
(Lemma~\ref{lem:bulk}(b)). So the arclength parametrisation exists once $\Ups$ has finite length (shown below), and it is $C^1$ with
unit speed.

\emph{Pointwise data.} On the forward zone, Lemma~\ref{lem:virt} gives $1\le w\le W_{\mathrm v}$ and the bound on $\psi$, and
$|T-\nhat|=|e^{i\psi}-1|\le|\psi|$; $u=\int_{\alpha_0}^{\Theta}w$ gives $\xi\le u\le W_{\mathrm v}\xi$. On the tail zone,
Lemma~\ref{lem:tail}(c) gives $1\le w\le W_{\mathrm t}$ and the bound on $\psi$, and $u'=\int_\Theta^{2\pi+\be}w$. On
$\{u\ge u_{\mathrm v}\eps\}$ the forward labels have $\xi\ge u/W_{\mathrm v}\ge\xi_1\eps$, so they are true and Lemma~\ref{lem:fwd}(d)
gives the curvature; Lemma~\ref{lem:tail}(c) on the tail zone and Lemma~\ref{lem:bulk}(a),(b) on the bulk.

\emph{Length.} At every label $\ne2\pi$, $w\le W_{\mathrm v}$ on the forward zone, $w\le W_{\mathrm t}$ on the tail zone and
$w\le W_{\mathrm b}W_{\mathrm{all}}(2/\circab\theta)^{2v}$ on the bulk. The function $\theta\mapsto(2/\circab\theta)^{2v}$ is
symmetric and decreasing in the circular distance $\circab\theta$ to $0$; so its integral over an arc of length $\Delta\le2\pi$ is at
most $2\int_0^{\Delta/2}(2/\chi)^{2v}d\chi=\Delta(4/\Delta)^{2v}/(1-2v)$ (the \emph{arc bound}). The bulk labels form an arc of length
$2\pi-2\eta>4$; by the arc bound, its contribution is at most
$W_{\mathrm b}W_{\mathrm{all}}(2\pi-2\eta)\bigl(4/(2\pi-2\eta)\bigr)^{2v}/(1-2v)\le W_{\mathrm b}W_{\mathrm{all}}(2\pi-2\eta)/(1-2v)$; with
$W_{\mathrm v}\eta$ and $W_{\mathrm t}\eta$ for the forward and tail zones, $U\le U_{\max}$. Further
$U\ge|\Ups(\alpha_0+\pi)-X_0|\ge|e^{i(\alpha_0+\pi)}-e^{i\alpha_0}|-2d_{\mathrm s}=2-2d_{\mathrm s}\ge1$, since all candidates and
$X_0=x_{\alpha_0}(\ttd)$ have moved by at most $v\ttd\le d_{\mathrm s}$ (Lemma~\ref{lem:tang}(b)), and \hyptag{SC38}.

As $(2/\circab\theta)^{2v}\ge1$ on $I_\eta$ (there $\circab\theta\le\pi/2+\eta<2$), $w\le W_x(2/\circab\theta)^{2v}$ at every label
$\ne2\pi$, and the arc bound gives the bound on $l$. If
$\Delta\le\Delta_0$, then, as $\Delta\mapsto\Delta^{1-2v}$ is increasing and $4/\Delta_0=4W_x\hat L/\lambda_0\le8/\lambda_0$ \hyptag{SC41},
$l\le W_x\Delta_0(4/\Delta_0)^{2v}/(1-2v)\le W_x\Delta_0(8/\lambda_0)^{2v}/(1-2v)\le W_x\Delta_0\hat L=\lambda_0$; this is the contrapositive
of the last claim.
\end{proof}

\begin{lemma}\label{lem:Nbound}
For all pairs $u_a<u_b$ of $\xz$ with $\eps/2\le\lambda:=u_b-u_a\le\lambda_0$: $N(a,b)\le c_N\lambda^{3/2-\muT}$. More precisely,
$N\le c_{\mathrm{near}}\lambda^{3/2-\mu}$ if $u_a<\max(\lambda,u_{\mathrm v}\eps)$ \textup(near $X_0$\textup);
$N\le c_{\mathrm{tail}}\lambda^{3/2-\muT}$ if $u'_b<\lambda$ and $\xe<\lambda$ \textup(near $\xend$\textup); and
$N\le c_{\mathrm{far}}\lambda^{3/2-\muT}$ otherwise.
\end{lemma}

\begin{proof}
\emph{Near $X_0$.} If $\lambda\ge u_{\mathrm v}\eps$, then $[u_a,u_b]\subset[0,2\lambda)$ and $u_b<2\lambda_0\le\eta$ \hyptag{SC39}; otherwise
$u_a<u_{\mathrm v}\eps\le2u_{\mathrm v}\lambda$, so $[u_a,u_b]\subset[0,(2u_{\mathrm v}+1)\lambda)$, and $u_b<2u_{\mathrm v}\eps\le
2u_{\mathrm v}\epsbar\le\eta$ \hyptag{SC40}. In both cases $[u_a,u_b]$ is an interval of length $\lambda$ in
$[0,(2u_{\mathrm v}+1)\lambda]$ ($u_{\mathrm v}\ge\frac12$, \hyptag{SC46}) contained in $\{u<\eta\}$, which lies in the forward zone
(there $u\ge\xi$). By Lemma~\ref{lem:static}(b) with $\uhat=\nhat$ and Lemma~\ref{lem:len} ($|T-\nhat|\le\theta_{\mathrm v}\xi^{1-\mu}\le
\theta_{\mathrm v}u^{1-\mu}$),
\[
N^2\le\theta_{\mathrm v}^2\int_{u_a}^{u_b}u^{2-2\mu}du\le\theta_{\mathrm v}^2\lambda^{3-2\mu}
\frac{(2u_{\mathrm v}+1)^{3-2\mu}-(2u_{\mathrm v})^{3-2\mu}}{3-2\mu}\le\theta_{\mathrm v}^2\lambda^{3-2\mu}\frac{12u_{\mathrm v}^2+6u_{\mathrm v}+1}{3-2\mub},
\]
since the integrand is increasing (the right-most interval maximises the integral) and $s\mapsto A^s-B^s$ is non-decreasing for
$A>B\ge1$. So $N\le c_{\mathrm{near}}\lambda^{3/2-\mu}$.

\emph{Near $\xend$.} Here $u'_a=u'_b+\lambda<2\lambda$ and $\xe<\lambda\le\lambda_0$, so $\xe+u'_a<2\lambda_0+\min(\bar\xi_{\mathrm e},\lambda_0)
\le\eta$ \hyptag{SC39}; as $u'\ge\xi-\xe$ on the tail zone, the pair lies in the tail zone, where $\xi\le s:=\xe+u'$ and
$|T-\nhat|\le|\psi|\le\theta_{\mathrm t}\xi^{1-\muT}\le\theta_{\mathrm t}s^{1-\muT}$. As $u$ runs over $[u_a,u_b]$, $s$ runs over an interval
of length $\lambda$ in $[0,3\lambda]$; as above,
$N^2\le\theta_{\mathrm t}^2\lambda^{3-2\muT}(3^{3-2\muT}-2^{3-2\muT})/(3-2\muT)\le19\,\theta_{\mathrm t}^2\lambda^{3-2\muT}/(3-2\muTb)$, i.e.\
$N\le c_{\mathrm{tail}}\lambda^{3/2-\muT}$.

\emph{Far:} $u_a\ge\max(\lambda,u_{\mathrm v}\eps)$, and $u'_b\ge\lambda$ or $\xe\ge\lambda$. By Lemma~\ref{lem:len} the tangent angle is
locally Lipschitz on $[u_a,u_b]$ with almost everywhere derivative $k$. Forward-zone labels of the pair have $\xi\ge u/W_{\mathrm v}\ge
\lambda/W_{\mathrm v}$, so $|k|\le k_{\mathrm f}W_{\mathrm v}^{\mub}\lambda^{-\mu}$; tail-zone labels have
$\xi\ge\xe+u'/W_{\mathrm t}\ge\max(\xe,u'_b/W_{\mathrm t})\ge\lambda/W_{\mathrm t}$, so $|k|\le k_{\mathrm t}W_{\mathrm t}^{\muTb}\lambda^{-\muT}$;
bulk labels have $|k|\le k_{\mathrm b}$. As $\lambda<1$ and $\mu\le\muT$, the tangent angle is $K$-Lipschitz on $[u_a,u_b]$ with
$K:=k_*\lambda^{-\muT}$, and $K\lambda\le k_*\lambda_0^{1-\muTb}\le\pi$ \hyptag{SC45}. By Lemma~\ref{lem:static}(c),
$N\le(K/\sqrt{12})\lambda^{3/2}=c_{\mathrm{far}}\lambda^{3/2-\muT}$.

The cases are exhaustive, $\lambda^{3/2-\mu}\le\lambda^{3/2-\muT}$, and $c_N=\max(c_{\mathrm{far}},c_{\mathrm{near}},c_{\mathrm{tail}})$.
\end{proof}

\begin{theorem}\label{thm:touchdown}
Let $0<v\le5.8\cdot10^{-3}$, $0<\eps\le10^{-3}$, and let $P$ follow design~T. Then $\ttd\le T_2$, and the curve $\xz$ of
Definition~\ref{def:x0} is $C^1$ with $|\partial_u\xz|\equiv1$ and $1\le U\le U_{\max}$. With $\nhat$ of Definition~\ref{def:T},
$u'=U-u$, $T$ the unit tangent of $\xz$, $\lambda=u_b-u_a$, and $l,c,N$ referring to $\xz$:
\begin{itemize}
\item[\textup{(T0)}] $\xz(0)=P(\ttd)$; every $u\notin\Ghost$ is a true label; every direction $\theta\in[0,2\pi)$ is seen during $[0,\ttd]$
or is the direction of a true label.
\item[\textup{(T1)}] $N(a,b)\le c_N\lambda^{3/2-\muTb}$ for $\eps/2\le\lambda\le\lambda_0$.
\item[\textup{(T2)}] $c(a,b)\ge\min\bigl(c_{\mathrm{sep}},\kc(u_a+U-u_b)\bigr)$ for $\lambda>\lambda_0$.
\item[\textup{(T3)}] $|T(u)-\nhat|\le\theta_{\mathrm v}\,u^{1-\mub}$ for $u\le G_{\mathrm T}$.
\item[\textup{(T4)}] $|T(u)-\nhat|\le c_{T4}(U-u+z_0)^{1-\muTb}$ for $U-u\le G_{\mathrm T}$.
\item[\textup{(T5)}] $P(\ttd)-\xz(U)=z_0\nhat+y_0\nhat^\perp$ with $\czlo\eps^{1/(2-2v)}\le z_0\le\czhi\eps^{1/2}$ and
$|y_0|\le c_y(z_0+\eps)^{2-\muTb}$.
\end{itemize}
Here $T_2=3.77231$, $U_{\max}=7.17768$, $\lambda_0=0.0337254$, $c_{\mathrm{sep}}=0.0276152$, $\kc=0.951517$, $c_N=3.41207$,
$\muTb=0.0172997$, $\mub=0.00583384$, $\theta_{\mathrm v}=1.153$, $c_{T4}=1.12275$, $\czlo=1.20086$, $\czhi=1.67043$,
$c_y=0.697925$ and $G_{\mathrm T}=0.030573$.
\end{theorem}

\begin{proof}
By Table~\ref{tab:td} (Appendix~\ref{app:const}) the parameters $\vbar=5.8\cdot10^{-3}$, $\eta=0.11639$, $\epsbar=10^{-3}$,
$\xi_1=1.16171$, $\lambda_0=0.0337254$, $G_{\mathrm T}=0.030573$ satisfy \hyptag{SC1}--\hyptag{SC48}, and the numbers displayed are the
constants
$T_2$, $U_{\max}$, $c_{\mathrm{sep}}$, $\kc$, $c_N$, $\muTb$, $\mub$, $\theta_{\mathrm v}$, $\theta_{\mathrm t}/h_{\mathrm t}$,
$h_{\mathrm t}\xi_-$, $z_{\mathrm T}\xi_+$, $c_y$ of Definition~\ref{def:tdconst}, rounded down for $c_{\mathrm{sep}}$, $\kc$, $\czlo$ and up
for the others. So it suffices to prove the theorem with these constants of Definition~\ref{def:tdconst} in place of the numbers. Also
$\mu\le\mub$ and $\muT\le\muTb$, and $\ttd\le\tc\le T_2$ by Lemma~\ref{lem:tang}(b). Regularity and $1\le U\le U_{\max}$ are
Lemma~\ref{lem:len}.

\emph{\textup{(T0)}.} $\xz(0)=x_{\alpha_0}(\ttd)=P(\ttd)$ (Lemma~\ref{lem:tang}(b)). Let $u\notin\Ghost$. If
$\Theta(u)\notin[\alpha_0,\alpha_0+\xi_1\eps)$, then $\Theta(u)\ne2\pi$ and $\xz(u)=S_{\Theta(u)}(\ttd)$, with $S_{\Theta(u)}$ defined on
$[0,\ttd]$ (Lemma~\ref{lem:run1}(e)). If $\Theta(u)\in(\alpha_0,\alpha_0+\xi_1\eps)$, the true candidate has had $\rho\ge\eps$ on
$[0,\ttd]$, so it coincides with the virtual one (Lemma~\ref{lem:truefront}) and $\xz(u)=S_{\Theta(u)}(\ttd)$. So every $u\notin\Ghost$ is
true. Directions: those in $[-\as,\as]$ are seen in run~1 (Lemma~\ref{lem:run1}(a)), those in $[\be,\alpha_0]$ in run~2
(Lemma~\ref{lem:tail}(b)); $\theta\in(\alpha_0,\alpha_0+\xi_1\eps)$ is seen if its label is in $\Ghost$ and is the direction of a true
label otherwise; every other $\theta\in[0,2\pi)$ is the direction of a true label.

\emph{\textup{(T1)}} is Lemma~\ref{lem:Nbound}, as $\lambda^{3/2-\muT}\le\lambda^{3/2-\muTb}$ for $\lambda<1$.

\emph{\textup{(T2)}.} Let $\lambda>\lambda_0$. By Lemma~\ref{lem:len}, $\Delta:=\Theta_b-\Theta_a\in(\Delta_0,2\pi-\xe]$.
\emph{If $\Delta\le2\pi-\Delta_0$,} then $\xz(u_a)$, $\xz(u_b)$ are positions at $\ttd$ as in Lemma~\ref{lem:bulk}(c), with initial points
$e^{i\Theta_a}$, $e^{i\Theta_b}$ at distance $2\sin(\Delta/2)\ge2\sin(\Delta_0/2)=c_{\mathrm{sep}}>2\epsbar\ge2\eps$ ($\Delta_0\le\eta$
\hyptag{SC42}, \hyptag{SC43}); so $c\ge c_{\mathrm{sep}}$. \emph{If $\Delta>2\pi-\Delta_0$,} then $\xi_a:=\Theta_a-\alpha_0\ge0$ and
$\xi_b:=2\pi+\alpha_0-\Theta_b\ge\xe$ satisfy $\xi_a+\xi_b=2\pi-\Delta<\Delta_0\le\eta$: $u_a$ lies in the forward zone and $u_b$ in the
tail zone. In the frame $(X_0;\nhat,\nhat^\perp)$, $z_a=\int_0^{u_a}\cos\psi\,du\ge\cos(\theta_{\mathrm v}\Delta_0^{1-\mub})u_a$, as
$|\psi|\le\theta_{\mathrm v}\xi^{1-\mu}\le\theta_{\mathrm v}\Delta_0^{1-\mub}\le\pi/2$ on $[0,u_a]$ \hyptag{SC44}; and
$z_b\le-h_{\mathrm t}\xi_b\le-(h_{\mathrm t}/W_{\mathrm t})u'_b$ by Lemma~\ref{lem:tail}(c) and $u'_b\le W_{\mathrm t}(\xi_b-\xe)$. So
$c\ge z_a-z_b\ge\kc(u_a+u'_b)$.

\emph{\textup{(T3)}.} Let $u\le G_{\mathrm T}$. As $G_{\mathrm T}<\eta$ \hyptag{SC47}, $u$ lies in the forward zone and $\xi\le u<1$; so
$|T-\nhat|\le\theta_{\mathrm v}\xi^{1-\mu}\le\theta_{\mathrm v}u^{1-\mu}\le\theta_{\mathrm v}u^{1-\mub}$.

\emph{\textup{(T5)}.} $\xz(U)=\xend$, so by Lemma~\ref{lem:tail}(d) $z_0=-z_{\mathrm E}$, $y_0=-y_{\mathrm E}$ and
$h_{\mathrm t}\xe\le z_0\le z_{\mathrm T}\xe$. With Lemma~\ref{lem:tail}(a), $h_{\mathrm t}\xi_-\eps^{1/(2-2v)}\le z_0\le
z_{\mathrm T}\xi_+\eps^{1/(2-2v)}\le z_{\mathrm T}\xi_+\eps^{1/2}$. Next $h_{\mathrm t}\le\sA_{\min}\le1$, so
$|y_0|\le y_{\mathrm T}\xe^{2-\muT}\le y_{\mathrm T}(z_0/h_{\mathrm t})^{2-\muT}\le c_yz_0^{2-\muT}\le c_y(z_0+\eps)^{2-\muTb}$, as
$z_0+\eps\le z_{\mathrm T}\bar\xi_{\mathrm e}+\epsbar<1$ \hyptag{SC48}.

\emph{\textup{(T4)}.} Let $u'\le G_{\mathrm T}$ and $s:=\xe+u'$. Then $\xi\le s\le\bar\xi_{\mathrm e}+G_{\mathrm T}\le\eta$ \hyptag{SC47}: $u$
lies in the tail zone, and $|T-\nhat|\le\theta_{\mathrm t}\xi^{1-\muT}\le\theta_{\mathrm t}s^{1-\muT}$. By (T5), $s\le u'+z_0/h_{\mathrm t}\le
(u'+z_0)/h_{\mathrm t}$; so $|T-\nhat|\le(\theta_{\mathrm t}/h_{\mathrm t})(u'+z_0)^{1-\muT}\le(\theta_{\mathrm t}/h_{\mathrm t})(u'+z_0)^{1-\muTb}$,
using $u'+z_0\le G_{\mathrm T}+z_{\mathrm T}\bar\xi_{\mathrm e}\le1$ \hyptag{SC48}.
\end{proof}

\begin{lemma}[global chords]\label{lem:T2star}
In the setting of Theorem~\ref{thm:touchdown}, let $u_a<u_b$ be labels of $\xz$ with $\lambda=u_b-u_a>\lambda_0$, and put
$d_c:=u_a+U-u_b$. Then, with the constants $W_x$, $\hat L$ and $\kc$ of Definition~\ref{def:tdconst},
\[
c(a,b)\ge\min\Bigl(\kc d_c,\ 2\sin\Bigl(\frac12\min\Bigl(\pi,\frac{\min(\lambda,d_c)}{W_x\hat L}\Bigr)\Bigr)\Bigr).
\]
The product $W_x\hat L$ is at most $1.22122$.
\end{lemma}

\begin{proof}
Let $\Theta_a<\Theta_b$ be the labels of $u_a,u_b$ and $\Delta:=\Theta_b-\Theta_a$; by Lemma~\ref{lem:len}, $\Delta>\Delta_0$. \emph{Case
$\Delta>2\pi-\Delta_0$.} The proof of (T2) (its second case) gives $c\ge\kc d_c$. \emph{Case $\Delta\le2\pi-\Delta_0$.} As in the proof of (T2),
$\xz(u_a)$, $\xz(u_b)$ are positions at $\ttd$ as in Lemma~\ref{lem:bulk}(c), with initial points at distance $2\sin(\tilde\Delta/2)$,
$\tilde\Delta:=\min(\Delta,2\pi-\Delta)\in[\Delta_0,\pi]$; since $2\sin(\tilde\Delta/2)\ge c_{\mathrm{sep}}>2\epsbar\ge2\eps$ \hyptag{SC42},
\hyptag{SC43}, Lemma~\ref{lem:bulk}(c) gives $c\ge2\sin(\tilde\Delta/2)$. By the proof of Lemma~\ref{lem:len}, $w\le W_x(2/\circab\theta)^{2v}$ at
every label $\theta\ne2\pi$, and by the arc bound the integral of $(2/\circab\theta)^{2v}$ over any arc of length $\Delta'\le2\pi$ is at most
$\Delta'(4/\Delta')^{2v}/(1-2v)$. For $\Delta'\ge\Delta_0$, $(4/\Delta')^{2v}\le\max(1,4/\Delta_0)^{2v}\le(8/\lambda_0)^{2v}$ (as
$4/\Delta_0=4W_x\hat L/\lambda_0\le8/\lambda_0$ by \hyptag{SC41}, and $8/\lambda_0\ge1$ as $\lambda_0<1$), and $(8/\lambda_0)^{2v}/(1-2v)\le\hat L$ ($v\le\vbar$); so
the integral of $w$ over such an arc is at most $W_x\hat L\Delta'$. First, $\lambda=\int_{\Theta_a}^{\Theta_b}w\,d\theta\le W_x\hat L\Delta$.
Second, $d_c=\int_{\alpha_0}^{\Theta_a}w+\int_{\Theta_b}^{2\pi+\be}w$, with $\be<\alpha_0$ (Lemma~\ref{lem:tail}(a)); modulo $2\pi$ both label sets lie
in the arc $[\Theta_b,\Theta_a+2\pi]$ of length $2\pi-\Delta\ge\Delta_0$ (as $\Theta_b\le2\pi+\be<2\pi+\alpha_0\le2\pi+\Theta_a$) and are
disjoint, and the pointwise bound depends only on $\theta$ mod $2\pi$; so $d_c\le W_x\hat L(2\pi-\Delta)$. Hence
$\tilde\Delta\ge\min(\lambda,d_c)/(W_x\hat L)$, and since $\tilde\Delta\le\pi$ and $x\mapsto2\sin(x/2)$ is increasing on $[0,\pi]$, the claim
follows. The value of $W_x\hat L$ is in Table~\ref{tab:td}.
\end{proof}


% sec8.tex -- writer W3.
\section{Proof of the main results}\label{sec:proof}

\subsection{Constants of the assembly}\label{sec:asmconst}

The \emph{main tuple} $p=(\gamma,R,m,s_0,\fF,\kappa,c_A,K_R,K_f,\Keps,\eta_1,\qc,\nu,\dL,\delta_1,A_0)$ is that of
Table~\ref{tab:core}; the \emph{stage-A tuple} $\pA$ is $p$ with $\gamma$ replaced by $\gD$, so $M_{\mathrm A}=mR^{-\gD}$, and
$C_A^{\mathrm A}$, $C_B^{\mathrm A}$ are the tilt constants $C_A$, $C_B$ of Definition~\ref{def:param} computed for $\pA$. The free
assembly parameters are $\bar\omega$, $c_k$, $\delta_0$; the touch-down constants are those of Theorem~\ref{thm:touchdown}; $a'$ is the
number of Table~\ref{tab:asm} and $C_0:=c_N$. Put
\begin{equation}\label{eq:asm}
\begin{gathered}
\zeta:=(1-\bar\omega^2)^{1/2}-\vz,\qquad\etaN:=\aD(4)\frac{\vz}\zeta,\qquad\etaD:=\frac{\vz}\zeta,\qquad
\We:=\frac{c_AK_RR}{\kc},\\
\UA:=\We+\delta_0+R,\qquad\UE:=\We+R,\qquad
\Dbar:=\frac{\kappa\We+\epsS}{1-\vz},\qquad G_f:=\We+\epsS+\delta_0,\\
z_0^{\max}:=\czhi\epsS^{1/2},\qquad X_*:=\UE+z_0^{\max}+\Dbar,\qquad
c_{\mathrm{sep}}^{\mathrm A}:=\min\bigl(c_{\mathrm{sep}},\kc(1-\UA)\bigr),\\
\theta_0:=c_{T4}X_*^{1-\muTb-\gD},\qquad\theta_y:=c_y2^{2-\muTb}(z_0^{\max})^{1-\muTb-\gD},\qquad
\kD:=c_k\frac{4(1+\delta_t)}\zeta,\\
\mathrm{(TS)}:=\Bigl(1+\frac{\Dbar}{\kD R}\Bigr)^{\etaN}\Bigl[C_0R^{a'}+\vz(c_n\Theta_b+c_\beta\Theta_{\mathrm p})R^{\gD}\Bigr],\qquad
\Tsharp:=T_2+\frac{\kappa U_{\max}}{1-\vz},
\end{gathered}
\end{equation}
with the heading constants \eqref{eq:dep-headconst} and the departure constants \eqref{eq:dep-const} ($\delta_t$ does not depend on
$\kD$; $\Theta_{\mathrm p}$, $c_n$, $c_\beta$ do).

\medskip\noindent\textbf{Start conditions.}\namedlabel{hyp:C}{(C1)--(C6)} For a tuple $p$ and data $(a',C_0,\lambda_0,c'_{\mathrm{sep}})$:
\begin{itemize}
\item[(C1)] $R\le\lambda_0$;\qquad(C2) $1-\frac12\bigl(C_0\min(K_RR,\lambda_0)^{a'}\bigr)^2\ge c_A$;
\item[(C3)] $\lambda_0\bigl(1-\frac12(C_0\lambda_0^{a'})^2\bigr)\ge c_AK_RR$;\qquad(C4) $c'_{\mathrm{sep}}\ge c_AK_RR$;\qquad
(C5) $\gamma+\nu\le a'$;
\item[(C6)] $C_0R^{a'}\le s_0m(K_f/K_R)^\nu$.
\end{itemize}

\noindent\textbf{Side conditions of the assembly.}\namedlabel{hyp:A}{(A1)--(A6)}
\begin{itemize}
\item[(A1)] (P0)--(P7) for $(p,U_{\max})$, and (C1)--(C6) for $p$ with $(a',C_0,\lambda_0,c_{\mathrm{sep}})$;
\item[(A2)] (P0)--(P7) for $(\pA,\UA)$, (C1)--(C6) for $\pA$ with $(a',C_0,\lambda_0,c_{\mathrm{sep}}^{\mathrm A})$, and
$1-\UA\ge\We$;
\item[(A3)] (heading) $G_f\le G_{\mathrm T}$, $\epsS\le\delta_0\le R/2$, $\bar\sigma\le\frac12$, $\rho_*\ge\epsS$, $m^2<2$,
$\gD\le1-\mub$, $\UA\ge\We+2\epsS+\delta_0$, $\kappa\le2$, $K_R\ge1$;
\item[(A4)] (departure) (D3), $\gD\le1-\muTb$, $\czlo\epsS^{1/(2-2\vz)}\ge\epsS/\zeta$, $\UE\le G_{\mathrm T}$;
\item[(A5)] (tail start) $\gD\ge\gamma+\nu+\etaN$, $\gD+\nu\le a'$, $\mathrm{(TS)}\le s_0m(K_f/K_R)^\nu$,
$K_RR\le1-\UE-\We-\epsS$, and (PT): $(1+\delta_t)\LJ(U_{\max})<\kappa$;
\item[(A6)] $\vz\le\vbar$, $\epsS\le\min(\epsbar,R/\Keps)$, $2\We+R+\epsS\le1$, $K_RR<\lambda_0$, $a'\le1-\muTb$, $\Tsharp\le10.13$.
\end{itemize}
All of them hold; they are checked in Appendix~\ref{app:const}.

\subsection{Start data from touch-down}

\begin{lemma}[Start data]\label{lem:start}
Let $\xz$ be the curve of Theorem~\ref{thm:touchdown}, and assume \textup{(T1)}, \textup{(T2)} of that theorem, a tuple $p$ satisfying
\textup{(C1)--(C6)} for $(a',C_0,\lambda_0,c'_{\mathrm{sep}})$ with $a'\le1-\muTb$, $C_0=c_N$, $K_RR<\lambda_0$, and a set
$\mathcal P$ of pairs such that $c(a,b;\ttd)\ge c'_{\mathrm{sep}}$ for the pairs of $\mathcal P$ with $\lambda>\lambda_0$. Then
$c(a,b;\ttd)\ge c_A\min(\lambda,K_RR)$ for the pairs of $\mathcal P$ with $\lambda\ge\eps/2$, and
\textup{(S2)} holds for all pairs with $\eps/2\le\lambda\le R$.
\end{lemma}

\begin{proof}
At $\ttd$ the labels are arclength, so $l=\lambda$ and $c=\lambda-E=\lambda-\frac12N^2$. For $\eps/2\le\lambda\le\lambda_0$
($\lambda_0<1$), (T1) and $a'\le1-\muTb$ give $N\le C_0\lambda^{1/2+a'}$.

\emph{(S2).} Let $\eps/2\le\lambda\le R$ ($\le\lambda_0$ by (C1)). By (C6), (C5) and $\lambda/R\le1$,
\[
N\le C_0R^{a'}\Bigl(\frac\lambda R\Bigr)^{a'}\lambda^{1/2}\le s_0m\Bigl(\frac{K_f}{K_R}\Bigr)^\nu\Bigl(\frac\lambda R\Bigr)^{\gamma+\nu}\lambda^{1/2}
=s_0M\lambda^{1/2+\gamma}\Bigl(\frac{K_f\lambda}{K_RR}\Bigr)^\nu .
\]
Since $K_f\lambda/(K_RR)\le\min\bigl(1,K_f\lambda/\min(a,K_RR)\bigr)$ (as $K_f<K_R$ and $\lambda\le R$), this is (S2), also when $a=0$.

\emph{Chords.} Let $\varphi(\lambda):=\lambda\bigl(1-\frac12C_0^2\lambda^{2a'}\bigr)$, so $c\ge\varphi(\lambda)$ for
$\eps/2\le\lambda\le\lambda_0$; $\varphi$ is concave on $[0,\infty)$. If $\lambda\le K_RR$ ($<\lambda_0$), then
$c\ge\lambda\bigl(1-\frac12(C_0(K_RR)^{a'})^2\bigr)\ge c_A\lambda$ by (C2). If $K_RR<\lambda\le\lambda_0$, then
$c\ge\min(\varphi(K_RR),\varphi(\lambda_0))\ge c_AK_RR$ by concavity, (C2) and (C3). If $\lambda>\lambda_0$ and the pair is in
$\mathcal P$, then $c\ge c'_{\mathrm{sep}}\ge c_AK_RR$ by (C4).
\end{proof}

\subsection{The conditional theorem}

\begin{proposition}[Conditional form]\label{prop:cond}
Let $d=1$, $0<v\le\vz$ and $0<\eps\le\epsS$. Let $P$ be design~T (Definition~\ref{def:T}) on $[0,\ttd]$, followed by the chase
(Definition~\ref{def:chase}) on the working front from the curve $\xz$ of Theorem~\ref{thm:touchdown} at $\ttd$, from $a_0=0$ and cut
at $U$, and constant after it stops. Then $P$ is well defined, $P(0)=0$, $P$ is $1$-Lipschitz, $P$ depends only on $(v,\eps)$,
and every direction $\alpha\in[0,2\pi)$ is seen by $\Tstop\le\ttd+\kappa U/(1-v)\le\Tsharp$.
\end{proposition}

\begin{proof}
We use the conditions (A1)--(A6). All constants are numbers computed at $(\vz,\epsS)$; in the arguments below the actual $v$ and
$\eps$ enter only through $v\le\vz$ and $\eps\le\epsS$ inside expressions that are visibly monotone (positive factors,
$1/(1-v)$, $e^{cv}$), through $\eps^{1/(2-2v)}/\eps$, which is non-increasing in $\eps\in(0,1)$ and in $v\in(0,\frac12)$, and through
$D\le\Dbar$, $z_0\le z_0^{\max}$. The touch-down constants are valid for all $v\le\vz$, $\eps\le\epsS\le\epsbar$
(Theorem~\ref{thm:touchdown}). Labels $u$, $u'=U-u$ and $\lambda$ refer to $\ttd$, where they are arclength.

\emph{Step 1 (stage A).} Let $P_{\mathrm A}$ be design~T followed by the chase on the same working front cut at $\UA$ (note
$\UA\le1-\We<1\le U$ by (A2)); it is the chase of Theorem~\ref{thm:end} on the front restricted to $[0,\UA]$. A pair in $[0,\UA]$
with $\lambda>\lambda_0$ has $a+U-b\ge U-\UA\ge1-\UA$, so (T2) gives $c\ge c_{\mathrm{sep}}^{\mathrm A}$. Lemma~\ref{lem:start}
with $\pA$, $c'_{\mathrm{sep}}=c^{\mathrm A}_{\mathrm{sep}}$ and $\mathcal P$ the pairs in $[0,\UA]$ ((A2), (A6)) gives (S1) and
(S2). With (P0)--(P7) for $(\pA,\UA)$ and $\eps\le R/\Keps$, Theorem~\ref{thm:end} with excuse length $0$ applies: conclusions
(a)--(e) hold for $P_{\mathrm A}$ with the tuple $\pA$.

\emph{Step 2 (common legs; $\tW$; duration).} Let $a_j,t_j$ be the bases and leg starts of $P$. Suppose $P=P_{\mathrm A}$ on
$[0,t_j]$ and $a_j<\We$. By Lemma~\ref{lem:chase}(4) the two working fronts agree on $[\ttd,t_j]$, so the bases agree, and since
$l_{t_j}(a_j,\UA)\ge\UA-a_j>\delta_0+R>\eps$ (Lemma~\ref{lem:flow}(a)), the rule of Definition~\ref{def:chase} gives the same
$a_{j+1}<\UA$ for the cuts at $\UA$ and at $U$, and leg $j+1$ is non-final for both. Both paths then run along the same ray
until they meet $x_{a_{j+1}}$, and by causality the fronts still agree, so $P=P_{\mathrm A}$ on $[0,t_{j+1}]$. By
Theorem~\ref{thm:end}(b) for $P_{\mathrm A}$ each such leg advances the base by at least $\eps/\kappa$ and at most $\eps$. Hence
$\jW$ exists, $a_{\jW}\in[\We,\We+\eps)$, $\jW<\kappa\We/\eps+1$, each leg before $\tW$ lasts at most $\eps/(1-v)$
(Lemma~\ref{lem:chase}(1)), and
\[
P=P_{\mathrm A}\ \text{on }[0,\tW],\qquad D:=\tW-\ttd\le\frac{\kappa\We+\eps}{1-v}\le\Dbar .
\]

\emph{Step 3 (heading).} Lemma~\ref{lem:heading} applies to $P_{\mathrm A}$ with $\nhat$ of Definition~\ref{def:T}: its hypothesis on
$T(\cdot,\ttd)$ is (T3), as $G_f\le G_{\mathrm T}$; the conclusions of Theorem~\ref{thm:end} hold by Step~1; (A3) holds. For the legs
$j+1\le\jW$ (non-final, $a_j<\We$, unit speed in direction $\uhat_{j+1}$) we get, by Step~2,
\[
|P'|=1,\qquad|P'-\nhat|\le\hH(t-\ttd)\qquad\text{for a.e.\ }t\in[\ttd,\tW].
\]

\emph{Step 4 (departure).} The departure arc is $s\mapsto x(U-s,t)$, $s\in[0,\UE]$ ($\UE\le U$): a virtual flow of $P$ from
$\xz(U-\cdot)$, unit speed at $\ttd$, whose label $0$ is $\xz(U)$. Apply Theorem~\ref{thm:departure} with $\uhat:=-\nhat$ and
$D=\tW-\ttd$. (D0): by (T5), $\xz(U)-P(\ttd)=z_0\uhat+y_0\uhat^\perp$ (as $\uhat^\perp=-\nhat^\perp$); by (A4) and monotonicity,
$z_0\ge\czlo\eps^{1/(2-2v)}\ge\eps/\zeta$, in particular $z_0\ge\eps$; $z_0\le\czhi\eps^{1/2}\le z_0^{\max}$; and as
$1-\muTb-\gD\ge0$,
\[
|y_0|\le c_y(2z_0)^{2-\muTb}=c_y2^{2-\muTb}z_0^{1-\muTb-\gD}z_0^{1+\gD}\le\theta_yz_0^{1+\gD}.
\]
(D1): the tangent of the arc is $-T(U-s)$ and $|-T-\uhat|=|T-\nhat|$; by (T4) ($s\le\UE\le G_{\mathrm T}$) and $s+z_0\le X_*$,
$|T-\nhat|\le c_{T4}(s+z_0)^{1-\muTb}\le\theta_0(s+z_0)^{\gD}$. (D2) on $[0,D]$ by Step~3:
$|\uhat\times P'|=|\nhat\times(P'-\nhat)|\le\hH$ and $P'\cdot\uhat=-1+\frac12|P'-\nhat|^2<0$, as $\hH\le\omega_0\le\frac12$.
(D3) is part of (A4). Note that $l,c,E,N$ of a pair do not depend on the orientation of the labels.

\emph{Step 5 ((TW) and (PT)).} Let $(a,b)$ be a tail pair with $\eps/2\le\lambda\le R$, i.e.\ $U-b<\We$. Then $U-a<\We+R=\UE$, so the
pair lies in the departure arc, with labels $U-b<U-a$. At $\tW$ ($\toff=D$), Theorem~\ref{thm:departure}(iv) gives
$l_{\tW}(a,b)\le(1+\delta_t)\lambda$. Next, $g_W=a-a_{\jW}\ge(U-\UE)-(\We+\eps)\ge1-\UE-\We-\epsS\ge K_RR\ge K_f\lambda$ by (A5), so
$\min(g_W,K_RR)=K_RR$. By Theorem~\ref{thm:departure}(v), (T1) ($\lambda\le R\le\lambda_0$, $a'\le1-\muTb$) and
$1+D/(\kD\lambda)\le(R/\lambda)(1+\Dbar/(\kD R))$, with $K_\Theta:=c_n\Theta_b+c_\beta\Theta_{\mathrm p}$,
\[
N(a,b;\tW)\le\Bigl(1+\frac{\Dbar}{\kD R}\Bigr)^{\etaN}\Bigl(\frac R\lambda\Bigr)^{\etaN}\lambda^{1/2}\bigl[C_0\lambda^{a'}+\vz K_\Theta\lambda^{\gD}\bigr].
\]
The bound required by (TW) is $s_0M\lambda^{1/2+\gamma}(K_f\lambda/(K_RR))^\nu$. The ratio of the two is at most
\[
\frac{(1+\Dbar/(\kD R))^{\etaN}}{s_0m(K_f/K_R)^\nu}\Bigl[C_0R^{a'}\Bigl(\frac\lambda R\Bigr)^{a'-\etaN-\gamma-\nu}
+\vz K_\Theta R^{\gD}\Bigl(\frac\lambda R\Bigr)^{\gD-\etaN-\gamma-\nu}\Bigr]\le\frac{\mathrm{(TS)}}{s_0m(K_f/K_R)^\nu}\le1,
\]
since both exponents are $\ge0$ ($\gamma+\nu+\etaN\le\gD\le a'$) and $\lambda\le R$. So (TW) holds, and (PT) is part of (A5).

\emph{Step 6 (the end theorem on $[0,U]$).} Apply Theorem~\ref{thm:end} with $p$, $U_0=U_{\max}$ and excuse length $\We$: (P0)--(P7)
hold by (A1); $U\le U_{\max}$; $\eps\le R/\Keps$ and $U\ge1\ge2\We+R+\eps$ by (A6). (S1): if $(a,b)$ is good and $\lambda>\lambda_0$,
then $a\ge\We$ or $U-b\ge\We$, so by (T2) $c\ge\min(c_{\mathrm{sep}},\kc\We)=\min(c_{\mathrm{sep}},c_AK_RR)=c_AK_RR$ ((C4) for $p$);
Lemma~\ref{lem:start} with $\mathcal P$ the good pairs and $c'_{\mathrm{sep}}=c_AK_RR$ gives (S1) and (S2). (TW), (PT): Step~5. Hence
all legs are defined, the chase stops at $\Tstop\le\ttd+\kappa U/(1-v)\le T_2+\kappa U_{\max}/(1-\vz)=\Tsharp$, and the direction of
every true label of $[0,U]$ is seen by $\Tstop$ (Theorem~\ref{thm:end}(b)).

\emph{Step 7 (all directions; the path).} By (T0), every direction is seen during $[0,\ttd]$ or is the direction of a true label, so
every direction is seen by $\Tstop$. $P$ is design~T (two unit-speed runs) on $[0,\ttd]$, $P(0)=0$, it is $1$-Lipschitz afterwards
(Lemma~\ref{lem:chase}(0)) and constant after $\Tstop$. Design~T depends only on $(v,\eps)$, hence so do $\ttd$, $\xz$ and $U$, and
each leg is determined by $\xz$ and by $P$ at earlier times (Lemma~\ref{lem:chase}(4)); so $P$ depends only on $(v,\eps)$.
\end{proof}

\noindent\textbf{Coverage.} Table~\ref{tab:cover} summarises which result controls which labels and pairs at which times.
\begin{table}[ht]
\centering\small
\begin{tabular}{@{}p{0.39\textwidth}p{0.2\textwidth}p{0.33\textwidth}@{}}
\toprule
object & times & established by\\
\midrule
directions not carried by true labels (ghost set, run-1 catches) & $[0,\ttd]$ & (T0) of Theorem~\ref{thm:touchdown}\\
stage-A chase on $[0,\UA]$: (J1), (J2), (CA), tilts, distances & $[\ttd,\tW]$ & Theorem~\ref{thm:end}, excuse $0$; data Lemma~\ref{lem:start}\\
$P=P_{\mathrm A}$; heading $|P'-\nhat|\le\hH$ & $[\ttd,\tW]$ & Step~2; Lemma~\ref{lem:heading}\\
departure arc $\{u'\le\UE\}$; data (TW) of tail pairs & $[\ttd,\tW]$; at $\tW$ & Theorem~\ref{thm:departure}; Step~5\\
body pairs of $[0,U]$ & $[\ttd,\Tstop]$ & Theorem~\ref{thm:end}; (S1), (S2) by Lemma~\ref{lem:start}\\
tail pairs & excused before $\tW$; $[\tW,\Tstop]$ & Theorem~\ref{thm:end} with (TW), (PT)\\
true labels of $[0,U]$ (seen) & by $\Tstop\le\Tsharp$ & Theorem~\ref{thm:end}(b)\\
\bottomrule
\end{tabular}
\caption{Which result covers which labels, pairs and times.}\label{tab:cover}
\end{table}

\subsection{Proofs of Theorem~\ref{thm:main} and Corollary~\ref{cor:exact}}

\begin{proof}[Proof of Theorem~\ref{thm:main}]
Let $d>0$, $\eps\in(0,d)$, $v\in(0,\vz]$, and put $\hat\eps:=\min(\eps/d,\epsS)\in(0,\epsS]$. Let $P$ be the path of
Proposition~\ref{prop:cond} for $(v,\hat\eps)$ and $H(t):=dP(t/d)$; this is design~T followed by the chase for the vision radius
$\eps'=\min(\eps,\epsS d)=d\hat\eps$, scaled to distance $d$. Then $H(0)=0$, $|H'|\le1$ a.e., and $H$ depends only on $(v,\eps,d)$.
The equation $S'=v(S-H)/|S-H|$ is invariant under $(S,H,t)\mapsto(dS,dH,dt)$: if $S$ is the true candidate of direction $\alpha$ for
$P$ (initial point $e^{i\alpha}$), then $t\mapsto dS(t/d)$ is the one for $H$ (initial point $de^{i\alpha}$). By
Proposition~\ref{prop:cond} there is $\hat t\le\Tsharp$ with $|S(\hat t)-P(\hat t)|\le\hat\eps$, so $t:=d\hat t\le\Tsharp d\le T_0d$
by (A6), and $|S_\alpha(t)-H(t)|\le d\hat\eps\le\eps$. The constants $\vz$ and $T_0$ do not involve $\eps$.
\end{proof}

\begin{proof}[Proof of Corollary~\ref{cor:exact}]
By the scaling of the previous proof (an exact catch is preserved) we may take $d=1$. Fix $v\in(0,\vz]$ and check the hypotheses
of Lemma~\ref{lem:compact} with $\bar T:=10.13$.
\emph{Directions:} $D_N:=\{2\pi i/2^N:0\le i<2^N\}$, $N\ge1$, are finite, nested, and their union is dense.
\emph{Radii:} $\eps_N:=2^{-N}\epsS\in(0,1)$, $\eps_N\to0$.
\emph{Paths:} $H_N$ is the path of Theorem~\ref{thm:main} for $(v,\eps_N,1)$ (the same $v$ for all $N$; its vision radius is
$\min(\eps_N,\epsS)=\eps_N$); $H_N(0)=0$ and $|H_N'|\le1$ a.e.\ on $[0,\infty)\supset[0,\bar T]$.
\emph{Proximity with a common time:} by Theorem~\ref{thm:main}, for every $\alpha$, in particular every $\alpha\in D_N$, the true
candidate $S_\alpha$ of $H_N$ comes within $\eps_N$ of $H_N$ at some time $t\le\bar T$, and $\bar T$ does not depend on $N$.
Lemma~\ref{lem:compact} gives a path $H$ on $[0,\bar T]$ with $H(0)=0$ and $|H'|\le1$ that catches every direction exactly by time
$\bar T=10.13$; extended as a constant after $\bar T$, it is a path on $[0,\infty)$ as required.
\end{proof}


\begin{thebibliography}{1}
\bibitem{Hartman} P.~Hartman, \emph{Ordinary Differential Equations}, Wiley, New York, 1964; 2nd ed., Classics in Applied
Mathematics~38, SIAM, Philadelphia, 2002.
\end{thebibliography}

\appendix
% appA.tex -- writer W5. Appendix A: constants and side conditions. The tables are generated by verify_constants.py.
\section{Constants and side conditions}\label{app:const}

The proofs use finitely many explicit numbers. The \emph{inputs} are: $\vz=6\cdot10^{-5}$, $\epsS=4\cdot10^{-11}$ and the bound
$10.13$ of Theorem~\ref{thm:main}; the touch-down parameters $(\vbar,\eta,\epsbar,\xi_1,\lambda_0,G_{\mathrm T})$; the touch-down constants
as printed in Theorem~\ref{thm:touchdown}, with $a'=0.99982$ and $C_0=c_N$; the free parameters of the end theorem and the dependent
quantities $\dL,\delta_1,\nu,A_0,\fF,\gD$ (rounded up), all in Table~\ref{tab:core}; and $\bar\omega$, $c_k$, $\delta_0/R$
(Table~\ref{tab:asm}). Every other constant is defined by a closed formula: those of Definition~\ref{def:tdconst} from the touch-down
parameters, those of Section~\ref{sec:end} (``Parameters'') from the tuple $p$ resp.\ $\pA$ and $U_0=U_{\max}$ resp.\ $U_0=\UA$, and
those of Section~\ref{sec:proof} from all of these. Each side condition is an inequality between two such numbers.

The script \texttt{verify\_constants.py} (Python~3 with \texttt{mpmath}, no other dependency) evaluates all formulas in $50$-digit
arithmetic at exactly the printed inputs and checks $201$ inequalities, in six groups: (SC1)--(SC48) of Definition~\ref{def:tdconst}
($52$ inequalities, as (SC16), (SC37), (SC47), (SC48) consist of two each); that each touch-down constant used in
Theorem~\ref{thm:touchdown} lies on the safe side of its exact value ($12$; $c_{\mathrm{sep}}$, $\kc$, $\czlo$ rounded down, the others
up); (P0)--(P7) for $p$ and for $\pA$ ($2\times33$); (C1)--(C6) for $p$ and for $\pA$ ($2\times6$); (A2)--(A6) of
Section~\ref{sec:proof} ($35$, including (TS), (PT), $a'\le1-\muTb$ and $\Tsharp\le10.13$); and $24$ further elementary inequalities that
the proofs use in passing (such as $\kappa\le\CTL^2(4-\kappa^2)$, $c_\tau\ge1$, $\zeta\le1$, $W_{\mathrm v},W_{\mathrm t}\ge1$,
$\sA_{\max}<2$, $m_0\le1$ and $e_\rho/\sqrt3\le1$ for the $\arcsin$ in Definition~\ref{def:tdconst}, $\pi/2+\eta<2$ and $2\pi-2\eta>4$
in Lemma~\ref{lem:len}). Several of them are trivially true (positivity and range conditions in (P0), (SC2), (SC5), (SC35),
(SC47) first part); they are checked for completeness. The computations after the touch-down use the printed touch-down constants,
never their exact values, so the printed numbers are exactly the numbers to which the proofs refer. By Lemma~\ref{lem:mono} and the
monotonicity remark in the proof of Proposition~\ref{prop:cond}, the conditions at $(\vz,\epsS)$ imply those at every $v\le\vz$,
$\eps\le\epsS$. The script also confirms that the simplified formulas for $\PsiI$ and $\DIIl$ in Section~\ref{sec:end} agree with the
general ones on the branch selected by $\gamma>\frac12$ and $K_f+1\le K_R$. It writes Tables~\ref{tab:core}--\ref{tab:side} below (this
file is generated, not typed) and a full log \texttt{verify\_constants.out}, which lists every condition with both sides and its relative
slack, and every constant of Definition~\ref{def:tdconst}, Section~\ref{sec:end} and Section~\ref{sec:proof} to $20$ digits or more.
It ends with the line \texttt{ALL 201 CONDITIONS HOLD} and exits with a non-zero status if any condition fails. To run it:
\begin{center}
\texttt{python -I verify\_constants.py}\qquad(about one second; output \texttt{appA\_tables.tex}, \texttt{verify\_constants.out}).
\end{center}
The smallest relative slacks are $0$ ($\vz=\vbar$), $10^{-8}$ ($a'\le1-\muTb$) and below $10^{-5}$ for the conditions that defined
a rounded input; Table~\ref{tab:side} lists the $40$ conditions with slack below $2\cdot10^{-3}$. All other conditions have slack at
least $3.9\cdot10^{-3}$; among them (TS) has $1.2\cdot10^{-2}$ and (P6) has $0.10$.

% appA_tables.tex -- GENERATED by verify_constants.py. Do not edit by hand.
\begin{table}[!htbp]\centering\footnotesize\setlength{\tabcolsep}{3pt}
\begin{tabular}{@{}ll@{\hspace{9pt}}ll@{\hspace{9pt}}ll@{\hspace{9pt}}ll@{\hspace{9pt}}ll@{\hspace{9pt}}ll@{}}\toprule
\multicolumn{12}{@{}l}{\emph{Free parameters} (both tuples; $\pA$ differs from $p$ only by $\gamma\mapsto\gD$), and \emph{dependent quantities, rounded up}}\\\midrule
$\gamma$ & $0.614267$ & $R$ & $0.001014$ & $m$ & $0.0561202$ & $s_0$ & $0.0543222$ & $\kappa$ & $1.01594$ & $c_A$ & $0.996095$ \\
$K_R$ & $31.9496$ & $K_f$ & $5.03588$ & $\Keps$ & $999.98$ & $\eta_1$ & $0.834088$ & $\qc$ & $2.05829$ & $\gD$ & $0.614666$ \\
$\dL$ & $1.0000602$ & $\delta_1$ & $1.0001202$ & $\nu$ & $3.10668\cdot10^{-4}$ & $A_0$ & $1.01292$ & $\fF$ & $0.136962$ &  &  \\
\bottomrule\end{tabular}
\par\medskip
\begin{tabular}{@{}lll@{\hspace{12pt}}lll@{\hspace{12pt}}lll@{}}\toprule
constant & $p$ & $\pA$ & constant & $p$ & $\pA$ & constant & $p$ & $\pA$\\\midrule
$M$ & $3.87446$ & $3.88513$ & $\chig$ & $1.01600$ & $1.01600$ & $c_6$ & $2.00796$ & $2.00796$ \\
$\sigma_q$ & $2.79201$ & $2.79057$ & $C_A$ & $6.39375$ & $6.39227$ & $C_B$ & $9.86362$ & $9.86365$ \\
$\CTL$ & $16.2574$ & $16.2559$ & $\CmuA$ & $42.5316$ & $42.5124$ & $\ConeA$ & $6.57340$ & $6.57191$ \\
$c_\rho$ & $0.397157$ & $0.397157$ & $c_R$ & $30.8079$ & $30.8079$ & $K_1$ & $1.96865$ & $1.96865$ \\
$K_2$ & $30.3245$ & $30.3245$ & $\camp$ & $5.17779$ & $5.17779$ & $B_{2U}$ & $2.22894\cdot10^{-3}$ & $2.22894\cdot10^{-3}$ \\
$B_{2\eta}$ & $23.8777$ & $23.8777$ & $r_f$ & $1.32939$ & $1.32939$ & $C_p$ & $12.6407$ & $12.6417$ \\
$B_3$ & $413.253$ & $413.285$ & $\PsiI$ & $3.61659$ & $3.61659$ & $\TI$ & $11.4184$ & $11.4184$ \\
$\Ta$ & $4.17251$ & $4.17251$ & $\Tb$ & $6.37953$ & $6.38210$ & $k_0$ & $9.94343$ & $9.94388$ \\
$\DI$ & $485.642$ & $485.423$ & $\DIIs$ & $6837.12$ & $6837.18$ & $\DIIl$ & $2960.78$ & $2962.24$ \\
$B_5$ & $5162.04$ & $5162.82$ & $I_5$ & $0.662754$ & $0.662229$ & $\Xfar(U_0)$ & $0.0128403$ & $6.25751\cdot10^{-4}$ \\
$\LJ(U_0)$ & $1.01427$ & $1.00197$ & $\bar\theta$ & $0.897490$ & $0.897471$ &  &  &  \\
\bottomrule\end{tabular}
\caption{Parameters of the end theorem (Section~\ref{sec:end}) and derived constants for the main tuple $p$ ($U_0=U_{\max}$) and the stage-A tuple $\pA$ ($U_0=\UA$), at $\vz=6\cdot10^{-5}$ with the touch-down constants of Theorem~\ref{thm:touchdown}. Derived values to six digits; the check uses $50$ digits.}\label{tab:core}
\end{table}
\begin{table}[!htbp]\centering\footnotesize\setlength{\tabcolsep}{3pt}
\begin{tabular}{@{}ll@{\hspace{10pt}}ll@{\hspace{10pt}}ll@{\hspace{10pt}}ll@{}}\toprule
$Q$ & $1.84339$ & $\sA_{\min}$ & $0.994021$ & $\sA_{\max}$ & $1.00015$ & $W_{\mathrm{all}}$ & $1.00030$ \\
$W_{\mathrm b}$ & $1.00102$ & $W_{\mathrm v}$ & $1.00646$ & $u_{\mathrm v}$ & $1.01350$ & $W_{\mathrm t}$ & $1.00599$ \\
$W_x$ & $1.00646$ & $\hat L$ & $1.00078$ & $\Delta_0$ & $0.0330607$ & $k_*$ & $1.21277$ \\
$c_{\mathrm{far}}$ & $0.350095$ & $c_{\mathrm{near}}$ & $2.72159$ & $c_{\mathrm{tail}}$ & $2.60831$ & $\theta_{\mathrm t}$ & $1.03637$ \\
$h_{\mathrm t}$ & $0.994006$ & $z_{\mathrm T}$ & $1.00037$ & $y_{\mathrm T}$ & $0.612388$ & $\xi_-$ & $1.27791$ \\
$\xi_+$ & $1.56504$ & $\bar\xi_{\mathrm e}$ & $0.0494909$ & $C_{\mathrm{hi}}$ & $0.612345$ & $C_{\mathrm{lo}}$ & $0.408294$ \\
\bottomrule\end{tabular}
\par\medskip
\begin{tabular}{@{}ll@{\hspace{10pt}}ll@{\hspace{10pt}}ll@{}}\toprule
$T_2$ & $3.7324640\to3.73247$ & $U_{\max}$ & $6.2931948\to6.29320$ & $c_{\mathrm{sep}}$ & $0.033059147\to0.0330591$ \\
$\kc$ & $0.98808912\to0.988089$ & $c_N$ & $2.7215935\to2.72160$ & $\muTb$ & $1.7998920\cdot10^{-4}\to1.79990\cdot10^{-4}$ \\
$\mub$ & $6.0003600\cdot10^{-5}\to6.00037\cdot10^{-5}$ & $\theta_{\mathrm v}$ & $1.0700235\to1.07003$ & $c_{T4}$ & $1.0426250\to1.04263$ \\
$\czlo$ & $1.2702545\to1.27025$ & $\czhi$ & $1.5656236\to1.56563$ & $c_y$ & $0.61979693\to0.619797$ \\
\bottomrule\end{tabular}
\caption{Constants of Definition~\ref{def:tdconst} at $\vbar=6\cdot10^{-5}$, $\eta=0.1$, $\epsbar=10^{-3}$, $\xi_1=1.007$, $\lambda_0=0.0333$, $G_{\mathrm T}=0.05$ (six digits; the check uses $50$ digits). Top: the ingredients of the constants of Theorem~\ref{thm:touchdown}; all 116 constants of Definition~\ref{def:tdconst} are listed in \texttt{verify\_constants.out}. Bottom: the constants used in Theorem~\ref{thm:touchdown} and afterwards, exact value $\to$ value used; $c_{\mathrm{sep}},\kc,\czlo$ are rounded down, the others up.}\label{tab:td}
\end{table}
\begin{table}[!htbp]\centering\footnotesize\setlength{\tabcolsep}{3pt}
\begin{tabular}{@{}ll@{\hspace{10pt}}ll@{\hspace{10pt}}ll@{\hspace{10pt}}ll@{}}\toprule
\multicolumn{8}{@{}l}{\emph{Free:} $\bar\omega=0.338254$, $c_k=1.21399$, $\delta_0=0.149785\,R$; and $\vz=6\cdot10^{-5}$, $\epsS=4\cdot10^{-11}$, $a^\prime=0.99982$, $C_0=c_N$}\\\midrule
$\zeta$ & $0.940995$ & $\etaN$ & $8.76732\cdot10^{-5}$ & $\etaD$ & $6.37623\cdot10^{-5}$ & $\We$ & $0.0326594$ \\
$\delta_0$ & $1.51882\cdot10^{-4}$ & $\UA$ & $0.0338253$ & $\UE$ & $0.0336734$ & $\Dbar$ & $0.0331820$ \\
$z_0^{\max}$ & $9.90191\cdot10^{-6}$ & $X_*$ & $0.0668653$ & $G_f$ & $0.0328113$ & $c^{\mathrm A}_{\mathrm{sep}}$ & $0.0330591$ \\
$\theta_0$ & $0.367834$ & $\theta_y$ & $0.0292979$ & $C_0$ & $2.72160$ & $\bar\sigma$ & $0.0351162$ \\
$\theta_G$ & $0.286852$ & $k_\sigma$ & $1.01819$ & $C_3$ & $26.1196$ & $c_\tau$ & $1.00158$ \\
$\rho_*$ & $8.58751\cdot10^{-4}$ & $K_{\mathrm{rot}}$ & $52.2591$ & $\omega_0$ & $0.150567$ & $\theta_{h1}$ & $25.3201$ \\
$\vartheta_h$ & $33.5158$ & $\theta_{0y}$ & $0.367834$ & $x_c$ & $1.51643\cdot10^{-4}$ & $L_h$ & $6.20168$ \\
$\Theta_0$ & $0.371554$ & $\Theta_b$ & $36.3821$ & $K_w$ & $0.373655$ & $\delta_t$ & $2.24195\cdot10^{-5}$ \\
$\omega_T$ & $0.0704551$ & $\bar\omega_y$ & $5.55555\cdot10^{-3}$ & $\kD$ & $5.16057$ & $\Theta_{\mathrm p}$ & $34.3034$ \\
$c_n$ & $4.74078$ & $c_\beta$ & $3.30993$ & $c_n\Theta_b+c_\beta\Theta_{\mathrm p}$ & $286.021$ & $C_0R^{a^\prime}$ & $2.76313\cdot10^{-3}$ \\
$\mathrm{(TS)}$ & $3.01155\cdot10^{-3}$ & $s_0m(K_f/K_R)^\nu$ & $3.04682\cdot10^{-3}$ & $(1+\delta_t)\LJ(U_{\max})$ & $1.01430$ & $\Tsharp$ & $10.1264$ \\
\bottomrule\end{tabular}
\caption{Constants of the assembly (Section~\ref{sec:proof}, Lemma~\ref{lem:heading}, Theorem~\ref{thm:departure}), at $(\vz,\epsS)$ with the touch-down constants as used.}\label{tab:asm}
\end{table}
\begin{table}[!htbp]\centering\scriptsize\setlength{\tabcolsep}{2.5pt}\renewcommand{\arraystretch}{1.05}
\begin{tabular}{@{}>{\raggedright\arraybackslash}p{96pt}rrr@{\hspace{10pt}}>{\raggedright\arraybackslash}p{96pt}rrr@{}}\toprule
condition & left & right & slack & condition & left & right & slack\\\midrule
(A6) $\vz\le\vbar$ & $6\cdot10^{-5}$ & $6\cdot10^{-5}$ & $0$ & $c_{T4}$ (up) & $1.0426$ & $1.0426$ & $4.8\cdot10^{-6}$ \\
(A6) $a^\prime\le1-\muTb$ (printed $\muTb$) & $0.99982$ & $0.99982$ & $10^{-8}$ & $\theta_{\mathrm v}$ (up) & $1.07$ & $1.07$ & $6.1\cdot10^{-6}$ \\
$c_y$ (up) & $0.6198$ & $0.6198$ & $1.1\cdot10^{-7}$ & (C2)\,[$p$] & $0.99611$ & $0.99609$ & $1.3\cdot10^{-5}$ \\
$\kc$ (down) & $0.98809$ & $0.98809$ & $1.2\cdot10^{-7}$ & (C2)\,[$\pA$] & $0.99611$ & $0.99609$ & $1.3\cdot10^{-5}$ \\
(P1) $e^{2\vz/(1-\vz)}\le\delta_1$\,[$p$] & $1.0001$ & $1.0001$ & $1.9\cdot10^{-7}$ & (P2) $c_\rho K_f\ge2$\,[$p$] & $2$ & $2$ & $1.7\cdot10^{-5}$ \\
(P1) $e^{2\vz/(1-\vz)}\le\delta_1$\,[$\pA$] & $1.0001$ & $1.0001$ & $1.9\cdot10^{-7}$ & (P2) $c_\rho K_f\ge2$\,[$\pA$] & $2$ & $2$ & $1.7\cdot10^{-5}$ \\
(P1) $1/(1-\vz)\le\dL$\,[$p$] & $1.0001$ & $1.0001$ & $2\cdot10^{-7}$ & (P0) $\dL\ge1$\,[$p$] & $1.0001$ & $1$ & $6\cdot10^{-5}$ \\
(P1) $1/(1-\vz)\le\dL$\,[$\pA$] & $1.0001$ & $1.0001$ & $2\cdot10^{-7}$ & (P0) $\dL\ge1$\,[$\pA$] & $1.0001$ & $1$ & $6\cdot10^{-5}$ \\
(P5)\,[$p$] & $0.13696$ & $0.13696$ & $7.8\cdot10^{-7}$ & (P0) $\delta_1\ge1$\,[$p$] & $1.0001$ & $1$ & $1.2\cdot10^{-4}$ \\
$U_{\max}$ (up) & $6.2932$ & $6.2932$ & $8.2\cdot10^{-7}$ & (P0) $\delta_1\ge1$\,[$\pA$] & $1.0001$ & $1$ & $1.2\cdot10^{-4}$ \\
(P3) $\camp\vz\le\nu$\,[$p$] & $3.1067\cdot10^{-4}$ & $3.1067\cdot10^{-4}$ & $9.7\cdot10^{-7}$ & (P4)\,[$p$] & $0.012834$ & $0.012837$ & $2.8\cdot10^{-4}$ \\
(P3) $\camp\vz\le\nu$\,[$\pA$] & $3.1067\cdot10^{-4}$ & $3.1067\cdot10^{-4}$ & $9.7\cdot10^{-7}$ & (P0) $\eta_1\le K_f/(K_f+1)$\,[$p$] & $0.83409$ & $0.83432$ & $2.8\cdot10^{-4}$ \\
(A5) $\gD\ge\gamma+\nu+\etaN$ & $0.61467$ & $0.61467$ & $1.1\cdot10^{-6}$ & (P0) $\eta_1\le K_f/(K_f+1)$\,[$\pA$] & $0.83409$ & $0.83432$ & $2.8\cdot10^{-4}$ \\
$c_{\mathrm{sep}}$ (down) & $0.033059$ & $0.033059$ & $1.4\cdot10^{-6}$ & (A6) $\Tsharp\le10.13$ & $10.126$ & $10.13$ & $3.6\cdot10^{-4}$ \\
$T_2$ (up) & $3.7325$ & $3.7325$ & $1.6\cdot10^{-6}$ & (SC17) $\rho_{\mathrm c}\xi_1>1$ & $1.0008$ & $1$ & $8.2\cdot10^{-4}$ \\
$\mub$ (up) & $6.0004\cdot10^{-5}$ & $6.0004\cdot10^{-5}$ & $1.7\cdot10^{-6}$ & (SC39) $2\lambda_0+\min(\bar\xi_{\mathrm e},\lambda_0)\le\eta$ & $0.0999$ & $0.1$ & $10^{-3}$ \\
$c_N$ (up) & $2.7216$ & $2.7216$ & $2.4\cdot10^{-6}$ & $c_\tau\ge1$ & $1.0016$ & $1$ & $1.6\cdot10^{-3}$ \\
$\czlo$ (down) & $1.2702$ & $1.2703$ & $3.5\cdot10^{-6}$ & (A5) (PT) & $1.0143$ & $1.0159$ & $1.6\cdot10^{-3}$ \\
$\czhi$ (up) & $1.5656$ & $1.5656$ & $4.1\cdot10^{-6}$ & (P7) $L_J(U_0)<\kappa$\,[$p$] & $1.0143$ & $1.0159$ & $1.6\cdot10^{-3}$ \\
$\muTb$ (up) & $1.7999\cdot10^{-4}$ & $1.7999\cdot10^{-4}$ & $4.4\cdot10^{-6}$ & $\zeta\le1-\bar\omega^2/2$ & $0.94099$ & $0.94279$ & $1.9\cdot10^{-3}$ \\
\bottomrule\end{tabular}
\caption{The 40 side conditions (of 201) with relative slack below $2\cdot10^{-3}$, in increasing order of slack (difference of the two sides divided by the larger one; sides to five digits); $[p]$, $[\pA]$ indicate the tuple, and ``(up), ``(down) the outward rounding of a touch-down constant (exact value against value used). The script checks all 201 conditions in $50$-digit arithmetic and lists each of them, with both sides and its slack, in \texttt{verify\_constants.out}.}\label{tab:side}
\end{table}



\end{document}
